% Section file for Chapter: Twisted Holography and Quantum Gravity
% Result (G8): Holographic Reconstruction

\section{Holographic reconstruction}
% label removed: sec:thqg-holographic-reconstruction
\label{ch:thqg-holographic-reconstruction}% alias for dangling \ref{ch:...}
\index{holographic reconstruction|textbf}
\index{holographic modular Koszul datum!reconstruction}
\index{extension tower|textbf}
\index{shadow obstruction tower!holographic reconstruction}

The shadow obstruction tower
$\Theta_\cA^{\leq 2} \to \Theta_\cA^{\leq 3} \to
\Theta_\cA^{\leq 4} \to \cdots$
is the projective system controlling the holographic
reconstruction of the boundary algebra~$\cA$ from its
shadow-tower data. Finite-order shadow data suffice exactly on the
finite-depth strata. The extension tower
$\{E_\cA(r)\}_{r \geq 2}$ is the dual inductive system: the
$r$-th stage~$E_\cA(r)$ parametrizes all solutions to
the Maurer--Cartan equation through degree~$r$, and the
extension from~$E_\cA(r)$ to~$E_\cA(r{+}1)$ is
controlled by the obstruction class
$o_{r+1}(\cA) \in H^2(\cA^{\mathrm{sh}}_{r+1,0})$.
There are two depths.  The obstruction depth is finite exactly when
the \(H^2\)-obstruction classes vanish from some degree onward.  The
support depth is finite exactly when the bar-intrinsic MC packet has no
\(H^1\)-lift coordinate above some degree.  Finite reconstruction of
\(\Theta_\cA\) uses the support depth.  Obstruction depth alone only
says that later extensions are unobstructed; it does not choose a point
of the torsor of lifts.

The dichotomy theorem is then instantiated for each shadow
archetype. For each of the four classes (Gaussian~(G),
Lie/tree~(L), contact/quartic~(C), mixed/infinite~(M)) we
give the chain-level computation of the obstruction classes,
with the all-degree non-cancellation inputs cited where they are used.
The Virasoro algebra receives
the most detailed treatment: the jet spaces
$J^r(\mathrm{Vir}_c)$ are computed with explicit basis
elements through $r = 8$, the growth rate is proved to be
polynomial, and the quintic obstruction is exhibited at
the chain level.

\begin{remark}[Scope restriction inherited from Volume~I]
The results of this section build on Volume~I
statements (Theorems \texttt{recursive-existence},
\texttt{nms-all-degree-master-equation},
\texttt{gaussian-rmax-two}, \texttt{lie-rmax-three},
\texttt{contact-rmax-four}, \texttt{virasoro-rmax-infinity},
\texttt{shadow-depth-dichotomy},
\texttt{completed-bar-cobar-strong};
Propositions \texttt{obstruction-extension-sequence},
\texttt{gram-wt4}, \texttt{virasoro-shadow-coefficients},
\texttt{independent-sum-factorization};
Corollaries \texttt{lambda-qp},
\texttt{mittag-leffler-shadow-tower},
\texttt{virasoro-quintic-shadow-explicit}) whose scope qualifiers
propagate to the present proofs. Load-bearing cross-volume citations
in proof bodies below carry the explicit ``Volume~I'' prose prefix;
the semantic content (shadow-tower master equation, obstruction
extension sequences, quintic and higher coefficient data) is inherited
without restatement.
\end{remark}

Four disciplines converge:
\emph{deformation theory} (formal moduli, obstruction classes, Postnikov towers),
\emph{representation theory} (Verma modules, singular vectors, Shapovalov determinants controlling jet-space dimensions),
\emph{combinatorial algebra} (graph-sum contributions to each obstruction class), and
\emph{information theory} (the growth rate of $\dim J^r(\cA)$, measuring reconstruction complexity).

% ================================================================
% 9.8.1: THE EXTENSION TOWER
% ================================================================

\subsection{The extension tower and formal moduli}
% label removed: subsec:extension-tower
\index{extension tower!formal moduli}

\begin{definition}[Extension tower]
% label removed: def:extension-tower
\index{extension tower|textbf}
Let $\cA$ be a modular Koszul chiral algebra
(Definition~\ref{V1-def:modular-koszul-chiral}) with
modular convolution dg~Lie algebra $\gAmod$.
The weight filtration $F^r \gAmod$ gives a tower
of quotients
\[
\gAmod / F^{r+1} \gAmod
\twoheadrightarrow
\gAmod / F^r \gAmod,
\qquad r \geq 2.
\]
The \emph{$r$-th extension space} is the Maurer--Cartan
moduli at level~$r$:
\begin{equation}% label removed: eq:extension-space-def
E_\cA(r) := \MC(\gAmod / F^{r+1} \gAmod)
/ {\sim},
\end{equation}
where two MC elements
$\theta_1, \theta_2 \in \MC(\gAmod / F^{r+1})$ are
\emph{gauge-equivalent} ($\theta_1 \sim \theta_2$)
if there exists $\xi \in (\gAmod / F^{r+1})^0$ such that
\begin{equation}% label removed: eq:gauge-equivalence-def
\theta_2
= e^{\ad(\xi)} \cdot \theta_1
+ \frac{e^{\ad(\xi)} - 1}{\ad(\xi)} \cdot d(\xi).
\end{equation}
The moduli space $E_\cA(r)$ is the quotient of
$\MC(\gAmod / F^{r+1})$ by this equivalence relation.
The \emph{extension tower} is the projective system
\begin{equation}% label removed: eq:extension-tower
\cdots \longrightarrow
E_\cA(r{+}1) \xrightarrow{\;\pi_{r+1,r}\;}
E_\cA(r) \longrightarrow \cdots
\longrightarrow E_\cA(2),
\end{equation}
where $\pi_{r+1,r}$ is induced by the quotient map
$\gAmod / F^{r+2} \to \gAmod / F^{r+1}$.
\end{definition}

\begin{remark}[Formal moduli interpretation]
% label removed: rem:extension-formal-moduli
\index{formal moduli!extension tower}
Each $E_\cA(r)$ is a formal moduli problem in the sense
of Lurie~\cite{Lurie-DAG-X}:
the tangent complex at the basepoint
$\Theta_\cA^{\leq r}$ is the truncated chiral Hochschild
complex $\ChirHoch^*(\cA; F^r/F^{r+1})$, and the
obstruction theory is controlled by
$H^2(\cA^{\mathrm{sh}}_{r+1,0})$. The projective limit
$\varprojlim_r E_\cA(r)$ is the full formal moduli
of the MC element $\Theta_\cA$; it is nonempty by the
bar-intrinsic construction
(Theorem~\ref*{V1-thm:mc2-bar-intrinsic}).
\end{remark}

\begin{definition}[Shadow jet space]
% label removed: def:shadow-jet-space
\index{shadow jet space|textbf}
\index{jet space!shadow}
The \emph{$r$-th shadow jet space} of $\cA$ is
\begin{equation}% label removed: eq:shadow-jet-space
J^r(\cA) :=
F^r {\gAmod}^1 / F^{r+1} {\gAmod}^1,
\end{equation}
the degree-$r$, cohomological-degree-$1$ graded piece of
the modular convolution algebra. Its
$d_2$-cohomology is the shadow algebra:
\begin{equation}% label removed: eq:shadow-jet-cohomology
\cA^{\mathrm{sh}}_{r,0} = H^1(J^r(\cA),\, d_2),
\qquad
H^2(\cA^{\mathrm{sh}}_{r+1,0})
= H^2(J^{r+1}(\cA),\, d_2),
\end{equation}
where $d_2 = d + [\Theta_\cA^{\leq 2}, -]$ is the
$\kappa$-twisted differential (the linearized MC operator
at the quadratic level).
\end{definition}

\begin{proposition}[Obstruction-extension sequence]
% label removed: prop:obstruction-extension-sequence
\index{obstruction class!extension sequence}
For each $r \geq 2$, there is an exact sequence
\begin{equation}% label removed: eq:obstruction-exact-sequence
H^1(J^{r+1}(\cA), d_2)
\xrightarrow{\;\iota\;}
E_\cA(r{+}1)
\xrightarrow{\;\pi_{r+1,r}\;}
E_\cA(r)
\xrightarrow{\;\delta\;}
H^2(J^{r+1}(\cA), d_2),
\end{equation}
where:
\begin{enumerate}[label=\textup{(\roman*)}]
\item $\delta(\Theta^{\leq r})$ is the obstruction class
 $o_{r+1}(\cA)$: the image under the connecting
 homomorphism of the MC failure at degree~$r{+}1$;
\item $\iota$ is the inclusion of the fiber: the set
 of lifts of a given element of $E_\cA(r)$ to
 $E_\cA(r{+}1)$, modulo gauge, is a torsor under
 $H^1(J^{r+1}(\cA), d_2)$ when the obstruction vanishes;
\item $\pi_{r+1,r}$ is the projection.
\end{enumerate}
\end{proposition}

\begin{proof}
This is the standard obstruction theory for
Maurer--Cartan elements in filtered dg~Lie algebras.
Let $\Theta^{\leq r} \in E_\cA(r)$ be a solution to the
MC equation modulo $F^{r+1}$. The residual
\[
R_{r+1} :=
\bigl(D_\cA \Theta^{\leq r}
+ \tfrac{1}{2}[\Theta^{\leq r}, \Theta^{\leq r}]\bigr)_{r+1}
\]
is a degree-$2$ element of $J^{r+1}(\cA)$.
By the Bianchi identity $D_\cA^2 = 0$ applied to
$R_{r+1}$, one obtains $d_2 R_{r+1} = 0$, so
$R_{r+1}$ is a $d_2$-cocycle. Define
$\delta(\Theta^{\leq r}) := [R_{r+1}] \in
H^2(J^{r+1}(\cA), d_2)$.

If $\delta(\Theta^{\leq r}) = 0$, then $R_{r+1} =
d_2 \Theta_{r+1}$ for some degree-$1$ element
$\Theta_{r+1} \in J^{r+1}(\cA)$, and
$\Theta^{\leq r+1} := \Theta^{\leq r} + \Theta_{r+1}$
solves the MC equation modulo $F^{r+2}$. The set of
such lifts forms a torsor under $H^1(J^{r+1}(\cA), d_2)$:
two lifts $\Theta_{r+1}$ and $\Theta'_{r+1}$ differ by
a $d_2$-cocycle, and their difference is a
$d_2$-coboundary if and only if they are gauge-equivalent.

If $\delta(\Theta^{\leq r}) \neq 0$, no lift exists:
the MC equation has no solution at degree~$r{+}1$
compatible with the given solution at degree~$r$.
But $\Theta_\cA$ exists unconditionally
(Theorem~\ref*{V1-thm:mc2-bar-intrinsic}), so the
obstruction is always resolved, but
each $o_{r+1}$ is a \emph{nontrivial} class that must
be absorbed by a new degree-$(r{+}1)$ correction.
\end{proof}

% ================================================================
% 9.8.2: THE DICHOTOMY THEOREM
% ================================================================

\subsection{The dichotomy theorem}
% label removed: subsec:dichotomy-theorem
\index{dichotomy theorem|textbf}
\index{Mittag--Leffler condition!dichotomy theorem}

\begin{theorem}[Obstruction depth and support depth of the shadow tower;
\(\ClaimStatusConditional\); licensing tags \(\gamma+\epsilon\)]
% label removed: thm:shadow-depth-dichotomy
\index{shadow depth!dichotomy}
Let $\cA$ be a modular Koszul chiral algebra. Let
\begin{equation}% label removed: eq:shadow-depth-definition-recall
d_{\mathrm{obs}}(\cA) :=
\inf\{r \geq 2 : o_{r'+1}(\cA) = 0
\;\text{for all}\; r' \geq r\}
\;\in\; \{2, 3, 4, \ldots\} \cup \{\infty\},
\end{equation}
where the infimum of the empty set is \(\infty\).  In the
bar-intrinsic gauge component define the support depth
\[
r^\Theta_{\max}(\cA):=
\inf\{r\ge2:\Theta_{\cA,n}=0\ \text{for all}\ n>r\},
\]
where \(\Theta_{\cA,n}\) is the degree-\(n\) component of the MC
element, equivalently the corresponding component of the shadow
reconstruction packet.  Then
\[
d_{\mathrm{obs}}(\cA)\le r^\Theta_{\max}(\cA),
\]
and finite reconstruction is governed by \(r^\Theta_{\max}\), not by
\(d_{\mathrm{obs}}\) alone:
\begin{enumerate}[label=\textup{(\roman*)}]
\item \emph{Finite support depth.}
If \(r^\Theta_{\max}<\infty\), then
\[
\Theta_\cA^{\leq r}=\Theta_\cA^{\leq r^\Theta_{\max}}=\Theta_\cA
\qquad (r\ge r^\Theta_{\max}).
\]
The bar-intrinsic branch \(E^\Theta_\cA(r)\) stabilizes for
\(r\ge r^\Theta_{\max}\).

\item \emph{Obstruction depth.}
If \(d_{\mathrm{obs}}<\infty\), then no new \(H^2\)-obstruction occurs
above degree \(d_{\mathrm{obs}}\).  The lifts above that degree may
still vary by the \(H^1(J^n(\cA),d_2)\)-torsors.  Thus
\(d_{\mathrm{obs}}<\infty\) implies finite reconstruction only under
the extra condition \(\Theta_{\cA,n}=0\) for \(n>d_{\mathrm{obs}}\), or
equivalently all higher lift coordinates vanish in the chosen gauge
slice.

\item \emph{Infinite support depth.}
If \(r^\Theta_{\max} = \infty\), then
$\Theta_\cA^{\leq r} \neq \Theta_\cA^{\leq r-1}$
for infinitely many values of~$r$. The extension
branch does not stabilize: the inverse limit
$\Theta_\cA = \varprojlim_r \Theta_\cA^{\leq r}$
is the limit of a genuinely infinite tower.
\end{enumerate}
\end{theorem}

\begin{proof}
The obstruction class \(o_{r+1}\) is the image of a truncation in
\(H^2(J^{r+1}(\cA),d_2)\) under the connecting morphism of the
obstruction-extension sequence.  If \(\Theta_{\cA,n}=0\) for every
\(n>r\), then the MC equation has no nonzero residual component above
degree \(r\), so all later obstruction classes vanish.  This gives
\(d_{\mathrm{obs}}\le r^\Theta_{\max}\).

The finite-support statement is the definition of
\(r^\Theta_{\max}\) together with completeness and separatedness of the
weight filtration: a filtered MC element equals the inverse limit of
its finite projections, and the projections are constant exactly after
the last nonzero homogeneous component.  This proves
\(\Theta_\cA^{\le r}=\Theta_\cA\) for
\(r\ge r^\Theta_{\max}\) and stabilization of the selected
bar-intrinsic branch.

If \(d_{\mathrm{obs}}<\infty\), the obstruction-extension sequence says
only that every later extension is unobstructed.  The fiber of
\(E_\cA(n)\to E_\cA(n-1)\) over a liftable point is a torsor under
\(H^1(J^n(\cA),d_2)\).  A nonzero torsor coordinate is a nonzero
degree-\(n\) component of the compatible MC packet.  Hence obstruction
vanishing does not imply truncation unless the higher lift coordinates
vanish.

Finally, \(r^\Theta_{\max}=\infty\) means that infinitely many
homogeneous components of the bar-intrinsic MC element are nonzero, so
the finite projections are not eventually constant.  The inverse limit
is the completed MC element by the bar-intrinsic construction
(Theorem~\ref{V1-thm:recursive-existence}) and the branch-level
Mittag--Leffler statement above.
\end{proof}

\begin{corollary}[Mittag--Leffler for the bar-intrinsic shadow branch;
\(\ClaimStatusConditional\); licensing tags \(\gamma+\epsilon\)]
% label removed: cor:mittag-leffler-shadow-tower
\index{Mittag--Leffler condition!shadow tower}
Let \(E^\Theta_\cA(r)\subset E_\cA(r)\) be the formal branch containing
the projection \([\Theta_\cA^{\le r}]\) of the bar-intrinsic MC element,
with transition maps induced from \(E_\cA(r+1)\to E_\cA(r)\).  The
inverse system \(\{E^\Theta_\cA(r)\}_{r\ge2}\) satisfies the
Mittag--Leffler condition. In particular:
\begin{enumerate}[label=\textup{(\roman*)}]
\item $\varprojlim^1_r E^\Theta_\cA(r) = 0$;
\item the natural map
 \[
 \MC(\widehat{\gAmod})_{\Theta}\big/\textup{gauge}
 \longrightarrow
 \varprojlim_r E^\Theta_\cA(r)
 \]
 is a bijection onto the compatible bar-intrinsic branch;
\item the gauge class of \(\Theta_\cA\) is uniquely determined by its
 compatible finite-order projections in this branch.
\end{enumerate}
No surjectivity statement is made for arbitrary components of
\(E_\cA(r)\); a different component may carry a non-vanishing
obstruction class at the next stage.
\end{corollary}

\begin{proof}
For each \(r\), the bar-intrinsic projection
\([\Theta_\cA^{\le r+1}]\) maps to
\([\Theta_\cA^{\le r}]\).  Hence the transition maps on the selected
branch \(E^\Theta_\cA(r+1)\to E^\Theta_\cA(r)\) are surjective onto
that branch.  This is the Mittag--Leffler condition for the countable
inverse system \(E^\Theta_\cA(\bullet)\), giving
\(\varprojlim^1 E^\Theta_\cA(r)=0\).

Completeness and separatedness of \(F^\bullet\gAmod\) identify an MC
element in the completed dg Lie algebra with a compatible system of
finite projections.  Gauge compatibility is taken inside the same
bar-intrinsic branch.  Therefore the displayed map is bijective onto
\(\varprojlim_r E^\Theta_\cA(r)\), and \([\Theta_\cA]\) is determined
by its compatible finite projections.  The obstruction-extension
sequence still controls other components, but it does not force those
components to lift.
\end{proof}

% ================================================================
% 9.8.3: SHADOW DEPTH AS A RECONSTRUCTION INVARIANT
% ================================================================

\subsection{Shadow depth as a reconstruction invariant}
% label removed: subsec:shadow-depth-invariant
\index{shadow depth!reconstruction invariant}

The support shadow depth \(r^\Theta_{\max}(\cA)\) measures the
complexity of the holographic reconstruction problem: how many
homogeneous orders of the bar-intrinsic packet are needed to determine
the MC gauge class \([\Theta_\cA]\).  The obstruction depth
\(d_{\mathrm{obs}}(\cA)\) measures only eventual \(H^2\)-vanishing.
The two agree on the canonical standard-family rows below, but they
are distinct invariants in a general filtered MC moduli problem.

\begin{definition}[Support shadow depth]
% label removed: def:shadow-depth-holographic
\index{shadow depth|textbf}
The \emph{shadow depth} of a modular Koszul chiral
algebra~$\cA$ is the support depth \(r^\Theta_{\max}(\cA)\) of the
bar-intrinsic MC packet.  Equivalently, it is the least \(r\) for
which all homogeneous components \(\Theta_{\cA,n}\) vanish for
\(n>r\); if no such \(r\) exists, the depth is \(\infty\).
The four \emph{shadow classes} are:
\begin{center}
\renewcommand{\arraystretch}{1.4}
\begin{tabular}{clcl}
\toprule
Class & Name & $r^\Theta_{\max}$ & Defining property \\
\midrule
\textbf{G} & Gaussian & $2$
 & bar-intrinsic packet has support only in degree \(2\) \\
\textbf{L} & Lie/tree & $3$
 & support ends at the Lie cubic; Jacobi kills the next obstruction \\
\textbf{C} & Contact/quartic & $4$
 & support ends at the contact quartic in the canonical gauge slice \\
\textbf{M} & Mixed/infinite & $\infty$
 & infinitely many packet components occur \\
\bottomrule
\end{tabular}
\end{center}
\end{definition}

\begin{remark}[Shadow depth versus Koszulness]
% label removed: rem:depth-vs-koszulness
\index{shadow depth!versus Koszulness}
Shadow depth does \emph{not} characterize Koszulness.
Classes~$\mathbf{G}$ (Heisenberg, \(r^\Theta_{\max} = 2\)),
$\mathbf{L}$ (affine Kac--Moody, \(r^\Theta_{\max} = 3\)), and
$\mathbf{C}$ ($\beta\gamma$ system, \(r^\Theta_{\max} = 4\)) are
chirally Koszul: bar cohomology is concentrated and
$m_k = 0$ for $k \ge 3$ on cohomology
(Proposition~\ref{V1-prop:standard-examples-modular-koszul}).
Class~$\mathbf{M}$ (Virasoro, \(r^\Theta_{\max} = \infty\)) is
chirally Koszul (PBW spectral sequence concentrates) but has
non-formal Swiss-cheese structure: the fourth-order pole forces
nonvanishing transferred operations $m_k^{\mathrm{SC}} \neq 0$
for all $k \geq 3$. The Koszul dual
$\mathrm{Vir}_c^! = \mathrm{Vir}_{26-c}$ carries
these higher $A_\infty$ operations. Shadow depth classifies the
\emph{complexity of the Lagrangian self-intersection}
within the standard families.
\end{remark}

% ================================================================
% 9.8.4: CLASS G; THE GAUSSIAN ARCHETYPE
% ================================================================

\subsection{Class~G: the Gaussian archetype}
% label removed: subsec:class-g-gaussian
\index{shadow archetype!Gaussian|textbf}
\index{Heisenberg algebra!shadow depth}

\begin{theorem}[$r^\Theta_{\max}(\mathcal{H}_k) = 2$]
% label removed: thm:gaussian-rmax-two
\index{shadow depth!Heisenberg}
For the Heisenberg algebra $\mathcal{H}_k$ at any nonzero
level $k \neq 0$, the support shadow depth is
\(r^\Theta_{\max}=2\).
Every obstruction class vanishes:
\begin{equation}% label removed: eq:gaussian-all-obstructions-vanish
o_{r+1}(\mathcal{H}_k) = 0
\qquad \text{for all } r \geq 2.
\end{equation}
The bar-intrinsic packet stabilizes immediately:
$\Theta_{\mathcal{H}_k}^{\leq 2} = \Theta_{\mathcal{H}_k}
= \kappaChHodge(\mathcal{H}_k)\,\eta \otimes \Lambda$ with
\(\kappaChHodge(\mathcal{H}_k) = k\).  All higher
\(H^1\)-lift coordinates vanish in the canonical Gaussian gauge slice.
\end{theorem}

\begin{proof}
The proof has two steps: show that all higher OPE products
vanish, and conclude that the graph-sum contributions
at degree $r \geq 3$ are identically zero.

\medskip
\emph{Step~1: All $m_n = 0$ for $n \geq 3$.}

The Heisenberg OPE is
\begin{equation}% label removed: eq:heisenberg-ope-recall
a(z)\,b(w) \sim \frac{k\,\langle a, b \rangle}{(z-w)^2}
\end{equation}
for $a, b \in V = \mathfrak{h}$ (the Cartan subalgebra),
with no simple pole (since $[a,b] = 0$) and no
higher-order poles. The vertex algebra product decomposes
as
\[
Y(a, z) = \sum_{n \in \Z} a_{(n)} z^{-n-1},
\]
where $a_{(n)}b = 0$ for $n \geq 1$ (no singular OPE
poles beyond the double pole) and $a_{(0)}b = 0$
(no Lie bracket). In the bar complex language:
\begin{itemize}
\item The bilinear product $m_2$ gives the Hessian
 $H_{\mathcal{H}} = \kappaChHodge(\mathcal{H}_k) = k$ (from the double pole).
\item The $n$-ary operations $m_n = 0$ for all $n \geq 3$.
 This is immediate: $m_n$ is built from
 iterated OPE residues, and each residue involves an
 $a_{(p)}b$ with $p \geq 0$. Since $a_{(0)} = 0$
 (no Lie bracket) and $a_{(1)} = k\langle a, b\rangle$
 (a scalar, not a new element of the algebra), no
 nontrivial chain of compositions can be formed beyond
 degree~$2$.
\end{itemize}

More precisely, the transferred cyclic minimal model
$\{m_n^{\mathrm{tr}}\}_{n \geq 1}$ of the
$A_\infty$-algebra $(\barB(\mathcal{H}_k), d)$ satisfies
$m_n^{\mathrm{tr}} = 0$ for $n \geq 3$. The reason: the
bar complex $\barB(\mathcal{H}_k)$ is the cofree
conilpotent coalgebra on $s^{-1}\bar{V}$ with quadratic
codifferential, and the homological perturbation lemma
produces no corrections beyond the quadratic term when the
perturbation itself is quadratic.

\medskip
\emph{Step~2: Graph-sum vanishing.}

We first record a combinatorial fact.

\emph{Lemma.} Every connected stable tree with
$n \geq 3$ leaves has at least one vertex of
valence $\geq 3$.

\emph{Proof of Lemma.} A tree with $n$ leaves and all
internal vertices of valence $\leq 2$ is a path, which
has exactly $2$ leaves. For $n \geq 3$, some vertex
must have valence $\geq 3$. \hfill$\square$

\medskip
The degree-$r$ component of $\Theta_\cA^{\leq r}$ is
a graph sum
\[
\Theta_r = \sum_{\Gamma \in \mathcal{G}_r}
\frac{1}{|\mathrm{Aut}(\Gamma)|}
\ell_\Gamma(\Theta),
\]
where $\mathcal{G}_r$ is the set of connected stable
graphs of weight~$r$, and $\ell_\Gamma$ is the graph
amplitude with vertex labels from $\{m_n^{\mathrm{tr}}\}$
and edge propagators $P = H^{-1} = 1/k$. Since
$m_n^{\mathrm{tr}} = 0$ for $n \geq 3$, every graph with
a vertex of valence $\geq 3$ contributes zero. The only
graphs of weight~$r \geq 3$ that could contribute would
need to use only bivalent vertices (from $m_2$), but by
the Lemma, a connected tree on $r \geq 3$ leaves must
have at least one vertex of valence $\geq 3$.
A connected graph on bivalent vertices alone is either a
single edge (weight~$2$) or a cycle. Cycles of length
$\geq 2$ have weight $\geq 3$ but contribute to the
genus expansion, not to the degree expansion at genus~$0$.
At genus~$0$, the tree graphs of weight~$r \geq 3$ must
have at least one vertex of valence~$\geq 3$, so all
contribute zero.

Hence $\Theta_r = 0$ for all $r \geq 3$, and
\begin{equation}% label removed: eq:heisenberg-theta-final
\Theta_{\mathcal{H}_k} = \Theta^{\leq 2}_{\mathcal{H}_k}
= k\,\eta \otimes \Lambda.
\end{equation}
All obstruction classes vanish:
$o_{r+1} = (D_\cA \Theta^{\leq r} +
\frac{1}{2}[\Theta^{\leq r}, \Theta^{\leq r}])_{r+1}
= 0$ because the MC equation is already satisfied
exactly at degree~$2$.
\end{proof}

\begin{remark}[The Gaussian class in generality]
% label removed: rem:gaussian-class-general
\index{Gaussian archetype!general members}
The argument of Theorem~\ref{V1-thm:gaussian-rmax-two}
applies to any chiral algebra with purely quadratic OPE:
the free fermion, the $bc$ ghost system, and more generally
any lattice vertex algebra $V_\Lambda$
(Theorem~\ref{V1-thm:lattice:curvature-braiding-orthogonal}).
In each case, the transferred $A_\infty$-structure is
formal ($m_n^{\mathrm{tr}} = 0$ for $n \geq 3$),
the modular $L_\infty$-algebra $\Definfmod(\cA)$ equals
its strict model, and the shadow obstruction tower terminates at
$r = 2$. The Gaussian class is characterized by:
\begin{equation}% label removed: eq:gaussian-characterization
\cA \in \textbf{G}
\quad\Longleftrightarrow\quad
\text{all OPE poles have order } \leq 2
\quad\Longleftrightarrow\quad
m_n^{\mathrm{tr}} = 0 \;\text{for } n \geq 3.
\end{equation}
\end{remark}

% ================================================================
% 9.8.5: CLASS L; THE LIE/TREE ARCHETYPE
% ================================================================

\subsection{Class~L: the Lie/tree archetype}
% label removed: subsec:class-l-lie-tree
\index{shadow archetype!Lie/tree|textbf}
\index{affine Kac--Moody!shadow depth}

\begin{theorem}[$r^\Theta_{\max}(\widehat{\mathfrak{g}}_k) = 3$]
% label removed: thm:lie-rmax-three
\index{shadow depth!affine Kac--Moody}
For the affine Kac--Moody algebra
$\widehat{\mathfrak{g}}_k$ at generic level $k$, the
support shadow depth is \(r^\Theta_{\max}=3\). The obstruction
classes satisfy:
\begin{equation}% label removed: eq:lie-obstruction-pattern
o_3(\widehat{\mathfrak{g}}_k) \neq 0,
\qquad
o_{r+1}(\widehat{\mathfrak{g}}_k) = 0
\;\;\text{for all } r \geq 3.
\end{equation}
In the canonical Lie/tree gauge slice, the higher
\(H^1(J^n(\widehat{\mathfrak g}_k),d_2)\)-lift coordinates vanish for
\(n>3\).  Thus the support packet, not merely the \(H^2\)-obstruction
profile, terminates at degree~\(3\).
\end{theorem}

\begin{proof}
The proof proceeds by computing the
transferred $A_\infty$-operations and verifying that the
Jacobi identity forces the quartic obstruction to vanish.

\medskip
\emph{Step~1: The OPE structure.}

The current algebra $\widehat{\mathfrak{g}}_k$ has OPE
\begin{equation}% label removed: eq:km-ope-recall
J^a(z)\,J^b(w)
\sim \frac{k\,\kappa^{ab}}{(z-w)^2}
+ \frac{f^{ab}{}_c\,J^c(w)}{z-w},
\end{equation}
where $\kappa^{ab}$ is the Killing form and $f^{ab}{}_c$
are the structure constants of $\mathfrak{g}$. The double
pole contributes the Hessian $H_{\widehat{\fg}} = k\kappa/2$.
The simple pole contributes the cubic operation
$m_2^{(0)}(J^a, J^b) = f^{ab}{}_c\,J^c$, which is
precisely the Lie bracket.

\medskip
\emph{Step~2: Nontrivial cubic shadow.}

The transferred cubic operation
$m_2^{\mathrm{tr}} \neq 0$ comes from the simple-pole
OPE. The cubic shadow is
\begin{equation}% label removed: eq:km-cubic-shadow
\mathfrak{C}_{\widehat{\fg}} =
f^{ab}{}_c\,x_a x_b x^c,
\end{equation}
the Lie bracket cubic. The obstruction class
$o_3 = [D_\cA \Theta^{\leq 2}]_3 \neq 0$ is nontrivial
because the Lie bracket is nontrivial ($\fg$ is non-abelian).
Hence $\Theta^{\leq 3} \neq \Theta^{\leq 2}$: the shadow
tower requires a cubic correction.

\medskip
\emph{Step~3: Quartic obstruction vanishes by Jacobi.}

The quartic obstruction is
\begin{equation}% label removed: eq:km-quartic-obstruction
o_4(\widehat{\fg}_k) =
\bigl[D_\cA \Theta^{\leq 3}
+ \tfrac{1}{2}[\Theta^{\leq 3}, \Theta^{\leq 3}]
\bigr]_4.
\end{equation}
This has two contributions:
\begin{enumerate}[label=(\alph*)]
\item the \emph{boundary bracket}
 $\frac{1}{2}\{\mathfrak{C}, \mathfrak{C}\}_H$:
 the sewing contraction of two cubic vertices through
 the propagator $P = (k\kappa)^{-1}$;
\item the \emph{local quartic vertex}
 $m_3^{\mathrm{tr}}$: the transferred ternary operation.
\end{enumerate}

We claim both contribute zero in cohomology.

\emph{Contribution~(a): the boundary bracket.}
The $H$-Poisson bracket of the Lie bracket cubic with
itself is
\begin{equation}% label removed: eq:km-boundary-bracket
\frac{1}{2}\{\mathfrak{C}, \mathfrak{C}\}_H
= \frac{1}{2} \cdot
\frac{\partial \mathfrak{C}}{\partial x^a}
\cdot P^{ab} \cdot
\frac{\partial \mathfrak{C}}{\partial x_b}.
\end{equation}
Writing out explicitly:
\[
\frac{\partial \mathfrak{C}}{\partial x^a}
= f^{bc}{}_a x_b x_c + 2f^{ab}{}_c x_b x^c.
\]
The full contraction produces a quartic polynomial in
the $x$-variables. The coefficient of each monomial
$x_a x_b x_c x_d$ involves a sum
\[
\sum_e P^{ae}f^{bc}{}_e + \text{permutations},
\]
which is the structure constant of the Lie bracket
composed with itself through the inverse Killing form.
The \emph{Jacobi identity}
\begin{equation}% label removed: eq:jacobi-identity-structure-constants
f^{ab}{}_e f^{ec}{}_d + f^{bc}{}_e f^{ea}{}_d
+ f^{ca}{}_e f^{eb}{}_d = 0
\end{equation}
forces the quartic polynomial
$\frac{1}{2}\{\mathfrak{C}, \mathfrak{C}\}_H$ to
lie in the image of $d_2$: it is a $d_2$-coboundary.

\emph{Contribution~(b): the local quartic vertex.}
The transferred ternary operation $m_3^{\mathrm{tr}}$ on
bar cohomology is computed by the homological perturbation
lemma:
\[
m_3^{\mathrm{tr}}(a,b,c) =
\pi \circ m_2(h \circ m_2(i(a), i(b)),\, i(c))
+ \pi \circ m_2(i(a),\, h \circ m_2(i(b), i(c)))
+ \cdots,
\]
where $\pi$ is the projection to bar cohomology, $i$ is
the inclusion, and $h$ is the contracting homotopy.
For the Kac--Moody algebra, the inner $m_2$ produces
either a Lie bracket (from the $(0)$-product
$J^a_{(0)}J^b = f^{ab}{}_c J^c$) or a scalar (from the
$(1)$-product $J^a_{(1)}J^b = k\kappa^{ab}\mathbf{1}$).
When it produces a scalar, the outer $m_2$ acts on
$(\text{element}) \otimes \mathbf{1}$, which gives zero
in the reduced bar complex. When it produces a Lie
bracket, the homotopy $h$ maps this back to a bar element,
and the outer $m_2$ acts.

For affine algebras, the Jacobi identity $\ell_3 = 0$
implies that $m_3^{\mathrm{tr}}$ is a $d_2$-coboundary:
$m_3^{\mathrm{tr}} = d_2(\text{something})$. The key
point is that $m_3^{\mathrm{tr}}$ and the $L_\infty$
bracket $\ell_3$ are \emph{distinct} operations: the
former lives on bar cohomology while the latter is the
homotopy-transfer bracket. They are related by the
Koszul sign rule, and the Jacobi identity
\begin{equation}% label removed: eq:ell3-jacobi-vanishing
\ell_3(a, b, c) = [a, [b, c]] - [[a,b], c]
- (-1)^{|a||b|}[b,[a,c]] = 0
\end{equation}
forces $\ell_3 = 0$, which implies $[m_3^{\mathrm{tr}}]
= 0$ in $H^2$. Hence $o_4 = [m_3^{\mathrm{tr}}] = 0$
in $H^2(J^4(\widehat{\fg}_k), d_2)$.

\emph{Combining:}
The total quartic obstruction is
$o_4 = [\frac{1}{2}\{\mathfrak{C}, \mathfrak{C}\}_H
+ m_3^{\mathrm{tr}}] = 0$
in $H^2(J^4(\widehat{\fg}_k), d_2)$, by the Jacobi identity.

\medskip
\emph{Step~4: All higher obstructions vanish.}

By induction: the only source of higher obstructions is
the graph-sum composition of lower-degree vertices. Since
the only nontrivial vertex is the cubic
$m_2^{(0)} = [\;,\;]$ (the Lie bracket), every graph of
degree $r \geq 4$ is a tree with only trivalent vertices.
The amplitude of such a tree involves iterated Lie
brackets, contracted through the propagator. The Jacobi
identity, applied at each internal vertex, shows that
the amplitude of any such tree is a $d_2$-coboundary:
it is the $d_2$-image of a lower-weight correction.

Concretely: at degree~$r$, the obstruction is a sum over
trees with $r-1$ leaves and internal edges labeled by
the propagator $P$. Each internal vertex is a Lie
bracket. The Jacobi identity gives a recursive formula
\begin{equation}% label removed: eq:jacobi-recursive-vanishing
o_{r+1} = d_2\Bigl(
\sum_{\text{trees } T \in \mathcal{T}_{r}}
\frac{1}{|\Aut(T)|}\,\ell_T
\Bigr) = 0,
\end{equation}
where $\mathcal{T}_r$ is the set of binary trees with
$r$ leaves. The cohomology class vanishes because the
Jacobi identity is exactly the $A_\infty$ formality
condition for the Lie algebra (the Koszul resolution of
$\fg$ is formal).

Therefore $o_{r+1} = 0$ for all $r \geq 3$, and
the canonical Lie/tree support packet has
\(r^\Theta_{\max}(\widehat{\fg}_k) = 3\).
\end{proof}

\begin{remark}[The Jacobi identity as a termination criterion]
% label removed: rem:jacobi-termination
\index{Jacobi identity!tower termination}
The vanishing $o_4 = 0$ is not accidental but reflects
a deep structural property: the Lie bracket $[\;,\;]$
satisfies an associativity condition (the Jacobi identity)
that forces all compositions of Lie brackets to be
algebraically determined by the bracket itself. In the
language of operads, $\Lie$ is a Koszul operad, and
the bar resolution of any Lie algebra is formal. The
shadow obstruction tower terminates at degree~$3$ because there are
no higher obstructions to formality.
\end{remark}

% ================================================================
% 9.8.6: CLASS C; THE CONTACT/QUARTIC ARCHETYPE
% ================================================================

\subsection{Class~C: the contact/quartic archetype}
% label removed: subsec:class-c-contact
\index{shadow archetype!contact/quartic|textbf}
\index{beta-gamma system@$\beta\gamma$ system!shadow depth}

\begin{theorem}[$r^\Theta_{\max}(\beta\gamma) = 4$]
% label removed: thm:contact-rmax-four
\index{shadow depth!beta-gamma@$\beta\gamma$}
For the $\beta\gamma$ system, the support shadow depth is
\(r^\Theta_{\max}=4\). The obstruction classes satisfy:
\begin{equation}% label removed: eq:contact-obstruction-pattern
o_3(\beta\gamma) \neq 0,
\qquad
o_4(\beta\gamma) \neq 0,
\qquad
o_{r+1}(\beta\gamma) = 0
\;\;\text{for all } r \geq 4.
\end{equation}
In the canonical contact gauge slice, all higher lift coordinates
\(\lambda_n\in H^1(J^n(\beta\gamma),d_2)\) vanish for \(n>4\).
\end{theorem}

\begin{proof}
The proof uses the cubic gauge triviality theorem
(Theorem~\ref{thm:cubic-gauge-triviality}) and
rank-one abelian rigidity.

\medskip
\emph{Step~1: The OPE structure and transferred operations.}

The $\beta\gamma$ system has OPE
\begin{equation}% label removed: eq:betagamma-ope-recall
\beta(z)\,\gamma(w) \sim \frac{1}{z-w}.
\end{equation}
This is a simple-pole OPE with no double pole, so
the curvature $m_0 = 0$ and the Hessian of the
modular deformation complex is determined by the normal-ordered
products. The stress tensor
$T = {:}\partial\beta \cdot \gamma{:}$ gives the
conformal structure.

The key structural feature: the $\beta\gamma$ system
has rank one (a single pair of conjugate fields), so
the shadow algebra is one-dimensional at each degree.
The transferred operations are:
\begin{align*}
m_2^{\mathrm{tr}} &\neq 0
 & &\text{(cubic, from simple-pole OPE)}, \\
m_3^{\mathrm{tr}} &\neq 0
 & &\text{(quartic, from normal-ordered composites)}, \\
m_4^{\mathrm{tr}} &= 0
 & &\text{(quintic, vanishes by rank-one rigidity)}.
\end{align*}

\medskip
\emph{Step~2: Cubic gauge triviality.}

The cubic shadow at degree~$3$ is nonzero, but the
$\beta\gamma$ system satisfies the hypothesis of
Theorem~\ref{thm:cubic-gauge-triviality}:
\begin{equation}% label removed: eq:betagamma-h1-vanishing
H^1(F^3 \gAmod / F^4 \gAmod,\, d_2) = 0.
\end{equation}
This vanishing holds because the weight-$3$ piece of the
convolution algebra, for a rank-one system, is
one-dimensional and the $d_2$-action on it is injective
(the $\kappa$-twisted differential has no kernel at
degree~$1$ in weight~$3$). By
Theorem~\ref{thm:cubic-gauge-triviality}:
\begin{enumerate}[label=(\alph*)]
\item The cubic term $\Theta_3$ is gauge-trivial:
 there exists a gauge transformation removing it.
\item After this gauge choice, the quartic term
 defines a \emph{canonical class}
 $[\Theta'_4] \in H^2(J^4(\beta\gamma), d_2)$.
\end{enumerate}
This canonical quartic class is the \emph{quartic
contact shadow}
$\mathfrak{Q}^{\mathrm{contact}}_{\beta\gamma}$.

\medskip
\emph{Step~3: Quartic terminality.}

The quartic contact shadow is nontrivial:
$\mathfrak{Q}^{\mathrm{contact}}_{\beta\gamma} \neq 0$
(Corollary~\ref*{V1-cor:nms-betagamma-mu-vanishing}).
The key computation is:
\begin{equation}% label removed: eq:betagamma-mu-zero
\mu_{\beta\gamma}
:= \langle \eta, m_3(\eta, \eta, \eta) \rangle = 0,
\end{equation}
where $\eta$ is the generator of the shadow algebra.
Proof: this vanishing by a parity argument.
On the rank-one primary line spanned by~$\eta$,
$m_3(\eta,\eta,\eta)$ must be proportional to~$\eta$
by conformal weight and parity constraints. Write
$m_3(\eta,\eta,\eta) = \lambda\,\eta$. Then
$\mu = \langle \eta,\, \lambda\,\eta \rangle
= \lambda\,\langle \eta, \eta \rangle$.
But $m_3$ is graded-symmetric in its inputs (cubic
operation on identical generators), while the cyclic
pairing is graded-antisymmetric for odd-degree elements.
Since $|\eta|$ is odd (bar-desuspended),
$\langle \eta,\, m_3(\eta,\eta,\eta) \rangle
= -\langle m_3(\eta,\eta,\eta),\, \eta \rangle$
by antisymmetry, hence $\mu = -\mu$, so $\mu = 0$.

\medskip
\emph{Step~4: Quintic obstruction vanishes.}

The quintic obstruction is
\begin{equation}% label removed: eq:betagamma-quintic-obstruction
o_5(\beta\gamma) =
\bigl[\{\mathfrak{C}_{\beta\gamma},
\mathfrak{Q}^{\mathrm{contact}}_{\beta\gamma}\}_{H}
+ m_4^{\mathrm{tr}}\bigr].
\end{equation}
Verify both contributions vanish.

\emph{The boundary bracket.}
After the cubic gauge triviality (Step~2), the
gauge-fixed MC element has $\Theta_3' = 0$, so
$\mathfrak{C}_{\beta\gamma} = 0$ in the chosen gauge.
The boundary bracket $\{0, \mathfrak{Q}\}_H = 0$.

\emph{The local quintic vertex.}
The transferred quartic operation $m_4^{\mathrm{tr}} = 0$
by rank-one rigidity: on a one-dimensional primary space,
the cyclic $4$-ary operation is forced to vanish by the
same argument as in~\eqref{V1-eq:betagamma-mu-zero},
extended to degree~$4$.

Therefore $o_5 = 0$. By the same argument, $o_{r+1} = 0$
for all $r \geq 4$: the only source of higher obstructions
would be graph-sum compositions, but with $\Theta_3' = 0$
(in the chosen gauge) and $m_n^{\mathrm{tr}} = 0$ for
$n \geq 4$, no nontrivial graphs exist. Hence
\(r^\Theta_{\max}(\beta\gamma) = 4\).
\end{proof}

\begin{remark}[The $\beta\gamma$ system as the contact atom]
% label removed: rem:betagamma-contact-atom
\index{beta-gamma system@$\beta\gamma$ system!as contact atom}
The $\beta\gamma$ system is the unique standard atom
whose shadow obstruction tower is nontrivial beyond the scalar
level yet terminates at finite degree. The quartic
contact shadow is the simplest nonlinear modular
invariant beyond $\kappa$: it measures the
self-interaction of the normal-ordered composite
${:}\beta\gamma{:}$ with the stress tensor.
\end{remark}

% ================================================================
% 9.8.7: CLASS M; THE MIXED/INFINITE ARCHETYPE
% ================================================================

\subsection{Class~M: the mixed/infinite archetype}
% label removed: subsec:class-m-mixed
\index{shadow archetype!mixed/infinite|textbf}
\index{Virasoro algebra!shadow depth}

This subsection records the local computation behind the class
$\mathsf M$ example. The infinite-depth statement is the
all-degree Virasoro theorem of Volume~I; the calculation here exhibits
the quintic obstruction and the persistent cubic-source branch. The
argument combines the explicit
Virasoro OPE with the Gram matrix computation
(Proposition~\ref{V1-prop:gram-wt4}) and the graph-sum
formula for the obstruction.

\begin{theorem}[$r^\Theta_{\max}(\mathrm{Vir}_c) = \infty$]
% label removed: thm:virasoro-rmax-infinity
\index{shadow depth!Virasoro}
For the Virasoro algebra $\mathrm{Vir}_c$ at generic
central charge $c \neq 0, -22/5$, the support shadow depth is
\(r^\Theta_{\max} = \infty\). The quintic obstruction is
\begin{equation}% label removed: eq:virasoro-quintic-nonzero
o_5(\mathrm{Vir}_c) =
\frac{480}{c^2(5c+22)}\,x^5 \neq 0,
\end{equation}
and by induction all higher obstruction classes are
nonzero:
\begin{equation}% label removed: eq:virasoro-all-obstructions-nonzero
o_{r+1}(\mathrm{Vir}_c) \neq 0
\qquad \text{for all } r \geq 2.
\end{equation}
\end{theorem}

The proof requires several intermediate results, which
we develop in order.

\subsubsection{The Virasoro shadow obstruction tower through degree~$4$}
% label removed: subsubsec:virasoro-tower-through-4

\begin{computation}[Virasoro tower data]
% label removed: comp:virasoro-tower-data
\index{Virasoro algebra!tower data}
The Virasoro algebra $\mathrm{Vir}_c$ has a single strong
generator~$T$ of conformal weight~$2$. The OPE is
\begin{equation}% label removed: eq:virasoro-ope-full
T(z)\,T(w) \sim
\frac{c/2}{(z-w)^4}
+ \frac{2T(w)}{(z-w)^2}
+ \frac{\partial T(w)}{z-w}.
\end{equation}
The shadow obstruction tower data through degree~$4$, on the
one-dimensional primary space spanned by~$x$
(dual to~$T$):

\begin{center}
\renewcommand{\arraystretch}{1.5}
\begin{tabular}{cll}
\toprule
Degree $r$ & Shadow component & Formula \\
\midrule
$2$ & Hessian $H = \kappa$ & $c/2$ \\
$2$ & Propagator $P = H^{-1}$ & $2/c$ \\
$3$ & Cubic shadow $\mathfrak{C}$ & $2x^3$ \\
$3$ & Source: $T_{(1)}T = 2T$ & (simple-pole OPE) \\
$4$ & Quartic contact $\mathfrak{Q}^{\mathrm{ct}}$ &
 $\displaystyle\frac{10}{c(5c+22)}\,x^4$ \\
$4$ & Source: $\Lambda = {:}TT{:}
 - \tfrac{3}{10}\partial^2 T$ & (quasi-primary composite)\\
\bottomrule
\end{tabular}
\end{center}

\noindent The quartic contact coefficient
$Q_0 = 10/[c(5c+22)]$ is the inverse of the BPZ norm of
the quasi-primary composite
$\Lambda = {:}TT{:} - \frac{3}{10}\partial^2 T$
(Corollary~\ref{V1-cor:lambda-qp}).
\end{computation}

\begin{proof}[Derivation of the quartic contact]
The quartic contact coefficient is determined by the
condition that $\Theta^{\leq 4}$ satisfies the MC equation
through degree~$4$. The quartic MC equation requires
absorbing the boundary bracket and local vertex.

\emph{The boundary bracket.}
With $\mathfrak{C} = 2x^3$ and $P = 2/c$:
\begin{align}
\tfrac{1}{2}\{\mathfrak{C}, \mathfrak{C}\}_H
&= \tfrac{1}{2} \cdot
 \frac{\partial(2x^3)}{\partial x} \cdot P \cdot
 \frac{\partial(2x^3)}{\partial x} \notag \\
&= \tfrac{1}{2} \cdot 6x^2 \cdot \frac{2}{c}
 \cdot 6x^2 \notag \\
&= \frac{36}{c}\,x^4. % label removed: eq:virasoro-boundary-bracket
\end{align}

\emph{The local quartic vertex.}
The quasi-primary composite
$\Lambda = {:}TT{:} - \frac{3}{10}\partial^2 T$
has BPZ norm
$\langle \Lambda | \Lambda \rangle = c(5c+22)/10$
(Corollary~\ref{V1-cor:lambda-qp}).
The local contribution is proportional to $1/\langle\Lambda|\Lambda\rangle$:
\begin{equation}% label removed: eq:virasoro-local-quartic
\mathfrak{Q}^{\mathrm{local}}
= \frac{4}{\langle\Lambda|\Lambda\rangle}\,x^4
= \frac{40}{c(5c+22)}\,x^4.
\end{equation}

\emph{The total quartic.}
By the all-degree master equation
$\nabla_H(\mathrm{Sh}_4) + \mathfrak{o}^{(4)} = 0$
(Volume~I Theorem~\ref{V1-thm:nms-all-degree-master-equation}),
the quartic shadow absorbs the total obstruction, giving:
\begin{equation}% label removed: eq:virasoro-contact-result
Q^{\mathrm{contact}}_{\mathrm{Vir}}
= \frac{10}{c(5c+22)}.
\end{equation}
\end{proof}

\subsubsection{The quintic obstruction at the chain level}
% label removed: subsubsec:quintic-chain-level

\begin{proof}[Local computation for Volume~I Theorem~\textup{\ref{V1-thm:virasoro-rmax-infinity}}]
The computation gives the quintic obstruction $o_5$ as a nonzero
class in $H^2(J^5(\mathrm{Vir}_c), d_2)$.

\medskip
\emph{Step~1: The quintic MC equation.}

The projection of the MC equation
$D_\cA \Theta + \frac{1}{2}[\Theta, \Theta] = 0$
onto degree~$5$ gives
\begin{equation}% label removed: eq:quintic-mc-equation
d_2 \Theta_5 +
\{\mathfrak{C}, \mathfrak{Q}^{\mathrm{ct}}\}_H
+ \text{(higher-order contributions)} = 0.
\end{equation}
The leading obstruction is the cubic-quartic bracket.

\medskip
\emph{Step~2: The cubic-quartic bracket.}

\begin{align}
\{\mathfrak{C}, \mathfrak{Q}^{\mathrm{ct}}\}_H
&= \frac{\partial(2x^3)}{\partial x}
 \cdot P \cdot
 \frac{\partial(Q_0\,x^4)}{\partial x} \notag \\
&= 6x^2 \cdot \frac{2}{c} \cdot 4Q_0\,x^3 \notag \\
&= \frac{48Q_0}{c}\,x^5 \notag \\
&= \frac{48}{c} \cdot \frac{10}{c(5c+22)}\,x^5 \notag \\
&= \frac{480}{c^2(5c+22)}\,x^5.
 % label removed: eq:quintic-cubic-quartic
\end{align}

\medskip
\emph{Step~3: Nonvanishing.}

The coefficient $480/[c^2(5c+22)]$ is nonzero for
generic $c$ (failing only at $c = 0$ and $c = -22/5$,
where the Virasoro algebra degenerates). We must verify
that $\{\mathfrak{C}, \mathfrak{Q}^{\mathrm{ct}}\}_H$
is \emph{not} a $d_2$-coboundary. The $H$-Poisson
bracket $[\mathfrak{C}, \mathrm{Sh}_r]_H$ lies in
$\cA^{\mathrm{sh}}_{r+1,0}$. If it were exact, there
would exist $\xi \in \cA^{\mathrm{sh}}_{r+1,-1}$ with
$d_2(\xi) = [\mathfrak{C}, \mathrm{Sh}_r]_H$. But
$\dim \cA^{\mathrm{sh}}_{r+1,-1} = 0$ for $r \geq 4$:
the weight-$({-}1)$ space is trivial in the cyclic
deformation complex at high degree (there are no degree-$1$
cochains in $J^{r+1}$ for $r \geq 4$ on the one-generator
primary line, since a single generator of conformal
weight~$2$ produces only even-weight monomials in the
shadow algebra). Hence the bracket
$[\mathfrak{C}, \mathrm{Sh}_r]_H$ is automatically a
nontrivial cohomology class.

Therefore
\begin{equation}% label removed: eq:quintic-obstruction-clean
o_5(\mathrm{Vir}_c) =
\frac{480}{c^2(5c+22)}\,x^5 \neq 0.
\end{equation}

\emph{Step~4: The quintic shadow.}
The all-degree master equation
$\nabla_H(\mathrm{Sh}_5) + o^{(5)} = 0$
(Volume~I Theorem~\ref{V1-thm:nms-all-degree-master-equation})
determines $\mathrm{Sh}_5$. On the one-generator primary
line with Hessian $H = c/2$, the covariant derivative
acts as $\nabla_H(S_r x^r) = 2rS_r x^r$. Therefore
$\mathrm{Sh}_5 = -o^{(5)}/(2 \cdot 5)$:
\begin{equation}% label removed: eq:sh5-virasoro-formula
\mathrm{Sh}_5(\mathrm{Vir}_c)
= -\frac{480}{c^2(5c+22)} \cdot \frac{1}{10}\,x^5
= -\frac{48}{c^2(5c+22)}\,x^5,
\end{equation}
in exact agreement with Volume~I
Corollary~\ref{V1-cor:virasoro-quintic-shadow-explicit}.

\medskip
\emph{Step~5: Cubic-source branch and the all-degree input.}

The quintic shadow
$\mathrm{Sh}_5 = -48/(c^2(5c+22))\,x^5 \neq 0$
acts as a new vertex in the graph sum at degree~$6$.
The sextic obstruction includes the cubic-quintic bracket:
\begin{align}
\{\mathfrak{C}, \mathrm{Sh}_5\}_H
&= 6x^2 \cdot \frac{2}{c} \cdot
 5\cdot\bigl(-\frac{48}{c^2(5c+22)}\bigr)\,x^4
 \notag \\
&= -\frac{2880}{c^3(5c+22)}\,x^6
\neq 0. % label removed: eq:sextic-leading-term
\end{align}
By induction in the cubic-source approximation: at each degree~$r$, the
cubic vertex
$\mathfrak{C} = 2x^3$ contracts with the $(r{-}1)$-st
shadow contribution to produce a nonzero $r$-th source:
\begin{equation}% label removed: eq:inductive-obstruction
o_{r+1}(\mathrm{Vir}_c) \ni
\{\mathfrak{C}, \mathrm{Sh}_r\}_H
= 6x^2 \cdot \frac{2}{c} \cdot r\,S_r\,x^{r-1}
= \frac{12r\,S_r}{c}\,x^{r+1},
\end{equation}
where $S_r$ denotes the coefficient of this cubic-source contribution.
This source is nonzero whenever $S_r \neq 0$. Since $S_5 \neq 0$, the
cubic-source branch starts, and its recursion gives
\begin{equation}% label removed: eq:shadow-recursive-formula
S_{r+1}^{\mathrm{cub}} =
-\frac{6r}{c(r+1)}\,S_r^{\mathrm{cub}},
\qquad r \geq 5.
\end{equation}
Since $6r/(c(r+1)) \neq 0$ for generic $c$ and all
$r \geq 5$, the cubic-source contribution is nonzero at all higher
degrees. The passage from this persistent source to exact nonvanishing
of every obstruction class uses the all-degree non-cancellation theorem
of Volume~I, Theorem~\ref{V1-thm:virasoro-rmax-infinity}. Therefore
\(r^\Theta_{\max}(\mathrm{Vir}_c)=\infty\).
\end{proof}

\begin{remark}[The cubic vertex as permanent source]
% label removed: rem:cubic-permanent-source
\index{Virasoro algebra!cubic as permanent source}
The essential mechanism is that the cubic vertex
$\mathfrak{C} = 2x^3$ acts as a \emph{permanent source}
in the shadow obstruction tower: its bracket with any nonzero shadow
produces a nonzero obstruction at the next level. This
is specific to the mixed archetype, where both
$\mathfrak{C} \neq 0$ and $\mathfrak{Q}^{\mathrm{ct}}
\neq 0$ coexist. In the Lie/tree archetype (Kac--Moody),
$\mathfrak{C} \neq 0$ but $\mathfrak{Q}^{\mathrm{ct}} = 0$,
so the quartic bracket $\{\mathfrak{C}, \mathfrak{C}\}_H$
is $d_2$-exact (by Jacobi) and no quintic is generated.
In the contact archetype ($\beta\gamma$),
$\mathfrak{C} = 0$ in the appropriate gauge, so the
bracket with $\mathfrak{Q}^{\mathrm{ct}}$ vanishes.
The coexistence of both nonzero $\mathfrak{C}$ and
$\mathfrak{Q}^{\mathrm{ct}}$ is the signature of
class~\textbf{M}.
\end{remark}

% ================================================================
% 9.8.8: VIRASORO JET SPACES
% ================================================================

\subsection{Virasoro jet spaces}
% label removed: subsec:virasoro-jet-spaces
\index{Virasoro algebra!jet spaces|textbf}
\index{shadow jet space!Virasoro}

The shadow jet space $J^r(\mathrm{Vir}_c)$ is the
degree-$r$ graded piece of the modular convolution algebra.
For the Virasoro algebra, the one-generator primary space
makes the jet spaces particularly tractable: $J^r$ is
identified with the space of degree-$r$ homogeneous
polynomials in the variable~$x$ (dual to the single
generator~$T$), twisted by the BPZ inner product
structure at each conformal weight.

The full jet space includes contributions from
\emph{all} conformal weights: descendants and
quasi-primary composites. At degree~$r$, the relevant
conformal weights range from $2r$ (the minimum, from
$T^r$) upward. The number of states at weight~$h$
in the Virasoro vacuum Verma module is the
partition function $p(h - 2)$ for $h \geq 2$.

\begin{computation}[Virasoro jet space dimensions]
% label removed: comp:virasoro-jet-dimensions
\index{Virasoro algebra!jet space dimensions}
The dimension of the shadow jet space $J^r(\mathrm{Vir}_c)$
on the primary line is $1$ for all $r \geq 2$ (since
there is one generator). The dimension of the
\emph{extended} jet space, including descendants and
quasi-primary composites at total conformal weight~$2r$,
is computed from the Verma module structure.

At degree~$r$ and total weight $w = 2r + d$ (excess~$d$),
each factor is a state in the vacuum Verma module at
weight $h_i \geq 2$ with $\sum h_i = w$. The states
at weight~$h$ number $p(h-2)$. The extended jet space
dimension at excess $d = 0$ is $1$; at excess $d = 1$ it
is $r$ (choice of which factor descends); at excess
$d = 2$ it is $r + \binom{r}{2} = r(r+1)/2$.

The first several values at excess $d = 0$:

\begin{center}
\renewcommand{\arraystretch}{1.4}
\begin{tabular}{crccl}
\toprule
$r$ & $\dim J^r$ & $\dim J^r_{\mathrm{ext}}(d{\leq}2)$
 & $\sum_{d=0}^{2r}$ & Obstruction \\
\midrule
$2$ & $1$ & $1 + 2 + 3 = 6$ & $7$
 & $\kappa = c/2$ \\
$3$ & $1$ & $1 + 3 + 6 = 10$ & $15$
 & $\mathfrak{C} = 2x^3$ \\
$4$ & $1$ & $1 + 4 + 10 = 15$ & $26$
 & $Q^{\mathrm{ct}} = 10/(c(5c{+}22))$ \\
$5$ & $1$ & $1 + 5 + 15 = 21$ & $42$
 & $\mathrm{Sh}_5 = -48/(c^2(5c{+}22))$ \\
$6$ & $1$ & $1 + 6 + 21 = 28$ & $63$
 & $\mathrm{Sh}_6 \neq 0$ \\
$7$ & $1$ & $1 + 7 + 28 = 36$ & $92$
 & $\mathrm{Sh}_7 \neq 0$ \\
$8$ & $1$ & $1 + 8 + 36 = 45$ & $131$
 & $\mathrm{Sh}_8 \neq 0$ \\
\bottomrule
\end{tabular}
\end{center}

\noindent The column $\dim J^r_{\mathrm{ext}}(d{\leq}2)$
counts states through conformal-weight excess~$2$.
\end{computation}

\begin{remark}[Shapovalov determinant and jet space
structure]
% label removed: rem:shapovalov-jet
\index{Shapovalov determinant!jet space structure}
The Shapovalov determinant at weight~$h$ in the Virasoro
vacuum Verma module is
\begin{equation}% label removed: eq:shapovalov-virasoro
\det S_h = C_h \prod_{\substack{r,s \geq 1 \\ rs \leq h}}
(h_{r,s}(c))^{p(h-rs)},
\end{equation}
where $h_{r,s}(c)$ are the Kac table entries
(parametrized by
$c = 13 - 6(t + t^{-1})$, $h_{r,s} =
((rt - s/t)^2 - (t - 1/t)^2)/4$).

The Shapovalov determinant controls the singularity
structure of the jet spaces: at central charges where
$\det S_h = 0$, the Verma module develops singular
vectors, and the quasi-primary composites at that weight
become linearly dependent. The quartic contact
coefficient $Q^{\mathrm{ct}} = 10/(c(5c+22))$ has
poles at $c = 0$ and $c = -22/5$, which are exactly
the zeros of the weight-$4$ Shapovalov determinant
$\det G = c^2(5c+22)/2$
(Proposition~\ref{V1-prop:gram-wt4}).
\end{remark}

% ================================================================
% 9.8.9: GROWTH RATE THEOREM
% ================================================================

\subsection{Growth rate of the shadow obstruction tower}
% label removed: subsec:growth-rate
\index{shadow tower!growth rate|textbf}
\index{growth rate!shadow jet spaces}

\begin{theorem}[Polynomial growth of Virasoro jet spaces]
% label removed: thm:jet-space-polynomial-growth
\index{Virasoro algebra!jet space growth rate}
For the Virasoro algebra at generic central charge,
the extended jet space dimensions grow
polynomially in~$r$:
\begin{equation}% label removed: eq:jet-polynomial-growth
\dim J^r_{\mathrm{ext}}(\mathrm{Vir}_c; d \leq D)
= O(r^D)
\qquad \text{as } r \to \infty,
\end{equation}
for each fixed conformal-weight excess~$D$.
On the primary line, $\dim J^r = 1$ for all $r$.
\end{theorem}

\begin{proof}
At conformal-weight excess $d = 0$: $\dim J^r = 1$
(the primary sector, spanned by $x^r$).

At excess $d = 1$: one factor carries
$L_{-1}$ descent (weight $h_i = 3$, one state), the
remaining $r - 1$ factors are at weight~$2$. The number
of choices is $r$, giving $\dim = r$.

At excess $d = 2$: either one factor at weight~$4$
($p(2) = 2$ states: $L_{-2}|0\rangle$ and $L_{-1}^2|0\rangle$)
and $r - 1$ at weight~$2$, giving $2r$ states; or
two factors at weight~$3$ ($p(1) = 1$ state each)
and $r - 2$ at weight~$2$, giving $\binom{r}{2}$ states.
Total: $2r + \binom{r}{2} = r(r+3)/2$.

In general, at excess~$d$ the dimension is a polynomial
in~$r$ of degree at most~$d$: the leading term is
$\binom{r}{d}/d!$ from distributing $d$ units of excess
among $r$ factors. By the multinomial theorem:
\begin{equation}% label removed: eq:extended-dimension-polynomial
\dim J^r_{\mathrm{ext}}(d)
= \sum_{\substack{\lambda_1 + \cdots + \lambda_r = d \\
\lambda_i \geq 0}}
\prod_{i=1}^r p(\lambda_i)
\sim c_d\,r^d,
\end{equation}
where $c_d = \sum_{|\lambda| = d} \prod p(\lambda_i) / d!$
is a combinatorial constant depending on the partition
function.

For the cumulative dimension through excess~$D$:
$\dim J^r_{\mathrm{ext}}(d \leq D) = \sum_{d=0}^D
c_d r^d + \text{lower-order terms} = O(r^D)$.
\end{proof}

\begin{remark}[Information-theoretic content]
% label removed: rem:information-content
\index{information theory!shadow depth}
The polynomial growth of $\dim J^r_{\mathrm{ext}}$
has an information-theoretic interpretation. To specify the
holographic reconstruction of $\cA$ through degree~$r$,
one needs $\dim J^r$ real parameters on the primary line.
For Gaussian algebras, $\dim J^r = 1$ for all $r$ (the
single parameter $\kappa$ determines the rank-one Gaussian shadow,
not the full tensor-channel package). For
Virasoro, each degree adds one primary-line parameter,
so $N(r) = r - 1$ through degree~$r$.

The \emph{total} number of parameters including descendants
is $\dim J^r_{\mathrm{ext}}$, which grows polynomially.
The ratio $\dim J^r_{\mathrm{ext}} / r$ measures the
\emph{marginal complexity}: the number of new parameters
per unit increase in degree.
\end{remark}

% ================================================================
% 9.8.10: RECURSIVE STRUCTURE OF THE SHADOW COEFFICIENTS
% ================================================================

\subsection{Recursive structure of the Virasoro
shadow coefficients}
% label removed: subsec:virasoro-shadow-recursive
\index{Virasoro algebra!shadow coefficients!recursive structure}

The recursive formula~\eqref{V1-eq:shadow-recursive-formula}
determines the cubic-source contribution to the shadow coefficients for
$r \geq 5$. We record the first several values and exhibit the
closed-form solution for that contribution.

\begin{proposition}[Virasoro shadow coefficients,
cubic-source approximation]
% label removed: prop:virasoro-shadow-coefficients
\index{shadow coefficients!Virasoro}
The primary-line shadow coefficients
$\mathrm{Sh}_r(\mathrm{Vir}_c) = S_r\,x^r$
have the following dominant contribution from
the cubic source term in the shadow master equation.
The cubic-source recursion gives
\begin{equation}% label removed: eq:virasoro-shadow-recursion
S_{r+1}^{\mathrm{cub}} = -\frac{6r}{c(r+1)}\,S_r^{\mathrm{cub}},
\qquad r \geq 5,
\end{equation}
with initial conditions from the explicit tower data:
\begin{align}
S_2 &= c/2, % label removed: eq:s2 \\
S_3 &= 2, % label removed: eq:s3 \\
S_4 &= 10/(c(5c+22)), % label removed: eq:s4 \\
S_5 &= -48/(c^2(5c+22)). % label removed: eq:s5
\end{align}
For $r \geq 5$, the cubic-source closed form is
\begin{equation}% label removed: eq:virasoro-shadow-closed-form
S_r^{\mathrm{cub}} =
\frac{(-1)^{r-4} \cdot 6^{r-5} \cdot 240}
{c^{r-3}(5c+22) \cdot r}.
\end{equation}
For $r \leq 5$ this agrees with the exact shadow
coefficients: $S_r^{\mathrm{cub}} = S_r$.
For $r \geq 6$, additional contributions from
$\{\mathrm{Sh}_j, \mathrm{Sh}_k\}$ bracket terms with
$j,k \geq 4$ enter the full recursion, and the exact
values differ; see
Theorem~\textup{\ref{thm:thqg-virasoro-tower-explicit}}
and~\S\textup{\ref{subsec:thqg-virasoro-tower}}.
\end{proposition}

\begin{proof}
The cubic-source recursion
$S_{r+1}^{\mathrm{cub}} = -\frac{6r}{c(r+1)}S_r^{\mathrm{cub}}$
is iterated from $r = 5$ to $r - 1$:
\begin{align*}
S_r^{\mathrm{cub}} &= \prod_{j=5}^{r-1}\bigl(-\frac{6j}{c(j+1)}\bigr)
 \cdot S_5 \\
&= \frac{(-6)^{r-5}}{c^{r-5}}
 \cdot \frac{\prod_{j=5}^{r-1} j}{\prod_{j=5}^{r-1}(j+1)}
 \cdot S_5 \\
&= \frac{(-6)^{r-5}}{c^{r-5}}
 \cdot \frac{5}{r}
 \cdot \bigl(-\frac{48}{c^2(5c+22)}\bigr) \\
&= \frac{(-1)^{r-4} \cdot 6^{r-5} \cdot 240}
 {c^{r-3}(5c+22) \cdot r}.
\end{align*}
\emph{Verification at $r = 5$:}
$(-1)^1 \cdot 6^0 \cdot 240 / (c^2(5c+22) \cdot 5)
= -48/(c^2(5c+22)) = S_5$. \checkmark

\emph{Verification at $r = 6$:}
$(-1)^2 \cdot 6 \cdot 240 / (c^3(5c+22) \cdot 6)
= 240/(c^3(5c+22)) = S_6^{\mathrm{cub}}$. \checkmark

Cross-check via cubic-source recursion:
$S_6^{\mathrm{cub}} = -\frac{6 \cdot 5}{c \cdot 6} S_5
= -\frac{5}{c} \cdot (-\frac{48}{c^2(5c+22)})
= \frac{240}{c^3(5c+22)}$. \checkmark

\smallskip
\noindent\textit{Caveat.}
The exact $S_6$ receives an additional
$\{\mathrm{Sh}_4, \mathrm{Sh}_4\}$ cross-term not captured
by this recursion. The cubic-source formula gives the
dominant contribution for $|c| \gg 1$ but is not exact for
$r \geq 6$; see
Theorem~\ref{thm:thqg-virasoro-tower-explicit}.
\end{proof}

\begin{corollary}[Cubic-source asymptotic shadow decay]
% label removed: cor:shadow-asymptotic-decay
\index{shadow coefficients!asymptotic decay}
The Virasoro shadow coefficients satisfy, in the
cubic-source approximation,
\begin{equation}% label removed: eq:shadow-asymptotic
|S_r^{\mathrm{cub}}| \sim \frac{240}{|c|^{r-3}|5c+22| \cdot r}
\cdot 6^{r-5}
\qquad \text{as } r \to \infty.
\end{equation}
For the cubic-source approximation, the ratio is
$|S_{r+1}^{\mathrm{cub}}|/|S_r^{\mathrm{cub}}| \to 6/|c|$.
The exact full tower receives corrections from the quartic and higher
bracket terms and must be read from the full recursion. Thus the
cubic-source convergence trichotomy is:
\begin{itemize}
\item For $|c| > 6$: the cubic-source series is geometrically
 decaying.
\item For $|c| < 6$: the cubic-source series is formally divergent;
 the completed obstruction tower remains an inverse-limit object in the
 completed convolution algebra topology
 \textup{(}Theorem~\ref{V1-thm:completed-bar-cobar-strong}\textup{)}.
\item For $|c| = 6$: borderline, with
 $|S_r^{\mathrm{cub}}| \sim C/r$.
\end{itemize}
\end{corollary}

\begin{remark}[Convergence and the analytic sewing programme]
% label removed: rem:convergence-sewing
\index{analytic sewing!shadow tower convergence}
The formal divergence of the cubic-source series for $|c| < 6$ does not
obstruct the algebraic theory: $\Theta_\cA$ is still defined as a
compatible element of the completed convolution algebra. For the
analytic sewing programme, convergence is an extra hypothesis: the
shadow tower must converge in an appropriate operator norm to yield
finite partition functions. The HS-sewing condition
(Theorem~\ref{thm:general-hs-sewing}), when its hypotheses hold,
provides the required estimate.
\end{remark}

% ================================================================
% 9.8.11: EXPLICIT VIRASORO SHADOW TABLE
% ================================================================

\subsection{Explicit Virasoro shadow table}
% label removed: subsec:virasoro-shadow-table
\index{Virasoro algebra!shadow table}

\begin{computation}[Virasoro shadow coefficients through
degree~$10$, cubic-source approximation]
% label removed: comp:virasoro-shadow-table
\index{Virasoro algebra!shadow coefficients!table}

The following table records the cubic-source
approximation~$S_r^{\mathrm{cub}}$ from
Proposition~\ref{V1-prop:virasoro-shadow-coefficients}.
For $r \leq 5$, these are the exact shadow coefficients.
For $r \geq 6$, the exact values receive corrections from
$\{\mathrm{Sh}_j, \mathrm{Sh}_k\}$ cross-terms with
$j,k \geq 4$; see
\S\ref{subsec:thqg-virasoro-tower} for exact values.

\begin{center}
\renewcommand{\arraystretch}{1.5}
\small
\begin{tabular}{cll}
\toprule
$r$ & $S_r^{\mathrm{cub}}$ & Value at $c = 26$ \\
\midrule
$2$ & $c/2$ & $13$ \\
$3$ & $2$ & $2$ \\
$4$ & $\displaystyle\frac{10}{c(5c+22)}$
 & $\displaystyle\frac{10}{3952}
 \approx 2.53 \times 10^{-3}$ \\
$5$ & $\displaystyle\frac{-48}{c^2(5c+22)}$
 & $\approx -4.67 \times 10^{-4}$ \\
$6^{\dagger}$ & $\displaystyle\frac{240}{c^3(5c+22)}$
 & $\approx 9.01 \times 10^{-5}$ \\
$7^{\dagger}$ & $\displaystyle\frac{-1440}{c^4(5c+22) \cdot 7/6}$
 & $\approx -1.48 \times 10^{-5}$ \\
$8^{\dagger}$ & $\displaystyle\frac{(-6)^3 \cdot 240}
 {c^5(5c+22) \cdot 8}$
 & $\approx 2.14 \times 10^{-6}$ \\
$9^{\dagger}$ & $\displaystyle\frac{(-6)^4 \cdot 240}
 {c^6(5c+22) \cdot 9}$
 & $\approx -2.77 \times 10^{-7}$ \\
$10^{\dagger}$ & $\displaystyle\frac{(-6)^5 \cdot 240}
 {c^7(5c+22) \cdot 10}$
 & $\approx 3.26 \times 10^{-8}$ \\
\bottomrule
\end{tabular}
\end{center}

\noindent{\footnotesize ${}^{\dagger}$Cubic-source
approximation; exact values for $r \geq 6$ differ
by $\{\mathrm{Sh}_4, \mathrm{Sh}_4\}$ and higher
cross-term corrections.}

\smallskip
\noindent At $c = 26$ (the bosonic string), the cubic-source
coefficients decay geometrically with ratio
$\approx 6/26 = 3/13 \approx 0.231$. No analytic convergence statement
for the exact tail is asserted without the sewing norm.
\end{computation}

\begin{remark}[The bosonic string cubic-source decay]
% label removed: rem:bosonic-string-convergent
\index{bosonic string!shadow tower}
At $c = 26$, the Koszul dual is
$\mathrm{Vir}_{26}^! = \mathrm{Vir}_0$, and the
complementarity pairing gives
$\kappa + \kappa' = 13$, with $\kappa = 13$, $\kappa' = 0$.
The cubic-source approximation decays rapidly; the exact completed
shadow tower remains governed by the full recursion and the analytic
sewing norm. This is the \emph{depth-zero resonance shadow}
(Remark~\ref{V1-rem:virasoro-resonance-model}).
\end{remark}

% ================================================================
% 9.8.12: DICTIONARY OF RECONSTRUCTION INVARIANTS
% ================================================================

\subsection{Dictionary of reconstruction invariants}
% label removed: subsec:reconstruction-dictionary
\index{reconstruction invariants!dictionary}

\begin{center}
\renewcommand{\arraystretch}{1.5}
\small
\begin{tabular}{ccccc}
\toprule
& \textbf{G} & \textbf{L} & \textbf{C} & \textbf{M} \\
& (Gaussian) & (Lie/tree) & (Contact) & (Mixed) \\
\midrule
$r^\Theta_{\max}$ & $2$ & $3$ & $4$ & $\infty$ \\
Atom & $\mathcal{H}_k$ & $\widehat{\fg}_k$
 & $\beta\gamma$ & $\mathrm{Vir}_c$ \\
$\Theta$ & $\kappa\,x^2$ & $+\,\mathfrak{C}\,x^3$
 & $+\,\mathfrak{Q}\,x^4$ & $+\sum_{r \geq 5} S_r\,x^r$ \\
\midrule
OPE poles & $\leq 2$ & $\leq 2$ & $\leq 1$ & $\leq 4$ \\
$m_n^{\mathrm{tr}} = 0$ & $n \geq 3$ & $n \geq 3$
 & $n \geq 5$ & never \\
Termination & Formality & Jacobi & Rank-1 rigidity
 & N/A \\
\midrule
Deformation & Smooth & Smooth & Smooth & Formal $\widehat{\phantom{x}}$ \\
Rep.\ theory & Fock & Verma irred.\ & Rank-1 Fock
 & Verma singular \\
Graph sums & 1 graph & trees & 1 graph + gauge
 & all graphs \\
Growth & $O(1)$ & $O(1)$ & $O(1)$
 & $O(r^D)$ \\
\bottomrule
\end{tabular}
\end{center}

\begin{remark}[Operadic complexity theorem]
% label removed: rem:operadic-complexity-holographic
\index{operadic complexity!holographic interpretation}
The operadic complexity theorem
(Vol~I, Theorem~\ref*{V1-thm:operadic-complexity-detailed})
states that
\(r^\Theta_{\max}(\cA)\) equals the support depth of the transferred
\(A_\infty\)-structure of~\(\cA\):
the minimal $n$ such that the transferred
$A_\infty$-structure on $H^*(\barB(\cA))$ has
$m_k^{\mathrm{tr}} = 0$ for all $k > n$. Equivalently,
\(r^\Theta_{\max}\) is the last nonzero homogeneous component of the
bar-intrinsic \(L_\infty\) deformation packet.  The obstruction depth
\(d_{\mathrm{obs}}\) records only eventual \(H^2\)-vanishing and may
be smaller. The proof identifies the
genus-$0$ shadow graph sum with the homotopy transfer
tree formula at all degrees. In the
holographic language: the gravitational complexity of the
bulk theory equals the homotopy-algebraic complexity of
the boundary algebra.
\end{remark}

% ================================================================
% 9.8.13: OBSTRUCTION TOWER FOR GENERAL W-ALGEBRAS
% ================================================================

\subsection{The obstruction tower for $\Walg_N$-algebras}
% label removed: subsec:wn-obstruction-tower
\index{W-algebra@$\Walg$-algebra!obstruction tower}

\begin{proposition}[$\Walg_N$ is class~\textbf{M} for
$N \geq 2$]
% label removed: prop:wn-class-m
\index{shadow depth!$\Walg_N$}
For the principal $\Walg_N$-algebra at generic level, read in the
generic completed ambient,
\(r^\Theta_{\max}(\Walg_N) = \infty\) for all $N \geq 2$.
The obstruction at degree~$5$ contains a nonzero local source from
the coexistence of the gravitational cubic
(from $T_{(1)}T = 2T$) and the quartic contact
(from $\Lambda = {:}TT{:} - \frac{3}{10}\partial^2 T$
and higher quasi-primaries).
\end{proposition}

\begin{proof}
The gravitational cubic $\mathfrak{C} \neq 0$ (from
$T_{(1)}T = 2T$, present for all $N$) and the quartic
contact $\mathfrak{Q}^{\mathrm{ct}} \neq 0$
(Theorem~\ref{V1-thm:w-principal-wn-contact-nonvanishing})
give a nonzero cubic-quartic bracket
$\{\mathfrak{C}, \mathfrak{Q}^{\mathrm{ct}}\}_H \neq 0$.
Exact non-cancellation and the all-degree infinite-depth statement are
read from the generic completed $\Walg_N$ recursion, with the Virasoro
stress-tensor subcomplex governed by
Theorem~\ref{V1-thm:virasoro-rmax-infinity}.
\end{proof}

\begin{remark}[The $\Walg_\infty$ limit]
% label removed: rem:winfty-limit-reconstruction
\index{W-infinity@$\Walg_\infty$!reconstruction}
At the $\Walg_\infty$ limit ($N \to \infty$), the
shadow obstruction tower is infinite at each finite stage but
one-dimensional at the minimal cyclic level at the limit
(Theorem~\ref{V1-thm:winfty-scalar}): the
Mittag--Leffler argument forces all shadow projections
into the one-dimensional $H^2_{\mathrm{cyc}}$.
This identifies only the minimal cyclic direction; it does not imply
that a single number $\kappaChHodge(\Walg_\infty)$ determines the entire
modular Koszul package.
\end{remark}

% ================================================================
% 9.8.14: THE POSTNIKOV FILTRATION
% ================================================================

\subsection{Obstruction tower and Postnikov filtration}
% label removed: subsec:obstruction-tower
\index{obstruction tower|textbf}
\index{Postnikov filtration!shadow tower}

\begin{definition}[Postnikov $k$-invariant of the
shadow obstruction tower]
% label removed: def:shadow-k-invariant
\index{Postnikov $k$-invariant!shadow tower}
The \emph{$r$-th $k$-invariant} of the shadow obstruction tower is
the obstruction class
\begin{equation}% label removed: eq:shadow-k-invariant
k_r(\cA) := o_{r+1}(\cA) \in H^2(J^{r+1}(\cA), d_2).
\end{equation}
The shadow obstruction tower is
\emph{$n$-truncated} if $k_r(\cA) = 0$ for all
$r \geq n$, i.e., if \(d_{\mathrm{obs}}(\cA)\le n\).
\end{definition}

\begin{proposition}[Structure of the Postnikov filtration]
% label removed: prop:postnikov-filtration-structure
\index{Postnikov filtration!structural properties}
\begin{enumerate}[label=\textup{(\roman*)}]
\item \emph{Obstruction profile.}
The $k$-invariants $\{k_r(\cA)\}_{r \geq 2}$ determine the
obstruction-depth profile \(d_{\mathrm{obs}}(\cA)\).  The MC gauge
class \([\Theta_\cA]\) is determined only after adjoining the
\(H^1\)-lift coordinates of the shadow reconstruction packet.

\item \emph{Functoriality.}
A quasi-isomorphism $\cA \xrightarrow{\sim} \cA'$ induces
isomorphisms $k_r(\cA) \cong k_r(\cA')$ for all $r$.

\item \emph{Additivity.}
For a direct sum $\cA = \cA_1 \oplus \cA_2$ with
vanishing mixed OPE:
$k_r(\cA) = k_r(\cA_1) + k_r(\cA_2)$ for all $r$.

\item \emph{Finiteness.}
Each $H^2(J^{r+1}(\cA), d_2)$ is finite-dimensional
for all modular Koszul chiral algebras with
finitely many strong generators.
\end{enumerate}
\end{proposition}

\begin{proof}
Part~(i): by definition \(k_r=o_{r+1}\), so the sequence
\(\{k_r\}\) is exactly the \(H^2\)-obstruction profile.  The
obstruction-extension sequence shows that a vanishing \(k_r\) leaves an
\(H^1(J^{r+1}(\cA),d_2)\)-torsor of lifts; therefore the full MC gauge
class requires the lift coordinates.  The bar-intrinsic branch with
those coordinates is the complete system of invariants by the
branch-level Mittag--Leffler corollary above.
Part~(ii):
Theorem~\ref{V1-thm:shadow-homotopy-invariance}.
Part~(iii):
Proposition~\ref{V1-prop:independent-sum-factorization}.
Part~(iv): finite-dimensionality of the bar complex at
fixed degree and conformal weight.
\end{proof}

% ================================================================
% 9.8.15: DEFORMATION-THEORETIC INTERPRETATION
% ================================================================

\subsection{Deformation-theoretic interpretation}
% label removed: subsec:deformation-interpretation
\index{deformation theory!shadow tower interpretation}

\begin{proposition}[The MC moduli as a formal moduli problem]
% label removed: prop:mc-formal-moduli
\index{formal moduli!MC moduli}
The assignment
\begin{equation}% label removed: eq:mc-formal-moduli
R \;\mapsto\; \MC(\gAmod \widehat{\otimes} \mathfrak{m}_R)
/ \mathrm{gauge}
\end{equation}
(where $R$ is an Artinian dg algebra and $\mathfrak{m}_R$
its maximal ideal) defines a formal moduli problem
whose tangent complex is $\gAmod[1]$ and whose
obstruction groups are
$H^{n+2}(\gAmod)$ for $n$-th order deformations.
\end{proposition}

\begin{proof}
This is the standard correspondence between dg~Lie
algebras and formal moduli problems
(Hinich, Lurie~\cite{Lurie-DAG-X},
Pridham). The functor satisfies Schlessinger's
conditions: $H^0$ gives automorphisms, $H^1$ gives
deformations, $H^2$ gives obstructions.
The weight filtration on $\gAmod$ provides the
Artinian stratification.
\end{proof}

The support depth \(r^\Theta_{\max}(\cA)\) is the length of the
bar-intrinsic MC packet along the weight direction.  The obstruction
depth \(d_{\mathrm{obs}}(\cA)\) is the length of the \(H^2\)-obstruction
profile.  When \(r^\Theta_{\max}<\infty\), the bar-intrinsic branch is
finite type in the weight direction; when
\(r^\Theta_{\max}=\infty\), the branch is a completed formal object.
Other branches of the formal moduli problem are controlled by the same
obstruction-extension sequence and are not classified by the
bar-intrinsic packet alone.

% ================================================================
% 9.8.16: HOLOGRAPHIC SYNTHESIS
% ================================================================

\subsection{Synthesis: holographic reconstruction from
obstruction data}
% label removed: subsec:holographic-synthesis
\index{holographic reconstruction!synthesis}

\begin{theorem}[Universal Holography, admissible shadow projection with realization datum;
\(\ClaimStatusConditional\); licensing tags \(\beta+\gamma+\delta\)]\label{thm:holographic-recon-shadow}
Let $\cA$ be a chiral algebra on a smooth projective curve $X$
equipped with a conformal vector $\omega$ at non-critical level.
Assume that the following data are fixed.
\begin{enumerate}[label=\textup{(H\arabic*)}]
\item \emph{Boundary realisation.} Either $\cA$ lies in the exact
boundary-linear Costello--Gwilliam/Costello--Li sector,
$\cA=V_k(\mathfrak g)$ is an affine Kac--Moody algebra with
$k\neq -h^\vee$ and the holomorphic Chern--Simons boundary model of
Theorem~\ref{thm:E3-topological-km}, or
$\cA=\mathcal{W}_k(\mathfrak g,f)$ is obtained by Drinfeld--Sokolov
reduction of $V_k(\mathfrak g)$ with $k\neq -h^\vee$ and carries the
Costello--Gaiotto DS boundary model.  This datum includes a chosen
holomorphic-topological BV/factorization model
\(\mathcal F_{\cA,\Xi}\) and a boundary comparison
\[
\eta^\partial_\cA\colon
\Obs^\partial(\mathcal F_{\cA,\Xi})\xrightarrow{\simeq}\cA .
\]
\item \emph{Bulk--centre comparison.} The chosen model lies in the
scope of the bulk--Hochschild comparison, and a map
\[
\chi_{\mathrm{HT},\cA}\colon
\Zderch(\cA)\to
\Obs^{\mathrm{bulk}}(\mathcal F_{\cA,\Xi})
\]
is part of the datum in the declared ambient.  In the
physical prefactorisation sector this is the comparison supplied by
Theorem~\ref{thm:bulk_hochschild}. In the boundary-linear sector the
compact-generation and derived-centre quasi-isomorphism hypotheses are
those of Theorem~\ref{thm:boundary-linear-bulk-boundary}. In the DS
sector the comparison is read through the DS--Hochschild bridge
Theorem~\ref{thm:ds-hochschild-bridge}: cohomologically for
class~$\mathsf M$, at chain level in the weight-completed/pro ambient,
and on the original complex only when the bounded conformal-weight-shift
HPL criterion of Remark~\ref{rem:class-M-hpl-convergence} is supplied.
\item \emph{Topological refinement.} If the conclusion is read as a
strict $E_3$-topological statement, the topologisation datum of
Convention~\ref{conv:chd-e3-chiral-meaning} is part of the input. For
affine Kac--Moody and DS $\mathcal W$-algebras this datum is supplied by
Theorems~\ref{thm:E3-topological-km},
\ref{thm:E3-topological-DS}, and
\ref{thm:E3-topological-DS-general}.
\end{enumerate}
Write
\[
\Xi_\cA=(\mathcal F_{\cA,\Xi},\eta^\partial_\cA,
\chi_{\mathrm{HT},\cA},\mathcal A(\cA)).
\]
Then the \(\Xi\)-realized admissible datum has the supplied relative
$3$d holomorphic-topological gauge theory
\(\mathcal{T}_{\cA,\Xi}:=\mathcal F_{\cA,\Xi}\) on
$X\times\mathbb{R}$ such that:
\begin{enumerate}
\item[(i)] \emph{Boundary restriction.}
\(\eta^\partial_\cA\colon
\mathrm{Obs}^\partial(\mathcal{T}_{\cA,\Xi})\simeq \cA\)
as $E_1$-chiral algebras.
\item[(ii)] \emph{Bulk--centre comparison.}
\[
\chi_{\mathrm{HT},\cA}\colon
Z^{\mathrm{der}}_{\mathrm{ch}}(\cA)\xrightarrow{\simeq}
\mathrm{Obs}^{\mathrm{bulk}}(\mathcal{T}_{\cA,\Xi})
\]
is an equivalence in the algebraic, cohomological, or completed/pro
ambient specified by the datum. Under hypothesis \textup{(H3)} this is an
equivalence of $E_3$-topological factorisation algebras on
$X\times\mathbb{R}$.
\item[(iii)] \emph{Functoriality.} On morphisms preserving the chosen
boundary-linear, affine, or DS realisation and the comparison package,
\((\cA,\Xi_\cA)\mapsto \mathcal{T}_{\cA,\Xi}\)
extends to a functor
\[
 \Phi_{\mathrm{hol}}\colon
 \mathsf{ChirAlg}^{\omega,\mathrm{adm},\Xi}_X
 \longrightarrow
 \mathsf{HT\text{-}QFT}^{\Xi}_{X\times\mathbb{R}}
\]
on the admissible subcategory, compatible with Drinfeld--Sokolov
reduction on the \(\Xi\)-compatible DS locus:
\[
\Phi_{\mathrm{hol}}(\mathcal{W}_k(\mathfrak{g},f),\Xi_{\mathrm{DS}})
\simeq
\Phi_{\mathrm{hol}}(V_k(\mathfrak{g}),\Xi_{V_k})
/\!/_{\mathrm{BV}}\mathrm{DS}_f .
\]
\item[(iv)] \emph{Class coverage.} Class $\mathrm{G}$ and the free-field
part of class $\mathrm{C}$ are realised by Costello--Li abelian
holomorphic Chern--Simons on the boundary-linear exact locus; class
$\mathrm{L}$ is realised by the affine holomorphic Chern--Simons model
of Theorem~\ref{thm:E3-topological-km}; class $\mathrm{M}$ (Virasoro,
$\mathcal{W}_N$) is realised by Costello--Gaiotto holomorphic Chern--
Simons with Drinfeld--Sokolov boundary condition, with bulk comparison
controlled by the DS--Hochschild bridge
$\mathrm{ChirHoch}^\bullet(\mathcal{W}_k(\mathfrak{g}))\simeq
H^\bullet_{\mathrm{DS}}(\mathrm{ChirHoch}^\bullet(V_k(\mathfrak{g})))$
in the cohomological or completed/pro sense specified in \textup{(H2)}.
\end{enumerate}
\end{theorem}

\begin{proof}
Hypothesis \textup{(H1)} supplies the BV theory and the boundary
comparison. In the
boundary-linear exact sector, boundary recovery is the
Costello--Gwilliam restriction theorem applied to the Costello--Li
model. For affine Kac--Moody it is the holomorphic Chern--Simons
boundary model of Theorem~\ref{thm:E3-topological-km}. In the DS sector,
Costello--Gaiotto identify the DS boundary condition for holomorphic
Chern--Simons, and Arakawa's theorem
$H^\bullet_{\mathrm{DS}}(V_k(\mathfrak{g}))\simeq
\mathcal{W}_k(\mathfrak{g})$ identifies the boundary chiral algebra.
The comparison with the boundary observables is \(\eta^\partial_\cA\).

Hypothesis \textup{(H2)} gives the bulk comparison. For
boundary-linear models, Theorem~\ref{thm:boundary-linear-bulk-boundary}
supplies the Morita comparison with the derived centre of the boundary
algebra. For a chosen physical prefactorisation model,
Theorem~\ref{thm:bulk_hochschild} supplies
\(\Psi_{\mathrm{Ran}}\), equivalently the inverse comparison
\(\chi_{\mathrm{HT},\cA}\). In the DS sector,
Theorem~\ref{thm:ds-hochschild-bridge} supplies exactly the stated
cohomological or completed/pro comparison, with original-complex
chain-level comparison only under the bounded-shift HPL criterion.
Thus \(\chi_{\mathrm{HT},\cA}\) identifies the algebraic closed sector
with \(\mathrm{Obs}^{\mathrm{bulk}}(\mathcal{T}_{\cA,\Xi})\) in the
asserted ambient.
Hypothesis \textup{(H3)} upgrades this algebraic open/closed bulk to
the strict $E_3$-topological factorisation algebra, by
Theorem~\ref{thm:chiral-higher-deligne} together with the named
topologisation theorems.

Functoriality is relative to the admissible data: a morphism preserving
the boundary-linear, affine, or DS presentation and the \(\Xi\)-package
is a morphism of boundary conditions in the Costello BV complex, hence
induces the boundary map on observables and the corresponding map on
Hochschild cochains and derived centres, compatible with
\(\chi_{\mathrm{HT}}\). The DS formula is the functoriality of the
Costello--Gaiotto DS boundary condition composed with
Theorem~\ref{thm:ds-hochschild-bridge} on the \(\Xi\)-compatible locus.
The class coverage is precisely the disjoint union of the three
constructions in \textup{(H1)}.
\end{proof}

\begin{remark}[Shadow-tower reconstruction]
The following theorem reconstructs the bar-intrinsic gauge class of
\(\Theta_\cA\) from a compatible shadow reconstruction packet.
Obstruction classes alone decide extendability; they do not select a
point of the \(H^1\)-torsor of lifts in the exact sequence
\[
H^1(J^{r+1}(\cA),d_2)\to E_\cA(r+1)\to E_\cA(r)\to
H^2(J^{r+1}(\cA),d_2).
\]
Thus the packet contains both the obstruction class and the lift
coordinate in a chosen bar-intrinsic gauge slice. Finite-order data
suffice exactly on the finite-depth strata; class~$\mathsf M$ is a
completed inverse-limit statement.
\end{remark}

\begin{theorem}[Filtered shadow-tower reconstruction of the bar-intrinsic MC gauge class;
\(\ClaimStatusConditional\); licensing tags \(\alpha+\gamma+\epsilon\)]
% label removed: thm:holographic-reconstruction
\label{thm:filtered-shadow-tower-reconstruction}
\index{holographic reconstruction!main theorem}
Let $\cA$ be a modular Koszul chiral algebra with
holographic modular Koszul datum
$\mathcal{H}(T) = (\cA, \cA^!, \mathcal{C}, r(z),
\Theta_\cA, \nabla^{\mathrm{hol}})$
\textup{(}Definition~\textup{\ref{V1-def:holographic-modular-koszul-datum})}.
Assume the modular convolution dg Lie algebra \(\gAmod\) is complete
and separated for the weight filtration \(F^\bullet\), and fix the
bar-intrinsic gauge component containing
\(\Theta_\cA\in F^2{\gAmod}^1\).  A \emph{shadow reconstruction packet}
is the compatible system
\[
\mathfrak S(\cA)=
\Bigl(
S_2,\;S_3,\;S_4,\;
\{(o_n(\cA),\lambda_n)\}_{n\geq 5}
\Bigr),
\]
where:
\begin{enumerate}[label=\textup{(\roman*)}]
\item \(S_2=H_\cA=\kappaChHodge(\cA)\), the Hessian
 \textup{(}degree~$2$\textup{)};
\item \(S_3=\mathfrak C_\cA\), the cubic shadow
 \textup{(}degree~$3$\textup{)};
\item \(S_4=\mathfrak Q^{\mathrm{ct}}_\cA\), the quartic contact
 \textup{(}degree~$4$\textup{)};
\item \(o_n(\cA)\in H^2(J^n(\cA),d_2)\) is the degree-\(n\)
 obstruction class, and
 \(\lambda_n\in H^1(J^n(\cA),d_2)\) is the lift coordinate in the
 \(H^1\)-torsor of lifts, measured in the chosen bar-intrinsic gauge
 slice.
\end{enumerate}
The compatible packet \(\mathfrak S(\cA)\) determines the gauge class
\([\Theta_\cA]\in \varprojlim_r E_\cA(r)\), and hence the
bar-intrinsic MC element \(\Theta_\cA\) in the chosen gauge component.
The reconstruction has the following scope.
\begin{itemize}
\item \emph{Finite support depth.} If \(r^\Theta_{\max}<\infty\), then the finite
 packet
 \[
 (S_2,S_3,S_4,\{(o_n,\lambda_n)\}_{5\leq n\leq r^\Theta_{\max}})
 \]
 determines \([\Theta_\cA]\), after the gauge-slice convention
 \(\lambda_n=0\) for \(n>r^\Theta_{\max}\).
\item \emph{Completed} if \(r^\Theta_{\max} = \infty\): the compatible system
 \(\{(o_n,\lambda_n)\}_{n\geq 5}\) determines
 \[
 [\Theta_\cA]=\varprojlim_r[\Theta_\cA^{\leq r}]
 \]
 in the \(F^\bullet\)-completed convolution algebra topology.
 Analytic convergence of scalar partition functions is a separate
 sewing/norm hypothesis.
\end{itemize}
The four shadow classes provide the classification:
\[
\begin{array}{lcl}
\textbf{G}\;\text{(Gaussian)}: & &
 \text{canonical gauge slice; } S_2 \text{ suffices.}\\
\textbf{L}\;\text{(Lie/tree)}: & &
 \text{canonical gauge slice; } (S_2,S_3) \text{ suffice.}\\
\textbf{C}\;\text{(Contact)}: & &
 \text{canonical gauge slice; } (S_2,S_3,S_4) \text{ suffice.}\\
\textbf{M}\;\text{(Mixed)}: & &
 \text{the full compatible sequence }
 \{(o_n,\lambda_n)\}_{n\ge5} \text{ is required.}
\end{array}
\]
\end{theorem}

\begin{proof}
The obstruction-extension exact sequence identifies the fiber of
\(E_\cA(r+1)\to E_\cA(r)\) over a liftable point with the
\(H^1(J^{r+1}(\cA),d_2)\)-torsor of lifts, and sends a point of
\(E_\cA(r)\) to the obstruction class in
\(H^2(J^{r+1}(\cA),d_2)\).  Inductively, \(S_2,S_3,S_4\) initialize
the truncation and the pair \((o_n,\lambda_n)\) decides whether the
next lift exists and which lift in the torsor is chosen.  Hence a
compatible packet determines a unique element of every \(E_\cA(r)\)
in the fixed bar-intrinsic gauge component.

The finite-depth assertion is
Theorem~\ref{V1-thm:shadow-depth-dichotomy} together with
Theorem~\ref{V1-thm:gaussian-rmax-two} (class~\textbf{G}),
Theorem~\ref{V1-thm:lie-rmax-three} (class~\textbf{L}),
Theorem~\ref{V1-thm:contact-rmax-four} (class~\textbf{C}),
and
Theorem~\ref{V1-thm:virasoro-rmax-infinity} (class~\textbf{M}).
For \(r^\Theta_{\max}<\infty\) the canonical gauge slice has
\(\lambda_n=0\) above \(r^\Theta_{\max}\), so no higher torsor coordinate
enters.  For class~\(\mathsf M\), no finite truncation determines the
tail, and the completed inverse-limit statement is exactly
Corollary~\ref{V1-cor:mittag-leffler-shadow-tower}.  It is a formal
statement in the \(F^\bullet\)-completed convolution algebra, not an
analytic partition-function convergence theorem.
\end{proof}

\begin{corollary}[$\Xi$-realized Universal Holography, strict $E_3$ projection;
\(\ClaimStatusConditional\); licensing tags \(\beta+\gamma+\delta\)]\label{thm:holographic-recon-e3}
Let \((\mathcal A,\Xi_{\mathcal A})\) satisfy the admissibility
hypotheses of Theorem~\ref{thm:holographic-recon-shadow}. If the
topologisation datum of hypothesis \textup{(H3)} is fixed, then
\[
Z^{\mathrm{der}}_{\mathrm{ch}}(\mathcal A)
\xrightarrow{\;\chi_{\mathrm{HT},\mathcal A}\;\simeq\;}
\mathrm{Obs}^{\mathrm{bulk}}(\mathcal T_{\mathcal A,\Xi})
\]
is the comparison equivalence of $E_3$-topological factorisation
algebras on
$X\times\mathbb R$. For affine Kac--Moody and principal or
non-principal DS $\mathcal W$-algebras, the required topologisation
datum is supplied by Theorems~\ref{thm:E3-topological-km},
\ref{thm:E3-topological-DS}, and
\ref{thm:E3-topological-DS-general}.
\end{corollary}
\begin{proof}
This is Theorem~\ref{thm:holographic-recon-shadow}\textup{(ii)}
after imposing hypothesis \textup{(H3)} on the \(\Xi\)-realized
comparison map \(\chi_{\mathrm{HT},\mathcal A}\). The named affine and
DS topologisation theorems verify \textup{(H3)} in the listed cases.
\end{proof}

% ================================================================
% 9.8.17: WORKED RECONSTRUCTIONS
% ================================================================

\subsection{Worked reconstructions}
% label removed: subsec:worked-reconstructions
\index{holographic reconstruction!worked examples}

\begin{example}[Heisenberg reconstruction]
% label removed: ex:heisenberg-reconstruction
\index{Heisenberg algebra!holographic reconstruction}
Input: $\kappaChHodge(\mathcal{H}_k) = k$.
Output: $\Theta_{\mathcal{H}_k} = k\,x^2$.
Reconstruction terminates at step~$1$.
\end{example}

\begin{example}[Affine Kac--Moody reconstruction]
% label removed: ex:km-reconstruction
\index{affine Kac--Moody!holographic reconstruction}
Input: $\kappaChHodge(\widehat{\fg}_k) = \dim(\fg)\cdot(k+h^\vee)/(2h^\vee)$,
$\mathfrak{C}_{\widehat{\fg}} = f^{ab}{}_c\,x_a x_b x^c$.
Output:
$\Theta_{\widehat{\fg}_k} = \bigl(\dim(\fg)\cdot(k+h^\vee)/(2h^\vee)\bigr)\,\kappa_{ab}x^a x^b
+ f^{ab}{}_c x_a x_b x^c$.
Reconstruction terminates at step~$2$.
\end{example}

\begin{example}[$\beta\gamma$ reconstruction]
% label removed: ex:betagamma-reconstruction
\index{beta-gamma system@$\beta\gamma$ system!holographic reconstruction}
Input: $\kappaChHodge(\beta\gamma) = 6\lambda^2 - 6\lambda + 1$,
$\mathfrak{Q}^{\mathrm{ct}}_{\beta\gamma} \neq 0$.
Output:
$\Theta_{\beta\gamma}^{\leq 4} = \kappa\,x^2 +
\mathfrak{Q}^{\mathrm{ct}}\,x^4$.
Reconstruction terminates at step~$3$ (the cubic is
gauge-trivial).
\end{example}

\begin{example}[Virasoro reconstruction]
% label removed: ex:virasoro-reconstruction
\index{Virasoro algebra!holographic reconstruction}
Input: $\kappa = c/2$, $\mathfrak{C} = 2x^3$,
$Q^{\mathrm{ct}} = 10/(c(5c+22))$, and the full shadow recursion. The
cubic-source contribution satisfies
$S_{r+1}^{\mathrm{cub}} = -6r/(c(r+1))\,S_r^{\mathrm{cub}}$, while the
exact coefficients for $r\geq 6$ also contain higher bracket terms.
Output:
$\Theta_{\mathrm{Vir}_c} = \sum_{r=2}^\infty S_r\,x^r$ in the completed
generic ambient, with $S_2,S_3,S_4,S_5$ as in
Proposition~\ref{V1-prop:virasoro-shadow-coefficients} and the tail
read through the full recursion.
At $c = 26$:
\[
\Theta_{\mathrm{Vir}_{26}} = 13x^2 + 2x^3
+ \frac{10}{26(152)}\,x^4
- \frac{48}{26^2(152)}\,x^5 + \cdots,
\]
where the displayed terms are exact through degree~$5$; analytic
geometric convergence of the tail is not asserted without the sewing
norm.
\end{example}

% ================================================================
% 9.8.18: SHAPOVALOV SINGULARITIES
% ================================================================

\subsection{Shapovalov singularities and reconstruction walls}
% label removed: subsec:shapovalov-singularities
\index{Shapovalov determinant!reconstruction walls}

\begin{proposition}[Shapovalov singularities of the
shadow obstruction tower]
% label removed: prop:shapovalov-shadow-singularities
\index{shadow tower!singularities}
For the Virasoro algebra, the shadow coefficient $S_r$
has poles at the zeros of the Shapovalov determinant
$\det S_h$ for $h \leq 2r$. Specifically:
\begin{enumerate}[label=\textup{(\roman*)}]
\item $S_4 = 10/(c(5c+22))$ has poles at $c = 0$ and
 $c = -22/5$, from the weight-$4$ Shapovalov zeros.
\item For $r \geq 5$, the cubic-source approximation
 $S_r^{\mathrm{cub}}$ has poles at $c = 0$
 (of order $r - 3$) and $c = -22/5$ (of order~$1$).
 No new poles appear from the cubic-source recursion.
 The full recursion, which includes
 $\{\mathrm{Sh}_j, \mathrm{Sh}_k\}$ cross-terms
 for $r \geq 6$, introduces higher powers of
 $(5c+22)$ in the denominator; see
 Theorem~\textup{\ref{thm:thqg-virasoro-tower-explicit}}.
\end{enumerate}
\end{proposition}

\begin{proof}
The cubic-source recursion
$S_{r+1}^{\mathrm{cub}} = -6r/(c(r+1))\,S_r^{\mathrm{cub}}$
introduces a pole at $c = 0$ at each step. The pole at
$c = -22/5$ comes from $S_4$ and is preserved but not
amplified by the cubic-source recursion.
\end{proof}

\begin{remark}[Singular vectors as reconstruction walls]
% label removed: rem:singular-vectors-walls
\index{singular vectors!as reconstruction walls}
At the Shapovalov zeros, the shadow obstruction tower develops
singularities: the quasi-primary composites become
linearly dependent, and the jet space dimension drops.
At $c = -22/5$, the weight-$4$ Verma module has a null
vector, so $\Lambda$ becomes a descendant and the quartic
contact shadow degenerates. At admissible levels
$c = c_{p,q}$ (the minimal model central charges), the
Virasoro algebra has finitely many primary fields, and
the shadow obstruction tower reduces to a finite-dimensional problem
controlled by the fusion rules.
\end{remark}

\begin{proposition}[Pole structure of the cubic-source shadow
approximation]
% label removed: prop:pole-structure-shadow-series
\index{shadow tower!pole structure}
The cubic-source approximation to the Virasoro shadow series
$\Theta_{\mathrm{Vir}_c}^{\mathrm{cub}}(x)
= \sum_{r=2}^\infty S_r^{\mathrm{cub}} x^r$
has the following pole structure as a function of~$c$:
\begin{enumerate}[label=\textup{(\roman*)}]
\item At $c = 0$: the terms $S_r^{\mathrm{cub}}$ with $r \geq 4$ have
 poles of increasing order, so the formal series has
 an essential singularity.
\item At $c = -22/5$: each $S_r^{\mathrm{cub}}$ with $r \geq 4$ has a
 simple pole, so the cubic-source series has a simple pole
 at $c = -22/5$.
\item At all other values of $c$: the cubic-source formal series
 is analytic (as a formal power series in~$x$ with
 $c$-dependent coefficients).
\end{enumerate}
The exact tower from degree~$6$ onward also contains the higher bracket
terms of Theorem~\textup{\ref{thm:thqg-virasoro-tower-explicit}}, and
its pole structure must be read from that full recursion.
\end{proposition}

\begin{proof}
For the cubic-source contribution, the closed form is
$S_r^{\mathrm{cub}} = (-1)^{r-4} 6^{r-5} \cdot 240 /
(c^{r-3}(5c+22) \cdot r)$
for $r \geq 5$: the pole at $c = 0$ has order $r - 3$
(growing with $r$), giving essential singularity in the
summed series; the pole at $c = -22/5$ has order~$1$
for all~$r$, giving a simple pole in the cubic-source summed series.
\end{proof}

\begin{remark}[Weight-$6$ Shapovalov analysis]
% label removed: rem:weight-6-shapovalov
\index{Shapovalov determinant!weight 6}
At conformal weight~$6$, the Virasoro vacuum Verma module
has a $4$-dimensional space with basis
$L_{-6}|0\rangle$, $L_{-4}L_{-2}|0\rangle$,
$L_{-3}^2|0\rangle$, $L_{-2}^3|0\rangle$.
The Shapovalov determinant is
\[
\det S_6 = C_6 \cdot c^4 \cdot (5c+22)^2 \cdot (2c+1)
\cdot (7c+68),
\]
where $C_6$ is a positive constant. The zero at
$c = -1/2$ is the weight-$6$ null vector of the Ising
model ($c = 1/2$ in the unitary parametrization, but
$c = -1/2$ appears in the dual channel); the zero at
$c = -68/7$ controls the $\Walg_4$ structure constant
(Theorem~\ref{V1-thm:c334}). These zeros affect the
multi-generator shadow obstruction tower at degree~$6$ and above.
They do not affect the cubic-source primary-line approximation, which
uses only the propagator $P = 2/c$ and the quadratic Hessian; the exact
primary-line coefficients are read through the full recursion.
\end{remark}

\begin{computation}[Shapovalov zeros and shadow
singularities]
% label removed: comp:shapovalov-zeros-shadow
\index{Shapovalov determinant!zeros versus shadow poles}
The following table lists the first several Shapovalov
zeros and their effect on the shadow obstruction tower:

\begin{center}
\renewcommand{\arraystretch}{1.4}
\begin{tabular}{clll}
\toprule
Weight $h$ & Zero location & Source & Effect on shadow \\
\midrule
$4$ & $c = 0$ (order $2$) & $\det G = c^2(5c+22)/2$
 & Pole in $S_r^{\mathrm{cub}}$ for $r \geq 4$ \\
$4$ & $c = -22/5$ & $\det G = c^2(5c+22)/2$
 & Simple pole in $S_r^{\mathrm{cub}}$ for $r \geq 4$ \\
$6$ & $c = -1/2$ & $\det S_6 \ni (2c+1)$
 & No effect on cubic-source primary line \\
$6$ & $c = -68/7$ & $\det S_6 \ni (7c+68)$
 & No effect on cubic-source primary line \\
$8$ & $c = -46/3$ & $\det S_8 \ni (3c+46)$
 & No effect on cubic-source primary line \\
\bottomrule
\end{tabular}
\end{center}

\noindent The cubic-source primary-line coefficients see only the
weight-$4$ Shapovalov zeros. The higher zeros affect the
\emph{extended} shadow obstruction tower (including descendants and
quasi-primary composites at higher weights), which controls the
multi-channel reconstruction for $\Walg_N$ with $N \geq 3$ and enters
the exact Virasoro tower through the full recursion from degree~$6$
onward.
The curvature of the multi-channel shadow connection equals
the propagator variance~$\delta_{\mathrm{mix}}$
(Volume~I, Theorem~\ref{V1-thm:propagator-variance}).
\end{computation}

% ================================================================
% 9.8.19: SPECIFIC CENTRAL CHARGES
% ================================================================

\subsection{Specific central charges}
% label removed: subsec:specific-central-charges
\index{Virasoro algebra!specific central charges}

\begin{example}[Ising model, $c = 1/2$]
% label removed: ex:ising-shadow
\index{Ising model!shadow tower}
$\kappa = 1/4$, $Q^{\mathrm{ct}} = 40/49$,
$S_5 = -384/49$.
The ratio $6/c = 12 > 1$: the shadow obstruction tower diverges
formally, reflecting the strong-coupling regime.
\end{example}

\begin{example}[Critical string, $c = 26$]
% label removed: ex:critical-string-shadow
\index{bosonic string!shadow tower!$c = 26$}
$\kappa = 13$, $Q^{\mathrm{ct}} = 5/1976$,
$|S_{r+1}^{\mathrm{cub}}/S_r^{\mathrm{cub}}|
\to 3/13 \approx 0.231$ for the cubic-source approximation.
The completed shadow tower is controlled formally; analytic convergence
of the exact tail requires the sewing norm. The Koszul dual
$\mathrm{Vir}_{26}^! = \mathrm{Vir}_0$ has $\kappa' = 0$:
a degenerate dual, the depth-zero resonance shadow
(Remark~\ref{V1-rem:virasoro-resonance-model}).
\end{example}

\begin{example}[Virasoro self-duality, $c = 13$]
% label removed: ex:virasoro-self-dual
\index{Virasoro algebra!self-duality at $c = 13$}
$\mathrm{Vir}_{13}^! = \mathrm{Vir}_{13}$
(self-dual). $\kappa = 13/2$,
$|S_{r+1}^{\mathrm{cub}}/S_r^{\mathrm{cub}}|
\to 6/13 \approx 0.462$ for the cubic-source approximation.
Self-duality means the complementarity splitting
$\mathbf{C}_g = \mathbf{Q}_g \oplus \mathbf{Q}_g^!$
satisfies $\mathbf{Q}_g \cong \mathbf{Q}_g^!$:
the two Lagrangians coincide.
\end{example}

% ================================================================
% 9.8.20: COMPLEXITY HIERARCHY
% ================================================================

\subsection{The reconstruction complexity hierarchy}
% label removed: subsec:complexity-hierarchy
\index{reconstruction complexity|textbf}

\begin{theorem}[Complexity hierarchy]
% label removed: thm:complexity-hierarchy
\index{complexity hierarchy!holographic}
The number of independent parameters $N(r)$ needed to
specify $\Theta_\cA^{\leq r}$ on the primary line is:
\begin{equation}% label removed: eq:complexity-hierarchy
N(r) = \begin{cases}
1 & \text{for class } \textbf{G}\;\;(\text{all } r), \\
2 & \text{for class } \textbf{L}\;\;(\text{all } r), \\
3 & \text{for class } \textbf{C}\;\;(\text{all } r), \\
r - 1 & \text{for class } \textbf{M}\;\;(r \geq 2).
\end{cases}
\end{equation}
\end{theorem}

\begin{proof}
For classes \textbf{G}, \textbf{L}, \textbf{C}: the bar-intrinsic
packet has support through \(r^\Theta_{\max}\), giving
\(r^\Theta_{\max}-1\) primary-line parameters in the canonical gauge
slice.
For class~\textbf{M}: each degree level adds one
primary-line parameter $S_r$, giving $r - 1$ through
degree~$r$.
\end{proof}

% ================================================================
% 9.8.21: CONCLUDING PERSPECTIVE
% ================================================================

\subsection{Concluding perspective}
% label removed: subsec:reconstruction-conclusion

The holographic reconstruction problem (determining the
bar-intrinsic MC gauge class from its shadow-tower data, with
finite-order data only on finite support-depth strata) is controlled by
the extension tower \(\{E_\cA(r)\}\), its obstruction classes
\(\{o_{r+1}(\cA)\}\), and its \(H^1\)-lift coordinates. The dichotomy
theorem gives the clean criterion: reconstruction terminates if and
only if the support shadow depth is finite. The four
shadow classes (Gaussian, Lie/tree, contact/quartic,
mixed/infinite) exhaust the structural possibilities
for the standard landscape.

The synthesis of deformation theory (the extension tower as a
Postnikov filtration), representation theory (the Shapovalov
determinant controlling jet space dimensions), combinatorial algebra
(graph-sum contributions to each obstruction), and information theory
(the polynomial growth rate of jet space dimensions) gives the formal
structure of the reconstruction problem on the stated ambient.

For the Virasoro algebra, the simplest class~\textbf{M}
example, the cubic-source recursion
$S_{r+1}^{\mathrm{cub}} = -6r/(c(r+1))S_r^{\mathrm{cub}}$ controls the
displayed approximation; the exact tower includes higher bracket terms
from degree~$6$ onward. The completed inverse-limit reconstruction is
formal. Analytic convergence in an operator or sewing norm is supplied
only under the hypotheses of the HS-sewing theorem
(Theorem~\ref{thm:general-hs-sewing}).

The gravitational-complexity interpretation
(Section~\ref{sec:thqg-gravitational-complexity}) gives
these results a conditional physical reading: the shadow depth of the
boundary algebra measures the homotopy-algebraic complexity of the
bulk model in the chosen HT construction. Gaussian boundaries give free
bulk theories; Lie/tree boundaries give semistrict higher-spin models;
contact boundaries give quartic vertices; mixed boundaries require an
infinite completed $L_\infty$ tower. Identifying that tower with a
metric-gravity saddle expansion is a separate physical hypothesis.

% ======================================================================
%
% TWISTED HOLOGRAPHY AND BRANE EXAMPLES
%
% ======================================================================

\subsection{Twisted holography examples: D3, M2, and ABJM}
\label{subsec:holographic-twisted-examples}
\index{twisted holography!examples|textbf}

The holographic modular Koszul datum
$\mathcal{H}(\mathcal{T})
= (\cA, \cA^!, \mathcal{C}, r(z), \Theta_\cA,
\nabla^{\mathrm{hol}})$
is computed for the standard brane systems of the
Costello--Li programme \cite{CL16,CostelloGaiotto18}.

\begin{proposition}[Twisted $\mathcal{N} = 4$ SYM datum]
\label{prop:twisted-n4-datum}
\index{twisted holography!D3 brane}
For the D3 brane \textup{(}twisted $\mathcal{N} = 4$ SYM
at rank~$N$\textup{)}:
\begin{enumerate}[label=\textup{(\roman*)}]
\item The boundary VOA is affine $\mathfrak{gl}_N$ at level $k = 1$
 with $\kappaChHodge(\mathfrak{gl}_N, 1)
 = (N^2{-}1)(N{+}1)/(2N) + 1$.
\item The holographic R-matrix is $r(z) = k\,\Omega_{\mathrm{tr}}/z$
 with $k = 1$.
\item The genus expansion
 \(F_g(\mathfrak{gl}_{N,1})
 = \kappaChHodge(\mathfrak{gl}_{N,1})\lambda_g^{\mathrm{FP}}\)
 \textup{(}uniform-weight\textup{)} computes twisted amplitudes.
\item The anomaly matching
 $\kappa_{\mathrm{eff}} = \kappaChHodge(\mathrm{matter})
 + \kappaChHodge(\mathrm{ghost}) = 0$
 is a consistency check.
\end{enumerate}
\end{proposition}

\begin{proposition}[ABJM holographic datum]
\label{prop:abjm-datum}
\index{ABJM theory!holographic datum}
For ABJM at rank~$N$ and level~$k$:
\begin{enumerate}[label=\textup{(\roman*)}]
\item The boundary VOA is the BRST reduction
 $\cA_{\mathrm{ABJM}}(N,k)
 = H^0_{\mathrm{BRST}}(V_k(\mathfrak{gl}_N) \otimes
 V_{-k}(\mathfrak{gl}_N) \otimes \mathrm{Sb}^{4N^2})$.
\item The modular characteristic is
 $\kappaChHodge(\cA_{\mathrm{ABJM}}) = -(N^2+1)$
 \textup{(}from $(N^2{-}1)$ for the $\mathrm{CS}_k + \mathrm{CS}_{-k}$
 sectors and $-2N^2$ for $4N^2$ symplectic bosons;
 $\kappa \sim -N^2$ at large~$N$\textup{)}.
\item In the large-$N$ matrix-model limit, assuming the ABJM
 Fermi-gas resummation, the shadow connection degenerates
 to the Airy connection and the ABJM partition function is
 $Z(N,k) = C_k^{-1/3}\mathrm{Ai}(C_k^{-1/3}(N - B_k))$
 with $C_k = 2/(\pi^2 k)$.
 The $N^{3/2}$ scaling
 $F \sim (\pi\sqrt{2k}/3)N^{3/2}$ is the hallmark
 of M2-brane degrees of freedom.
\end{enumerate}
\end{proposition}

\begin{remark}[Seiberg--Witten theory and the shadow connection]
\label{rem:holographic-sw}
\index{Seiberg--Witten!shadow connection}
The shadow connection
$\nabla^{\mathrm{sh}} = d - Q'_L/(2Q_L)\,dt$
for the Virasoro algebra is a Fuchsian connection
with three regular singular points, and its Picard--Fuchs
equation is compared with the $\mathrm{SU}(2)$ Seiberg--Witten
Gauss--Manin connection. The identification maps:
shadow metric zeros $\leftrightarrow$ singular fibers
$u = \pm\Lambda^2$;
Koszul monodromy $-1$ $\leftrightarrow$ SW monodromy;
$\kappa = c/2$ $\leftrightarrow$ prepotential coefficient.
Under this comparison, the instanton coefficients $F_k$ of the SW
prepotential are Taylor coefficients of the flat sections of the shadow
connection.
\end{remark}


% ======================================================================
%
% K3 x E: THE COSTELLO-LI PROTOTYPE
%
% ======================================================================

\subsection{\texorpdfstring{$K3 \times E$}{K3 x E}: the Costello--Li prototype}
\label{subsec:k3e-holographic-datum}
\index{K3 surface!twisted holography|textbf}
\index{Costello--Li programme!K3 x E@$K3 \times E$|textbf}
\index{holographic modular Koszul datum!K3 x E@$K3 \times E$|textbf}
\index{Kodaira--Spencer theory!K3 x E@$K3 \times E$}

The product $X = K3 \times E$ of a K3 surface with an elliptic curve
is a compact Calabi--Yau threefold. On the abelian harmonic branch of
the Costello--Li reduction, the boundary algebra is the rank-$24$
Heisenberg algebra, and the corresponding shadow tower has class
$\mathsf{G}$. This branch is a computable geometric test object for
the holographic modular Koszul formalism.

\subsubsection{The HT twist and Kodaira--Spencer reduction}
\label{subsubsec:k3e-ht-twist}
\index{HT twist!K3 x E@$K3 \times E$}

The holomorphic-topological twist of type~IIB on $K3 \times E$
decomposes the CY$_3$ into:
\begin{itemize}
\item \emph{Holomorphic direction:} the elliptic curve $E$
 ($\dim_\C = 1$).
\item \emph{Topological directions:} the K3 surface
 ($\dim_\C = 2$, $\dim_\R = 4$).
\end{itemize}
The B-model twist produces Kodaira--Spencer gravity on $X$, whose
field is the Beltrami differential
$\mu \in \Omega^{0,1}(T^{1,0}_X)$ with action
$S = \int_X \mu \,\dbar\mu + \tfrac{1}{3!}\int_X \mu^3$.
The observables are polyvector fields: by HKR,
\begin{equation}\label{eq:k3e-hh-hkr}
\HH^n(X) = \bigoplus_{p+q=n} \PV^{p,q}(X),
\qquad
\PV^{p,q}(X) = H^q\!\bigl(X, \textstyle\bigwedge^p T^{1,0}_X\bigr)
\overset{\text{CY}}{=} h^{3-p,q}(X).
\end{equation}

The Hodge diamond of $K3 \times E$ is computed by K\"unneth from
\begin{equation}\label{eq:k3e-hodge-kunneth}
h^{p,q}(K3 \times E) = \sum_{a+c=p,\; b+d=q}
h^{a,b}(K3) \cdot h^{c,d}(E).
\end{equation}
The nonzero entries are:
\begin{center}
\renewcommand{\arraystretch}{1.2}
\small
\begin{tabular}{c|cccccccccc}
$(p,q)$ & $(0,0)$ & $(1,0)$ & $(0,1)$ & $(2,0)$ & $(0,2)$ & $(1,1)$ & $(2,1)$ & $(1,2)$
 & $(3,0)$ & $(0,3)$ \\
\hline
$h^{p,q}$ & $1$ & $1$ & $1$ & $1$ & $1$ & $21$ & $21$ & $21$ & $1$ & $1$
\end{tabular}
\end{center}
with the remaining nonzero entries fixed by $h^{p,q} = h^{q,p}
= h^{3-p,3-q}$ (CY$_3$ Hodge symmetry). In particular:
\begin{equation}\label{eq:k3e-hodge-key}
h^{1,1} = h^{2,1} = 21, \qquad
\chi(K3 \times E) = \chi(K3)\cdot\chi(E) = 24 \cdot 0 = 0.
\end{equation}
The total polyvector field dimension is
$\dim \PV^*(X) = \dim \PV^*(K3) \cdot \dim \PV^*(E)
= 24 \cdot 4 = 96$.

\begin{remark}[Self-mirror with $\kappa_{\mathrm{KS}}^{(1)} = 0$]
\label{rem:k3e-self-mirror}
Since $h^{1,1} = h^{2,1} = 21$, the threefold $K3 \times E$
is self-mirror: the KS modular characteristic
$\kappa_{\mathrm{KS}}^{(1)} = h^{1,1} - h^{2,1} = 0$.
The KS weight-one Hodge row vanishes, and the BCOV constant-map
scalar also vanishes because $\chi(K3\times E)=0$. This is the entry
$K3 \times T^2$ in the KS landscape table of
\S\ref*{subsec:thqg-ks-landscape}.

The boundary chiral algebra $\cA_E$ discussed below is a
\emph{different} object from the KS boundary algebra: it arises
from reducing the B-model \emph{along} K3 to obtain a chiral
algebra \emph{on} $E$, rather than from the abstract KS
deformation theory. The two algebras have different $\kappa$
values: $\kappa_{\mathrm{KS}}^{(1)} = 0$ vs.\
$\kappaChHodge(\cA_E) = 24$. The discrepancy records that
$\kappa$ depends on the full algebra, not on a single
invariant of the geometry.
\end{remark}


\subsubsection{The boundary chiral algebra}
\label{subsubsec:k3e-boundary}
\index{boundary chiral algebra!K3 x E@$K3 \times E$}

On the abelian harmonic branch of the Costello--Li reduction along the
topological K3 directions, one obtains a chiral algebra $\cA_E$ on $E$
with generators from $\PV^*(K3)$. Each polyvector class
$\alpha \in \PV^{p,q}(K3) = H^{2-p,q}(K3)$ contributes a generator of
conformal weight $2 - p$:

\begin{center}
\renewcommand{\arraystretch}{1.2}
\small
\begin{tabular}{lccc}
\textbf{Sector} & \textbf{Hodge type} & \textbf{Count}
 & \textbf{Conformal weight} \\
\hline
$\PV^{0,0}(K3) = H^{2,0}(K3)$ & holomorphic 2-form & $1$ & $2$ \\
$\PV^{0,2}(K3) = H^{2,2}(K3)$ & volume form & $1$ & $2$ \\
$\PV^{1,1}(K3) = H^{1,1}(K3)$ & $(1{,}1)$-forms & $20$ & $1$ \\
$\PV^{2,0}(K3) = H^{0,0}(K3)$ & constant & $1$ & $0$ \\
$\PV^{2,2}(K3) = H^{0,2}(K3)$ & anti-hol.\ 2-form & $1$ & $0$ \\
\hline
\multicolumn{2}{c}{\textbf{Total}} & $\mathbf{24}$
 & $= \chi(K3)$
\end{tabular}
\end{center}

On this branch the central charge of the boundary algebra is
$c(\cA_E) = \chi(K3) = 24$.

\begin{proposition}[$K3 \times E$ boundary datum, abelian harmonic branch]
\label{prop:k3e-boundary-datum}
\index{K3 surface!boundary chiral algebra}
Assume the Costello--Li reduction of Kodaira--Spencer gravity on
$K3 \times E$ is taken on the abelian harmonic branch: the harmonic
projector is multiplicative on the K3 polyvector cohomology used in the
reduction, and exact Schouten brackets are killed after projection. The
boundary chiral algebra $\cA_E$ is then modeled by the free-field
algebra of $24$ bosons indexed by $H^*(K3, \Z)$:
\begin{equation}\label{eq:k3e-boundary-ope}
J^a(z)\, J^b(w) \;\sim\; \frac{\delta^{ab}}{(z-w)^2},
\qquad a, b = 1, \ldots, 24,
\end{equation}
% The Heisenberg OPE uses the unit level matrix delta^{ab}
% (positive definite, one boson per cohomology class at level 1).
% In the orthonormal basis, kappa = 24 * kappa(H_1) = 24 by additivity.
% The Mukai pairing of signature (4,20) enters the lattice VOA
% through the vertex operators e^alpha (topological sector), not
% through the Heisenberg OPE levels. In a Mukai basis the OPE
% would read G^{ab}/(z-w)^2, but kappa = rank = 24 regardless of
% basis (kappa is basis-independent).
where the $24$ generators are indexed by a basis of $H^*(K3, \Z)$, each
at unit level. The Mukai pairing of signature $(4, 20)$ determines the
lattice vertex operators and the topological sector, but the Heisenberg
OPE levels are positive definite.
The free-field assertion is the branch condition: after harmonic
projection, the Schouten bracket on the retained K3 sector is zero and
only the Heisenberg pairing remains in the OPE.

\begin{enumerate}[label=\textup{(\roman*)}]
\item $\kappaChHodge(\cA_E) = 24$ on this rank-$24$ Heisenberg branch
 \textup{(}$= \operatorname{rank}$ for unit-level free bosons\textup{)}.
\item Support shadow depth: class $\mathsf{G}$ \textup{(}Gaussian,
 \(r^\Theta_{\max} = 2\)\textup{)}. All higher shadows vanish: $C = Q = 0$.
\item The critical discriminant
\(\Delta = 8\kappaChHodge(\cA_E) S_4 = 0\)
 \textup{(}$S_4 = 0$ for free fields\textup{)}.
\item The shadow metric
\(Q_L(t) = (2\kappaChHodge(\cA_E))^2 = 2304\) is constant.
\end{enumerate}
\end{proposition}

\begin{proof}
The hypotheses impose a multiplicative harmonic projection and kill the
projected Schouten bracket on the retained K3 cohomology sector. The
projected interaction Lie algebra is therefore abelian. The reduction
along K3 produces a cohomological free-field chiral algebra on $E$ with
one Heisenberg boson per cohomology class of K3, each at unit level.
The Mukai lattice of signature $(4, 20)$ underlies the
topological sector: $(3, 19)$ from $H^2(K3)$ plus $(1, 1)$ from the
$H^0 \oplus H^4$ pairing.

For the modular characteristic: $\cA_E$ consists of 24 free
bosons at level~$1$, so
\(\kappaChHodge(\cA_E)=24 \cdot \kappaChHodge(\cH_1) = 24\)
by additivity (Volume~I, Corollary~\ref{V1-cor:kappa-additivity}).
The support shadow depth is \(r^\Theta_{\max}=2\) because the OPE is
purely quadratic and the canonical Gaussian gauge slice has no higher
lift coordinates.
\end{proof}


\subsubsection{The holographic modular Koszul datum
\texorpdfstring{$\cH(K3 \times E)$}{H(K3 x E)}}
\label{subsubsec:k3e-holographic-datum}

\begin{computation}[$K3 \times E$ holographic datum on the abelian harmonic branch]
\label{comp:k3e-holographic-datum}
\index{holographic modular Koszul datum!K3 x E@$K3 \times E$!components}
Under the hypotheses of Proposition~\ref{prop:k3e-boundary-datum},
the six components of
$\cH(K3 \times E) = (\cA, \cA^!, \mathcal{C}, r(z),
\Theta_\cA, \nabla^{\mathrm{hol}})$ are:
\begin{center}
\renewcommand{\arraystretch}{1.35}
\small
\begin{tabular}{lp{9.5cm}}
\textbf{Component} & \textbf{Value} \\
\hline
(a) $\cA$ & Free-field: $24$ bosons at unit level ($\delta^{ab}$),
 $c = 24$, $\kappa = 24$ \\
(b) $\cA^!$ & Verdier--Koszul partner: the chiral symmetric algebra on
 the dual current space, not the negative-level Heisenberg algebra;
 $\kappaChHodge(\cA^!) = -24$ and
 $\kappa + \kappa' = 0$ \\
(c) $\mathcal{C}$ & $\cA^!\text{-}\mathsf{mod}$;
 the derived center $\Zder^{\mathrm{ch}}(\cA_E)$
 has primitive Heisenberg Hochschild row
 $P_{\cA}^{\mathrm{prim}}(t) = 1 + 24\,t + t^2$ \\
(d) $r(z)$ & $r(z) = k\,\Omega_{\mathrm{tr}}/z$ at unit level $k = 1$,
 $\Omega = \sum \delta^{ab}\,J_a \otimes J_b$ CYBE trivially satisfied \textup{(}abelian\textup{)} \\
(e) $\Theta_\cA$ & $\Theta_\cA^{\leq 2} = \kappa = 24$;
 $\Theta_\cA^{(g,r)} = 0$ for $r \geq 3$
 \textup{(}class $\mathsf{G}$: tower terminates\textup{)} \\
(f) $\nabla^{\mathrm{hol}}$ & Trivial on the shadow metric
 \textup{(}$Q_L$ constant\textup{)};
 nontrivial at the holographic level with
 universal residue $\tfrac{1}{2}$
\end{tabular}
\end{center}

The genus-$g$ free energies are
$F_g(\cA_E) = \kappaChHodge(\cA_E) \cdot \lambda_g^{\mathrm{FP}}$
\textup{(}uniform-weight\textup{)}:
\begin{equation}\label{eq:k3e-genus-amplitudes}
F_1 = 1, \qquad
F_2 = \tfrac{7}{240}, \qquad
F_3 = \tfrac{31}{40320}.
\end{equation}

\emph{Three consistency checks.}
\begin{enumerate}
\item \emph{KS reduction}: The reduction of KS theory along K3 gives
 $24$ free bosons and $\kappa = \operatorname{rank} = 24$.
\item \emph{HKR comparison}: The equality $\dim \PV^*(K3) = 24$
 supplies the current row. The group $\HH^2(K3 \times E)$ has dimension
 $23$ in the geometric category; the chiral boundary center is the
 primitive Heisenberg row above.
\item \emph{Koszul complementarity}: $\kappaChHodge(\cA^!) = -24$;
 $F_1(\cA) + F_1(\cA^!) = 1 + ({-}1) = 0$ \textup{(}Theorem~C
 for free fields\textup{)}.
\end{enumerate}
\end{computation}


\subsubsection{Noncommutative deformations and the Brauer group}
\label{subsubsec:k3e-nc-deformations}
\index{noncommutative deformation!K3 x E@$K3 \times E$|textbf}
\index{Brauer group!K3 x E@$K3 \times E$}
\index{Hochschild cohomology!K3 x E@$K3 \times E$}

The deformation theory of $D^b(K3 \times E)$ is controlled by
$\HH^2(K3 \times E) = 23$, decomposed by HKR as:
\begin{equation}\label{eq:k3e-hh2-decomposition}
\HH^2(K3 \times E)
\;=\;
\underbrace{H^1(T_X)}_{h^{2,1} = 21}
\;\oplus\;
\underbrace{H^0\!\bigl(\textstyle\bigwedge^2 T_X\bigr)}_{h^{1,0} = 1}
\;\oplus\;
\underbrace{H^2(\cO_X)}_{h^{0,2} = 1}.
\end{equation}
The $21$-dimensional summand parametrizes complex-structure
deformations. The two remaining summands are the Poisson and gerby
rows of the HKR decomposition:
\begin{enumerate}[label=\textup{(\roman*)}]
\item The bivector direction $H^0(\bigwedge^2 T_X) \cong \C$
 is the first-order Poisson deformation; for $K3\times E$ it is
 generated by the inverse holomorphic symplectic form on K3.
\item The gerby $B$-field direction $H^2(\cO_X) \cong \C$ is the
 infinitesimal Brauer-twisting row. It is not a rescaling of the
 holomorphic volume form.
\end{enumerate}

Deformation quantization of $K3 \times E$ is one-parameter:
$H^0(\bigwedge^2 T_{K3 \times E}) = \C$, generated entirely by the
K3 symplectic form $\omega_{K3}^{-1} \in H^0(\bigwedge^2 T_{K3})$.
The elliptic curve contributes no Poisson structure
($\bigwedge^2 T_E = 0$ for $\dim_\C E = 1$), and $H^0(T_{K3}) = 0$
eliminates mixed terms. The star product is the Moyal product
in local Darboux coordinates on K3:
\begin{equation}\label{eq:k3e-star-product}
f \star_\hbar g
= f g + \hbar\,\{f, g\}_{K3}
+ \sum_{n \geq 2} \frac{\hbar^n}{n!}\,
\omega^{i_1 j_1} \cdots \omega^{i_n j_n}\,
(\partial_{i_1} \cdots \partial_{i_n} f)\,
(\partial_{j_1} \cdots \partial_{j_n} g).
\end{equation}

The analytic Brauer torsion is read from the exponential sequence. Since
$H^1(K3,\cO_{K3})=0$ and $\operatorname{Br}(E)=0$, there is no elliptic
or mixed Brauer summand. If the K3 factor has Picard rank $\rho$, then
\begin{equation}\label{eq:k3e-brauer}
\operatorname{Br}(K3 \times E)_{\mathrm{tors}}
\;\cong\;
(\Q/\Z)^{22-\rho}.
\end{equation}
For algebraic K3 surfaces one has $\rho\geq 1$; the Picard-rank-zero
case is the analytic non-projective limit.

\begin{remark}[Geometric Hochschild comparison]
\label{rem:k3e-chiral-hochschild}
\index{Hochschild cohomology!K3 x E@$K3 \times E$}
The following diamond is the HKR diamond of the geometric category
$D^b(K3\times E)$. It is a comparison target for the CDR and
factorization models, not an identification of
$\operatorname{ChirHoch}^\bullet(\Omega^{\mathrm{ch}}_{K3\times E})$
with geometric Hochschild cohomology. The full geometric Hochschild
diamond is
\begin{equation}\label{eq:k3e-hh-full}
(\HH^0, \HH^1, \HH^2, \HH^3, \HH^4, \HH^5, \HH^6)
= (1, 2, 23, 44, 23, 2, 1),
\end{equation}
with CY Serre duality $\HH^n \cong \HH^{6-n}$ verified computationally
via the Volume~III K3$\times$E noncommutative-deformation engine.
The boundary algebra $\cA_E$ used above is the rank-$24$ harmonic
reduction; its primitive chiral Hochschild row is
$1+24t+t^2$, not this geometric HKR diamond.
\end{remark}


\subsubsection{Swiss-cheese structure and the CDR shadow}
\label{subsubsec:k3e-swiss-cheese}
\index{Swiss-cheese!K3 x E@$K3 \times E$}
\index{chiral de Rham complex!K3 x E@$K3 \times E$}

The chiral de Rham complex $\Omega^{\mathrm{ch}}(K3 \times E)$
is a sheaf of vertex superalgebras on $K3 \times E$ whose
global sections carry an $\cN = 2$ superconformal structure
(Malikov--Schechtman--Vaintrob). Put
\begin{equation}\label{eq:k3-cdr-sigma-model}
H^*\!\bigl(K3, \Omega^{\mathrm{ch}}_{K3}\bigr)
\;=:\; V_{K3}
\qquad (c = 6,\; \cN = (4,4)\;\text{SCVA}).
\end{equation}
The exterior-product K\"unneth map for CDR sheaves gives
\begin{equation}\label{eq:k3e-cdr-kunneth}
H^*\!\bigl(K3 \times E, \Omega^{\mathrm{ch}}\bigr)
\;\cong\;
V_{K3} \otimes V_E,
\end{equation}
where $V_E=H^*(E,\Omega_E^{\mathrm{ch}})$ is the elliptic chiral de
Rham factor. The MSV central charges are $6$ on K3 and $3$ on~$E$, so
the untwisted product has $c=9$. The rank-$24$ Heisenberg boundary
algebra is not this full CDR factor; it is the harmonic K3 reduction
of Proposition~\ref{prop:k3e-boundary-datum}.

The HT twist changes the stress tensor by the $U(1)$ current and then
passes to BRST cohomology. It does not identify the untwisted CDR
central charge with zero. On the abelian harmonic branch, the boundary
shadow operations of arity $k\geq 3$ vanish on bar cohomology, because
$\cA_E$ is class~$\mathsf{G}$. No formality assertion is made for the
full CDR Swiss-cheese factorization algebra.

\begin{example}[K3 bulk-boundary pair as an \texorpdfstring{$\SCchtop$}{SCchtop} realisation]
\label{ex:k3e-scchtop-bulk-boundary}
The K3 specialisation of $\SCchtop$ is not concentrated at a single
smooth fibre. Its open colour is the harmonic boundary algebra
$\cA_E$ of Proposition~\ref{prop:k3e-boundary-datum}, obtained after
projecting the K3 factor before restriction to~$E$. Its closed colour
is the bulk-boundary pair
\[
 \bigl(C^\bullet_{\mathrm{ch}}(\cA_E,\cA_E),\, \cA_E\bigr)
\]
carried over the bi-based host
\[
 \Ran(E^{\mathrm{nod,sm}}_{24}) \times \overline{\mathcal{A}}_2.
\]
The index $24$ is the Euler count: an elliptic K3 has Euler number
$24$, and in a nodal model each nodal fibre contributes one to that
count, so the generic smoothing problem has exactly twenty-four nodal
walls. The nearby-cycles sector records transport
across these twenty-four walls, while the $\overline{\mathcal{A}}_2$
parameter records the Humbert bifurcation on which the Siegel form
$\Delta_5$ acquires its first monodromy wall
\textup{(}Gritsenko--Nikulin 1998, Theorem~2.1\textup{)}. In this form
the K3 bulk-boundary pair is an honest worked $\SCchtop$ example:
holomorphic fusion happens on $\Ran(E^{\mathrm{nod,sm}}_{24})$, ordered
topological propagation happens on the boundary interval, and the bulk
is seen through nearby cycles rather than through a single smooth
parameter value.
\end{example}

\begin{remark}[Factorization homology and the elliptic genus]
\label{rem:k3e-factorization-homology}
\index{factorization homology!K3 x E@$K3 \times E$}
\index{elliptic genus!K3 x E@$K3 \times E$}
The CDR Euler characteristic gives
$\chi(\Omega^{\mathrm{ch}}(K3 \times E)) = \chi(K3) \cdot \chi(E)
= 24 \cdot 0 = 0$: the Witten genus vanishes. This
vanishing of the elliptic genus does \emph{not} imply
$F_g(\cA_E) = 0$: the elliptic genus is a
\emph{refined} trace with $(-1)^F$ insertion, while
$F_g(\cA_E) = \kappaChHodge(\cA_E) \cdot \lambda_g^{\mathrm{FP}}$
\textup{(}uniform-weight\textup{)} is the bar complex
curvature at genus~$g$ (no grading insertion). The BPS partition
function is related to $(\Phi_{10}^{\mathrm{un}})^{-1}$, the reciprocal
of the unnormalised Igusa cusp form (DMVV formula for
$K3 \times E_\tau$). In the K3 Borcherds normalisation
\[
 \Phi_{10}^{\mathrm{un}}=\Delta_5^2,\qquad
 \mathrm{wt}(\Delta_5)=c_{\phi_{0,1}^{K3}}(0)/2=10/2=5 .
\]
This automorphic/BPS datum is external to the scalar FP genus tower; the
shadow
partition function $Z^{\mathrm{sh}} = \exp(\sum F_g \hbar^{2g})$
is the perturbative shadow of this non-perturbative BPS count.
\end{remark}


\subsubsection{Connection to Volume~I}
\label{subsubsec:k3e-vol1-connection}

On the abelian harmonic branch, the $K3 \times E$ example supplies the
following Volume~I comparison entries:

\begin{center}
\renewcommand{\arraystretch}{1.25}
\small
\begin{tabular}{lll}
\textbf{Vol~I result} & \textbf{$K3 \times E$ instance} & \textbf{Value} \\
\hline
Theorem~A (adjunction) & $D_{\mathrm{Ran}}(B(\cA_E))
 \simeq B(\cA_E^!)$ & free-field Verdier \\
Theorem~B (completed Positselski) & $\Omega(B(\cA_E)) \simeq \cA_E$
 in the coderived/completed ambient; the raw counit is Theorem~A's
 inversion face & immediate (class $\mathsf{G}$) \\
Theorem~C (complementarity) &
 $\kappaChHodge(\cA_E)+\kappaChHodge(\cA_E^!) = 0$;
 $F_g + F_g' = 0$ for all $g$ & $24 + ({-}24) = 0$ \\
Theorem~D (leading coeff.) & $F_g(\cA_E) =
 \kappaChHodge(\cA_E)\lambda_g^{\mathrm{FP}}$
 \textup{(}u.w.\textup{)} & $\kappaChHodge(\cA_E)=24$ \\
Theorem~H (Hochschild) & $P_{\cA}^{\mathrm{prim}}(t) =
 1 + 24\,t + t^2$ & $\HH^0 = 1,\; \HH^1 = 24,\; \HH^2 = 1$ \\
\hline
Shadow class & $\mathsf{G}$ (Gaussian) & \(r^\Theta_{\max} = 2\) \\
Shadow metric & \(Q_L = (2\kappaChHodge(\cA_E))^2 = 2304\) & constant \\
MC element & \(\Theta_\cA = \Theta_\cA^{\leq 2}
= \kappaChHodge(\cA_E)\)
 & terminates at degree $2$ \\
$r$-matrix & $r(z) = k\,\Omega_{\mathrm{tr}}/z$ \textup{(}abelian Casimir, $k=1$\textup{)} &
 single pole, $\dim = 24$
\end{tabular}
\end{center}

\noindent
The moduli space of K3 has dimension~$20$ (complex structure
deformations from $h^{1,1}(K3)$), plus $\dim \cM_E = 1$
and $21$ K\"ahler moduli, for a total of~$42$.
The period domain for K3 complex structures is a
type~$\mathrm{IV}_{20}$ bounded symmetric domain of
complex dimension~$20$, with arithmetic group
$O(\Lambda_{K3}) \cong O(3, 19; \Z)$.
The holographic shadow connection $\nabla^{\mathrm{hol}}$ on the
K3 moduli space is trivial at the shadow metric level
($Q_L$ constant for class~$\mathsf{G}$), but acquires nontrivial
monodromy at the full holographic level from the arithmetic structure
of the K3 period domain.

\begin{remark}[Comparison across CY$_3$ geometries]
\label{rem:k3e-cy3-comparison}
The $K3 \times E$ datum contrasts with the other standard CY$_3$
targets as follows:
\begin{center}
\renewcommand{\arraystretch}{1.25}
\small
\begin{tabular}{lp{0.27\linewidth}p{0.20\linewidth}p{0.19\linewidth}}
\textbf{CY$_3$ datum} & \textbf{Scalar row}
 & \textbf{Depth assertion} & \textbf{Infinitesimal row} \\
\hline
$K3 \times E$, abelian harmonic branch
& $\kappaChHodge=24$ & $\mathsf{G}$ & geometric HKR: $\dim \HH^2=23$ \\
Quintic, KS/BCOV lane
& $\kappa_{\mathrm{KS}}^{(1)}=-100,\;
   \kappa_{\mathrm{BCOV}}=-25/3$
& scalar projection only & $h^{2,1}=101$ \\
D3 ($\cN\!=\!4$ SYM)
& $\frac{N(N+1)}{2}$ & $\mathsf{L}$ & current row \\
ABJM & $-(N^2{+}1)$ & $\mathsf{M}$ & monopole row
\end{tabular}
\end{center}
The $K3 \times E$ row is the free-field Calabi--Yau row in the
table. The quintic row is a Kodaira--Spencer/BCOV row: its
$-100$ is the Hodge supertrace $h^{1,1}-h^{2,1}$, and its one-loop
BCOV scalar is $\chi/24=-25/3$. The full quintic boundary algebra still
contains Schouten and BCOV collision operations, so class~$\mathsf{G}$
is not assigned to it. In the uniform-weight scalar FP model, a
positive scalar characteristic gives nonzero scalar genus coefficients
only after the scalar projection. This is neither an elliptic genus nor
a BPS-state count.
\end{remark}

\begin{remark}[Bulk OPE structure constants via the Gerstenhaber bracket]
\label{rem:bulk-ope-gerstenhaber-bracket}
The derived chiral center $Z^{\mathrm{der}}_{\mathrm{ch}}(\cA)$
carries a Gerstenhaber bracket $\{-,-\}$ encoding the bulk OPE;
this bracket is the chiral analogue of the classical Hochschild
bracket on $\operatorname{ChirHoch}^*(\cA)$.
The shadow depth classification ($\mathsf{G}$/$\mathsf{L}$/$\mathsf{M}$)
controls the truncation level of this bracket.
We record the structure constants for three families.

\emph{Heisenberg $H_k$ (class~$\mathsf{G}$, shadow depth $r=2$).}
The space $\operatorname{ChirHoch}^*(H_k)$ has three generators: the
unit $\mathbf{1}$ in degree~$0$, the current $J$ in degree~$1$,
and the normal-ordered product ${:}JJ{:}$ in degree~$2$.
The Gerstenhaber bracket is determined at degree~$2$:
\[
 \{J, J\} \;=\; k,
\]
where $\kappa^{\mathrm{Heis}} = k$ is the Heisenberg level.
All higher brackets vanish; the only nonzero binary bracket is central.
The bulk dg Lie structure is therefore two-step nilpotent, and it
becomes abelian only at $k=0$.

\emph{Affine $V_k(\mathfrak{sl}_2)$ (class~$\mathsf{L}$, shadow depth $r=3$).}
The bar cohomology satisfies $\dim H^2(B(\mathfrak{sl}_2)) = 5$,
and $\kappa^{\mathrm{KM}}(V_k(\mathfrak{sl}_2)) = 3(k+2)/4$.
The Gerstenhaber bracket on
$Z^{\mathrm{der}}_{\mathrm{ch}}(V_k(\mathfrak{sl}_2))$ is
non-abelian:
\[
 \{J^a, J^b\}
 \;=\; f^{ab}{}_{c} \, J^c
 \;+\; k \, \delta^{ab},
\]
combining the simple-pole (structure constant) and double-pole
(level) contributions from the OPE
$J^a(z) J^b(w) \sim
 k\,\delta^{ab}/(z-w)^2 + f^{ab}{}_{c}\,J^c(w)/(z-w)$.
A non-trivial degree-$3$ component arises from the cubic Casimir
interaction, scaled by a rational function of~$k$.
The tower terminates: all brackets of degree $\geq 4$ are determined
by the lower ones.
At the critical level $k = -2$: $\kappa^{\mathrm{KM}} = 0$, the
center of the completed enveloping algebra jumps
(Feigin--Frenkel), and the Gerstenhaber bracket changes
qualitatively.

\emph{Virasoro $\mathrm{Vir}_c$ (class~$\mathsf{M}$, shadow depth $r = \infty$).}
$\operatorname{ChirHoch}^*(\mathrm{Vir}_c)$ is concentrated in
degrees $\{0,1,2\}$ with total dimension at most~$4$; it is neither
$\mathbf{C}[\Theta]$ nor the Gelfand--Fuchs cohomology (which is
infinite-dimensional).
The leading structure constant is $\kappa^{\mathrm{Vir}} = c/2$,
arising from the $T(z)T(w)$ OPE via $\mathrm{d}\log$-absorption
(OPE quartic pole $\to$ $r$-matrix cubic pole:
\(r_{\mathrm{Vir}}^{\mathrm{coll}}(z) = (c/2)/z^3 + 2T/z\)).
The bracket has non-vanishing components at every degree $m_k$,
$k \geq 2$: despite bounded cohomological amplitude, the algebraic
depth is infinite (class~$\mathsf{M}$).
At $c = 13$ (the self-dual point):
$\kappa^{\mathrm{Vir}} + \kappa^{\mathrm{Vir}!} = 13$.
At $c=0$: the bracket degenerates.

The pattern is sharp: shadow depth controls truncation of the
Gerstenhaber bracket on the derived chiral center.
Class~$\mathsf{G}$ brackets close at degree~$2$;
class~$\mathsf{L}$ at degree~$3$;
class~$\mathsf{M}$ never close.
This is consistent with Theorem~H, which bounds the cohomological
amplitude of $\operatorname{ChirHoch}^*(\cA)$ uniformly, while the
\emph{algebraic} depth of the bracket structure varies by family.
\end{remark}

% ================================================================
% 9.8.19: UNIVERSAL HOLOGRAPHY (master theorem)
% ================================================================

\subsection{Universal holography}\label{subsec:universal-holography}
\index{universal holography!master theorem}
\index{holomorphic-topological theory!functor from chiral algebras}

Fix an admissible boundary datum in the sense of
Theorem~\ref{thm:holographic-recon-shadow}: a chiral algebra $\cA$ on a
smooth projective curve $X$ with conformal vector $\omega$ at
non-critical level, together with a boundary-linear, affine, or DS
realisation and the ambient needed for the bulk--centre comparison.
The physical $3$d theory is not determined by this algebraic datum
alone.  A holographic object includes a realization package
\[
\Xi_{\cA}=(\mathcal F_{\cA,\Xi},\eta^\partial_{\cA},
\chi_{\mathrm{HT},\cA},\mathcal A(\cA)),
\]
where $\mathcal F_{\cA,\Xi}$ is a Costello--Gwilliam
holomorphic-topological BV/factorization model on the slab,
\(\eta^\partial_{\cA}\colon
\Obs^\partial(\mathcal F_{\cA,\Xi})\xrightarrow{\simeq}\cA\)
is the boundary comparison, and
\(\chi_{\mathrm{HT},\cA}\colon
\Zderch(\cA)\to\Obs^{\mathrm{bulk}}(\mathcal F_{\cA,\Xi})\)
is the bulk--Hochschild comparison in the declared ambient
\(\mathcal A(\cA)\).  On the $3$d slab $X\times\R$, the associated
theory is the supplied BV factorisation algebra
\begin{equation}\label{eq:universal-holography-construction}
\mathcal{T}_{\cA,\Xi}
\;:=\;
\mathcal F_{\cA,\Xi}
\quad\text{on}\quad X\times\R_{\ge0},
\qquad
\eta^\partial_\cA\colon
\mathrm{Obs}^{\partial}(\mathcal{T}_{\cA,\Xi})
\xrightarrow{\simeq}\cA .
\end{equation}
On the boundary component $X\times\{0\}$ impose the chosen holomorphic
boundary condition whose mode algebra is $\cA$; in the bulk
$X\times(0,\infty)$ run the corresponding Costello--Li or
Costello--Gaiotto holomorphic Chern--Simons BV quantisation contained
in \(\Xi_\cA\). When the
topologisation datum is present, the improvement
$T^{\mathrm{imp}}=[Q_{\mathrm{tot}},G']$ makes holomorphic translations
$Q$-exact
(Theorems~\ref{thm:E3-topological-km},
\ref{thm:E3-topological-DS},
\ref{thm:E3-topological-DS-general}).
This is $\mathcal{T}_{\cA,\Xi}$ relative to the chosen admissible input.

Four facts hold on this locus: the boundary recovers $\cA$; the bulk
is compared with the chiral derived centre
$Z^{\mathrm{der}}_{\mathrm{ch}}(\cA)$ by
\(\chi_{\mathrm{HT},\cA}\) in the algebraic, cohomological, or
completed/pro ambient specified by the input; the assignment is
functorial for morphisms preserving the chosen realisation and
comparison maps; and the standard shadow classes are covered by the
boundary-linear, affine, and DS constructions on the corresponding
\(\Xi\)-loci.

\begin{theorem}[$\Xi$-realized universal holography on the admissible boundary-linear/DS locus;
\(\ClaimStatusConditional\); licensing tags \(\beta+\gamma+\delta\)]%
\label{thm:holographic-recon-class-m}%
\index{universal holography!Theorem|textbf}%
\index{holomorphic-topological theory!universal functor}%
\index{Drinfeld--Sokolov reduction!holographic functor}%
Let $X$ be a smooth projective curve and let
$\mathrm{ChirAlg}^{\omega,\mathrm{adm},\Xi}_X$ denote the category whose
objects are conformal chiral algebras on $X$ at non-critical level
together with the admissible boundary-linear, affine, or DS datum of
Theorem~\ref{thm:holographic-recon-shadow} and a realization package
\(\Xi_\cA=(\mathcal F_{\cA,\Xi},\eta^\partial_{\cA},
\chi_{\mathrm{HT},\cA},\mathcal A(\cA))\); morphisms preserve the
boundary comparison, the bulk--Hochschild comparison, and the ambient.
Let $\mathrm{HT\text{-}QFT}^{\Xi}_{X\times\R}$ denote the category of
$3$d holomorphic-topological BV theories on $X\times\R$ equipped with
these comparison data. For every
$(\cA,\omega,\mathrm{adm},\Xi_\cA)\in
\mathrm{ChirAlg}^{\omega,\mathrm{adm},\Xi}_X$ in the generic/Koszul
exact locus, with class~$\mathsf M$ read in the cohomological or
weight-completed/pro ambient unless the bounded-shift HPL criterion is
supplied, the supplied model
$\mathcal{T}_{\cA,\Xi}:=\mathcal F_{\cA,\Xi}$ satisfies:
\begin{enumerate}[label=\textup{(\roman*)}]
\item\label{item:uh-boundary}
 $\eta^\partial_\cA\colon
 \mathrm{Obs}^{\partial}(\mathcal{T}_{\cA,\Xi})\xrightarrow{\simeq}\cA$
 as
 $E_1$-chiral algebras on $X$;
\item\label{item:uh-bulk}
	 $\chi_{\mathrm{HT},\cA}\colon
	 Z^{\mathrm{der}}_{\mathrm{ch}}(\cA)\xrightarrow{\simeq}
	 \mathrm{Obs}^{\mathrm{bulk}}(\mathcal{T}_{\cA,\Xi})$
	 is the bulk--centre comparison in the stated ambient; its strict
	 $E_3$-topological realisation is asserted only when the
	 topologisation hypothesis of
	 Convention~\textup{\ref{conv:chd-e3-chiral-meaning}} is part of
	 the admissible datum;
\item\label{item:uh-functor}
 the assignment
 $\Phi_{\mathrm{hol}}\colon
 (\cA,\omega,\mathrm{adm},\Xi_\cA)\mapsto\mathcal{T}_{\cA,\Xi}$
 is a functor
	 $\mathrm{ChirAlg}^{\omega,\mathrm{adm},\Xi}_X\to
	 \mathrm{HT\text{-}QFT}^{\Xi}_{X\times\R}$
	 intertwining Drinfeld--Sokolov reduction only for
	 \(\Xi\)-compatible DS data:
 $\Phi_{\mathrm{hol}}(\mathrm{DS}\,V^k(\fg),\Xi_{\mathrm{DS}})
 =
 \mathrm{DS\text{-}boundary}\,
 \Phi_{\mathrm{hol}}(V^k(\fg),\Xi_{V^k})$;
\item\label{item:uh-coverage}
	 class coverage is complete on the admissible locus: class
	 $\mathsf{G}$ and the free-field part of class $\mathsf{C}$ are
	 realised by Costello--Li abelian holomorphic Chern--Simons, class
	 $\mathsf{L}$ by affine holomorphic Chern--Simons
	 \textup{(}Theorem~\textup{\ref{thm:E3-topological-km}}\textup{)},
	 and class $\mathsf{M}$ by Costello--Gaiotto plus the
	 DS--Hochschild bridge in its cohomological or
	 weight-completed/pro form
	 \textup{(}Theorem~\textup{\ref{thm:ds-hochschild-bridge}}\textup{)}.
\end{enumerate}
\end{theorem}

\begin{proof}
\emph{Step} (i) \emph{-- boundary recovery.}
By construction~\eqref{eq:universal-holography-construction} the
boundary observables on $X\times\{0\}$ are the holomorphic boundary
cochains of the chosen BV theory, and the comparison with \(\cA\) is
the supplied map \(\eta^\partial_\cA\). For the DS route this is
Costello--Gaiotto~\cite[\S 4]{CostelloGaiotto18}; for boundary-linear
algebras it is Costello--Li~\cite{CL16}; for Kac--Moody it is the
holomorphic Chern--Simons boundary
(Theorem~\ref{thm:E3-topological-km}).
Combining these across shadow classes yields
$\eta^\partial_\cA\colon
\mathrm{Obs}^{\partial}(\mathcal{T}_{\cA,\Xi})\simeq\cA$ as
$E_1$-chiral algebras on $X$.

\emph{Step} (ii) \emph{-- bulk--centre comparison.}
The algebraic closed-sector object determined by the boundary algebra
is
\[
Z^{\mathrm{der}}_{\mathrm{ch}}(\cA)
=C^\bullet_{\mathrm{ch}}(\cA,\cA)
\]
in the ambient stated in the theorem.  The physical bulk factorisation
algebra on $X\times\R$ is the chosen
\(\Obs^{\mathrm{bulk}}(\mathcal{T}_{\cA,\Xi})\).  The bridge between
them is the comparison map \(\chi_{\mathrm{HT},\cA}\), not a
construction of a physical bulk from \(\cA\) alone.  In the
boundary-linear exact sector,
Theorem~\ref{thm:boundary-linear-bulk-boundary} supplies the
Morita/derived-centre comparison.  In the physical prefactorisation
scope, Theorem~\ref{thm:bulk_hochschild} supplies
\(\Psi_{\mathrm{Ran}}\), equivalently \(\chi_{\mathrm{HT},\cA}\).  In
the DS sector, Theorem~\ref{thm:ds-hochschild-bridge} gives the
cohomological or completed/pro comparison, with original-complex
chain-level comparison only under the bounded-shift HPL criterion.
If the topologisation datum of
Convention~\ref{conv:chd-e3-chiral-meaning} is present, then
Theorem~\ref{thm:chiral-higher-deligne} upgrades the algebraic
closed-sector comparison to the strict $E_3$-topological
factorisation-algebra comparison.

\emph{Step} (iii) \emph{-- functoriality and DS intertwining.}
A morphism $\varphi\colon\cA\to\cB$ in
$\mathrm{ChirAlg}^{\omega,\mathrm{adm},\Xi}_X$ determines a boundary
morphism between the chosen Costello BV complexes compatible with
\(\eta^\partial\), \(\chi_{\mathrm{HT}}\), and the ambient; the
induced map of algebraic closed-sector objects is the chiral
Hochschild/coderivation comparison supplied as part of the
\(\Xi\)-compatible morphism data; it is not automatic functoriality of
cochains along the algebra map \(\varphi\).
Intertwining with DS is the content of the cross-volume bridge
``DS-bar'' (Volume~I, \texttt{ds-koszul-intertwine}) composed with
the DS--Hochschild bridge~\ref{thm:ds-hochschild-bridge}:
$\Phi_{\mathrm{hol}}\circ\mathrm{DS}=\mathrm{DS}\circ\Phi_{\mathrm{hol}}$
on the \(\Xi\)-compatible locus, because DS is defined as a boundary
condition on holomorphic Chern--Simons and commutes with the supplied
bulk BV quantisation and comparison maps
(Costello--Gaiotto~\cite[\S 5]{CostelloGaiotto18}).

\emph{Step} (iv) \emph{-- class coverage.}
For class~$\mathsf{G}$ (Heisenberg): abelian holomorphic Chern--Simons
on $X\times\R$ with Lie-algebra the Cartan of the lattice; the chosen
bulk observables are the symmetric algebra on $H^{0,1}(X)$ shifted,
and \(\chi_{\mathrm{HT},\mathcal H_k}\) compares them with
$Z^{\mathrm{der}}_{\mathrm{ch}}(\mathcal{H}_k)$.
For class~$\mathsf{L}$ (affine Kac--Moody):
Theorem~\ref{thm:E3-topological-km}.
For class~$\mathsf{C}$ ($\beta\gamma$, symplectic fermion): contact
$W$-algebras realised as Costello--Li holomorphic boundary conditions
on abelian $3$d HT.
For class~$\mathsf{M}$ (Virasoro, $\cW_N$): Costello--Gaiotto construct
the $3$d HT theory whose boundary is $\cW^k(\fg)$ for principal DS;
the comparison
\(\chi_{\mathrm{HT},\cW^k(\fg)}\) identifies its physical bulk with
$Z^{\mathrm{der}}_{\mathrm{ch}}(\cW^k(\fg))$ cohomologically, and at
chain level only in the weight-completed/pro ambient unless the bounded
conformal-weight-shift HPL criterion is verified, by the
DS--Hochschild bridge~\ref{thm:ds-hochschild-bridge}.
The union of the four sub-constructions defines $\Phi_{\mathrm{hol}}$
on $\mathrm{ChirAlg}^{\omega,\mathrm{adm},\Xi}_X$.
\end{proof}

\begin{corollary}[Class~$\mathsf{M}$ cohomological and completed bulk--centre comparison;
\(\ClaimStatusConditional\); licensing tags \(\beta+\gamma+\delta\)]%
\label{thm:class-M-chain-bulk}%
\index{class M@class~$\mathsf{M}$!chain-level bulk--centre comparison|textbf}%
\index{W-algebras@$\cW$-algebras!bulk--centre comparison}%
\index{Virasoro algebra!bulk--centre comparison}%
Let $\cA=\cW^k(\fg)$ be a $\cW$-algebra obtained by principal
Drinfeld--Sokolov reduction of $V^k(\fg)$ at non-critical level, and
equip it with a class-\(\mathsf M\) realization datum
\(\Xi_{\cA}^{\mathrm{DS}}\) consisting of the Costello--Gaiotto
DS-boundary model, the boundary comparison \(\eta^\partial_\cA\), the
bulk--Hochschild comparison \(\chi_{\mathrm{HT},\cA}\), and the
weight-completed/pro ambient.  Then
\[
Z^{\mathrm{der}}_{\mathrm{ch}}(\cA)
\;\simeq\;
\mathcal{O}\bigl(M_{\mathrm{vac}}^{\mathrm{der}}\bigr)
\]
is the DS algebraic closed-sector comparison cohomologically.  The
physical bulk comparison is the supplied map
\[
\chi_{\mathrm{HT},\cA}\colon
Z^{\mathrm{der}}_{\mathrm{ch}}(\cA)
\xrightarrow{\simeq}
\mathrm{Obs}^{\mathrm{bulk}}(\mathcal{T}_{\cA,\Xi}).
\]
It is chain-level in the weight-completed/pro ambient, and on the
original complex only after the bounded conformal-weight-shift HPL
criterion is supplied.  Thus the class-\(\mathsf M\) global triangle is
a triangle of comparison maps in the declared ambient, not a
reconstruction of a physical bulk from \(\cA\) alone.
\end{corollary}

\begin{proof}
Apply Theorem~\ref{thm:holographic-recon-class-m}\ref{item:uh-bulk} to
$\cA=\cW^k(\fg)$ equipped with \(\Xi_{\cA}^{\mathrm{DS}}\).  The
physical comparison is
\(\chi_{\mathrm{HT},\cA}\colon
Z^{\mathrm{der}}_{\mathrm{ch}}(\cA)\to
\mathrm{Obs}^{\mathrm{bulk}}(\mathcal{T}_{\cA,\Xi})\).
Compose its source with the DS--Hochschild bridge
Theorem~\ref{thm:ds-hochschild-bridge}:
\[
Z^{\mathrm{der}}_{\mathrm{ch}}(\cW^k(\fg))
\;\xrightarrow{\;\mathrm{HKR}^{\mathrm{DS}}\;\simeq\;}\;
\mathcal{O}\bigl(M_{\mathrm{vac}}^{\mathrm{der}}\bigr),
\]
  cohomologically under the smooth chiral-HKR/derived-vacuum package
  and the completed DS coderivation transfer of
  Theorems~\ref{thm:chd-ds-hochschild} and
  \ref{thm:ds-hochschild-bridge}.  The finite input is the
  Kazhdan/PBW finite-weight profile of the BRST complex; generic
  non-critical level is not being used as a \(C_2\)-cofinite
  hypothesis. The same comparison is chain-level in the
  weight-completed/pro ambient, and on the original complex only after
  the bounded conformal-weight-shift HPL criterion of
  Remark~\ref{rem:class-M-hpl-convergence}; this is precisely the scope
  of Theorem~\ref{thm:ds-hochschild-bridge}.
The class~$\mathsf{M}$ global triangle, left open by the
boundary-linear argument
(Theorem~\ref{thm:boundary-linear-bulk-boundary}, which requires
linearity in normal directions that DS violates), is therefore closed
as a cohomological and completed/pro comparison triangle.
\end{proof}

\begin{remark}[Admissible construction]%
\label{rem:universal-holography-admissible-construction}%
The functor $\Phi_{\mathrm{hol}}$ is a construction on the
\(\Xi\)-realized admissible domain: each object includes a conformal
vector, a boundary-linear, affine, or DS origin, a physical BV model, a
boundary comparison, a bulk--Hochschild comparison, and an ambient. The
four-part statement (boundary comparison, bulk--centre comparison,
functoriality with DS, class coverage) is asserted only in that domain.
The bulk observables are identified by the named comparison theorems
for $\SCchtop$ acting on the pair
$(Z^{\mathrm{der}}_{\mathrm{ch}}(\cA),\cA)$.
The DS--Hochschild bridge~\ref{thm:ds-hochschild-bridge} closes the
class~$\mathsf{M}$ bulk comparison cohomologically and in the
completed/pro chain-level ambient; the original-complex chain-level
statement remains governed by the bounded-shift HPL criterion.
\end{remark}


%% ========================================================================
%% Verdier distance and chiral Hochschild comparison.
%% ========================================================================

\section{Holographic distance via the Verdier pairing}
\label{sec:holographic-verdier-distance}

The Verdier pairing gives a distance only after the physical error
filtration has been compared with the bar filtration.

\begin{theorem}[Verdier-realised holographic code distance;
\ClaimStatusProvedHere]
\label{thm:hc-verdier-distance}
Let $\cA$ be a chirally Koszul algebra on a smooth projective curve
$X$, and let $(\cA,\cA^!)$ be its Verdier--Koszul pair. Let
$V\colon \cH_{\mathrm{bdy}}\to\cH_{\mathrm{bulk}}$ be a holographic
encoding with physical error dg algebra $\mathcal E$. Suppose
$\mathcal E$ carries an exhaustive weight filtration and that there is a
filtered morphism
\[
 \operatorname{gr}\mathcal E
 \longrightarrow
 \operatorname{gr}\bigl(\mathrm{Bar}(\cA)\otimes\mathrm{Bar}(\cA^!)\bigr)
\]
which identifies the associated graded of the undetectable error ideal
with the radical of the Verdier evaluation
\[
 \mathrm{ev}_\cA\colon
 \mathrm{Bar}(\cA)\otimes\mathrm{Bar}(\cA^!)\longrightarrow \mathbf 1
\]
and preserves weight. Then
\[
 d_{\mathrm{QECC}}(\cA) \;=\; d_{\mathrm{Verdier}}(\cA),
\]
where
\[
 d_{\mathrm{Verdier}}(\cA)
 =
 \inf\{w\geq 1\mid
 \operatorname{Rad}(\mathrm{ev}_\cA)\cap
 \operatorname{gr}_w(\mathrm{Bar}(\cA)\otimes\mathrm{Bar}(\cA^!))\neq 0\}.
\]
The convention is $\inf\varnothing=+\infty$.
\end{theorem}

\begin{proof}
The code distance is the least weight of a nonzero undetectable error.
By the filtered comparison hypothesis, the associated graded of that
ideal is the Verdier radical, and the comparison preserves weight.
Thus the least physical weight of an undetectable error is the least bar
weight of a nonzero Verdier-radical class.
\end{proof}

\begin{remark}[Genus-zero calibrations]
\label{cor:hc-verdier-distance-examples}
The genus-zero computations give three calibration values:
\begin{itemize}
 \item Heisenberg $H_k$: $d_{\mathrm{Verdier}}(H_k) = +\infty$
 \textup{(}free-field, no obstruction\textup{).}
 \item Virasoro at Lee--Yang $\mathrm{Vir}_{c = -22/5}$:
 $d_{\mathrm{Verdier}} = 3$.
 \item HaPPY pentagon code: $d_{\mathrm{Verdier}} = 3$, matching the
 published QECC distance $d_{\mathrm{QECC}} = 3$ of the
 $[[5, 1, 3]]$ pentagon code on the pentagon model.
\end{itemize}
Each equality is an instance of Theorem~\ref{thm:hc-verdier-distance}
after the weight filtration of the physical error algebra is identified
with the displayed bar filtration.
\end{remark}

\begin{proposition}[BRST lower bound on the Verdier distance;
\ClaimStatusProvedHere]
\label{rem:brst-verdier-lower-bound}
Assume $\cA$ admits a quasi-free BRST resolution whose ghost-spin
filtration maps to the Verdier bar filtration and sends every nonzero
radical class to bar weight at least twice its ghost spin. If the
Verdier radical is zero, then $d_{\mathrm{Verdier}}(\cA)=+\infty$.
Otherwise, if
$\lambda_{\min}^{\mathrm{ghost}}$ is the smallest ghost spin of a
nonzero radical class, then
\[
 d_{\mathrm{Verdier}}(\cA) \;\geq\; 2 \lambda_{\min}^{\mathrm{ghost}},
\]
and equality holds precisely when a Verdier-radical class occurs at the
minimal ghost spin and has bar weight
$2\lambda_{\min}^{\mathrm{ghost}}$.
\end{proposition}

\begin{proof}
The Verdier distance is the least bar weight of a nonzero radical
class. The stated filtration inequality gives bar weight
$\geq 2\lambda_{\min}^{\mathrm{ghost}}$ for every such class. The last
assertion is the equality case.
\end{proof}

\begin{conjecture}[Genus-1 Verdier distance]
\label{conj:verdier-distance-g1}
At genus $1$, the Verdier-pairing distance is governed by the elliptic
bar pairing and the Felder form factor. No additive formula follows
from the genus-zero pairing without this elliptic comparison.
\end{conjecture}

\begin{conjecture}[Higher-genus Verdier distance]
\label{conj:verdier-distance-higher-genus}
At genus $g \geq 2$, the Verdier-pairing distance is computed by a
spectral sequence whose $E_2$ page is the genus-$g$ Hochschild cohomology
of the Koszul pair together with the multi-weight Verdier correction.
The QECC interpretation requires a fibrewise error filtration over
$\overline{\mathcal M}_g$.
\end{conjecture}


\section{The chiral Hochschild comparison theorem}
\label{sec:chiral-hochschild-trinity}

The Fulton--MacPherson, enveloping, and factorization-homology models
are compared by two quasi-isomorphisms.

\begin{theorem}[Chiral Hochschild comparison theorem;
\ClaimStatusProvedHere]
\label{thm:chiral-hochschild-trinity}
\phantomsection\label{prop:chiral-hochschild-comparison-locus}
Let $\cA$ be a logarithmic finite-type chiral algebra in a filtered
chiral ambient in which the geometric, algebraic, and
factorization-homology Hochschild models are defined. Suppose the
comparison maps
\[
 \Phi_{\mathrm{GA}}\colon
 C^\bullet_{\mathrm{chiral}}(\cA)
 \longrightarrow
 \mathrm{Ext}^*_{\cA^e}(\cA,\cA),
 \qquad
 \Phi_{\mathrm{AB}}\colon
 \mathrm{Ext}^*_{\cA^e}(\cA,\cA)
 \longrightarrow
 \mathrm{RHH}_{\mathrm{ch}}(\cA)
\]
are filtered quasi-isomorphisms compatible with the circle
factorization product. Then
\[
 \boxed{\;C^\bullet_{\mathrm{chiral}}(\cA)
 \;\xrightarrow[\Phi_{\mathrm{GA}}]{\sim}\;
 \mathrm{Ext}^*_{\cA^e}(\cA, \cA)
 \;\xrightarrow[\Phi_{\mathrm{AB}}]{\sim}\;
 \mathrm{RHH}_{\mathrm{ch}}(\cA)
 \;=\; \int_{S^1} \cA.\;}
\]
The geometric Fulton--MacPherson complex, the algebraic enveloping
Ext-complex, and the derived chiral Hochschild complex therefore carry
one chiral Hochschild object.
\end{theorem}

\begin{proof}
The assertion is the two-out-of-three statement for the displayed
zigzag of filtered quasi-isomorphisms. Compatibility with circle
factorization identifies the last complex with
$\int_{S^1}\cA$ and transports the product and bracket operations
through the zigzag.
\end{proof}

\begin{theorem}[$\kappa$-conductor as trace on chiral Hochschild;
\ClaimStatusProvedHere]
\label{cor:kappa-conductor-trinity-centre}
In the hypotheses of Theorem~\ref{thm:chiral-hochschild-trinity},
suppose the conductor operator $K_\cA$ is trace-class on the common
Hochschild object and that the Vol~I conductor identity
$K(\cA)=-c_{\mathrm{ghost}}(\mathrm{BRST}(\cA))$ holds in the same
ambient. Then
\[
 K(\cA) \;=\; \mathrm{tr}_{Z(\cA)}(K_\cA)
 \;=\; -c_{\mathrm{ghost}}(\mathrm{BRST}(\cA)).
\]
Here $Z(\cA)$ denotes any of the three quasi-isomorphic Hochschild
models of Theorem~\ref{thm:chiral-hochschild-trinity}.
\end{theorem}

\begin{proof}
The trace-class hypothesis defines $\mathrm{tr}_{Z(\cA)}(K_\cA)$ on
the common Hochschild object. The Vol~I conductor identity identifies
this trace with $-c_{\mathrm{ghost}}(\mathrm{BRST}(\cA))$. The
quasi-isomorphisms of Theorem~\ref{thm:chiral-hochschild-trinity}
transport the trace between the three Hochschild models.
\end{proof}

\begin{example}[K3 realisation of the Hochschild centre supertrace]
\label{ex:trinity-k3-centre-trace}
On the K3 Hall--Borcherds comparison locus, the centre trace is a
bridge between chiral and automorphic data. At $d=2$, the
Mukai-Heisenberg output of the K3 category is an $E_2$-chiral
Heisenberg object, so the Hochschild trace remains on the chiral side
and recovers the K3
$\kappaChHodge$-value. At $d=3$, on the K3 Borcherds row with
Borcherds-side algebra $\mathfrak g_{\Delta_5}$, the same Vol~II
centre trace specialises through the Universal $\Phi$-Trace Identity to
\[
 \Phi_\ast\bigl(K(\cA_{\mathfrak g_{\Delta_5}})\bigr)
 \;=\;
 \Phi_\ast\!\bigl(
 \mathrm{tr}_{Z(\cA_{\mathfrak g_{\Delta_5}})}
 (K_{\cA_{\mathfrak g_{\Delta_5}}})\bigr)
 \;=\;
 \kappa_{\mathrm{BKM}}(\mathfrak g_{\Delta_5})
 \;=\;
 \frac{c_{\phi_{0,1}^{K3}}(0)}{2}
 \;=\;
 5,
\]
where $c_{\phi_{0,1}^{K3}}(0)$ is the constant term of the K3
elliptic-genus Jacobi input in the Borcherds lift
\textup{(}Borcherds 1998; Gritsenko 1999\textup{)}, so
$c_{\phi_{0,1}^{K3}}(0)/2=10/2=5=\mathrm{wt}(\Delta_5)$, and the
dyonic square is $\Phi_{10}^{\mathrm{un}}=\Delta_5^2$. The centre is
therefore the place where the K3 datum can be read both as a Vol~II
supertrace and as a Vol~III Borcherds weight, once the K3
Hall--Borcherds comparison data are supplied.
\end{example}

\begin{remark}[Amplitude and occupation]
\label{rem:hz3-14-amplitude-occupation}
If one of the three Hochschild models of
Theorem~\ref{thm:chiral-hochschild-trinity} has cohomological amplitude
$[0,2]$, the other two have the same amplitude. For the Virasoro
central-extension sector, the occupied degrees are $\{0,2\}$.
\end{remark}

\begin{conjecture}[Comparison extensions: non-Koszul / trace-conductor /
 higher-genus / CY3 / pentagon-coherence]
\label{conj:trinity-extensions}
Five extensions of the comparison-locus statement record the axes
along which the equivalence may break: \textup{(i)} non-Koszul chiral
algebras may admit the comparison only after passing to a Koszul double cover;
\textup{(ii)} the supertrace identification
$K(\cA) = \mathrm{str}(K_\cA)$ extends
to non-Koszul trace-class objects; \textup{(iii)} at higher genus the
three models split into a stratification rather than a coincidence;
\textup{(iv)} the CY$_3$ chiral comparison follows from the Vol~III
$\Phi_3$ functor under CY-A$_3$; \textup{(v)} the Vol~II
$\SCchtop$ pentagon coherence at five presentations controls the
single-colour Hochschild comparison.
\end{conjecture}

\section{Synthesis}
\label{sec:thqghr-supplement}

\begin{remark}[K3 Hall--Borcherds comparison datum]
\label{rem:thqghr-hall-drinfeld}
At $K3\times T^2$ the Vol~II reconstruction is compatible with the
Vol~III Hall--Borcherds target
$\mathcal{D}_\hbar(\mathcal{Y}^{\mathrm{Hall}}_\hbar
(\mathrm{CoHA}_{K3 \times E}))$ only after the oriented hCS--Hall map,
the protected trace comparison, and the Borcherds-Serre quotient data
are supplied. The automorphic normalisation is
\[
 \Phi_{10}^{\mathrm{un}}=\Delta_5^2,\qquad
 \kappa_{\mathrm{BKM}}(\mathfrak g_{\Delta_5})
 =c_{\phi_{0,1}^{K3}}(0)/2=5 .
\]
The one-dimensionality of the $\Delta_5$ automorphic line controls the
Borcherds denominator; it does not by itself construct the full
Hall--Drinfeld double or the twisted 11D-SUGRA bulk.
\end{remark}

\section{DGGPR state sums and conditional gravitational readings}
\label{sec:thqghr-gravitational-partition-RT}

Witten 1989 \emph{CMP}~121 Thm.~2.4 identifies the semisimple
Reshetikhin--Turaev invariant $\tau_{\mathcal C_k}(M^3)$ with the
partition function of three-dimensional $G$-Chern--Simons theory
at level $k$:
\[
 Z_{\mathrm{CS}}(M^3; G, k)
 \;=\;
 \int \mathcal DA\,
 \exp\!\Bigl(\tfrac{ik}{4\pi}\!\int_{M^3}\mathrm{tr}
 (A\wedge\mathrm dA + \tfrac23 A\wedge A\wedge A)\Bigr)
 \;=\;
 \tau_{\mathcal C_k}(M^3).
\]
The non-semisimple DGGPR invariant $\tau_{\mathcal C}(M^3)$ attached
to $\mathcal C = \mathrm{Rep}^{\mathrm{fd}}(\mathfrak u_{\zeta_8}
(\widehat{\mathfrak m}_{\Delta_5}))$ is defined by Hopf-integral and
modified-trace topology (Hennings 1996; Kerler--Lyubashenko 2001;
De~Renzi--Geer--Patureau-Mirand). Its interpretation as a
three-dimensional HT gravitational path integral is an additional
holographic hypothesis: the exact-sector functor, the boundary/asymptotic
condition, and the relevant saddle expansion must be fixed. In the
formulas below
$Z_{\mathrm{grav}}(M^3;\mathbf H_{\Delta_5})$ denotes this conditional
reading of the already-defined DGGPR invariant.

\begin{theorem}[DGGPR state sums on $S^3$, $S^2\!\times\! S^1$, $\Sigma_g\!\times\! S^1$]
\label{thm:thqghr-rt-canonical-manifolds}
Under the admissible-locus holographic reconstruction of
Theorem~\ref{thm:universal-holography-master},
the DGGPR invariant $\tau_{\mathcal C}(M^3)$ computes the closed
topological sector of the proposed bulk 3d holomorphic-topological
gravity coupled to the boundary chiral algebra $\mathbf H_{\Delta_5}$.
Its interpretation as a gravitational path integral is conditional on
the exact-sector holography functor and the chosen boundary/asymptotic
condition. On canonical 3-manifolds:
\begin{align}
 Z_{\mathrm{grav}}(S^3; \mathbf H_{\Delta_5})
 &\;=\;\tau_{\mathcal C}(S^3)
 \;=\;\mathcal D(\mathcal C)^{-1},
 \label{eq:thqghr-Z-S3}\\
 Z_{\mathrm{grav}}(S^2\!\times\! S^1; \mathbf H_{\Delta_5})
 &\;=\;\tau_{\mathcal C}(S^2\!\times\! S^1)
 \;=\;|\Lambda_\infty|,
 \label{eq:thqghr-Z-S2S1}\\
 Z_{\mathrm{grav}}(\Sigma_g\!\times\! S^1; \mathbf H_{\Delta_5})
 &\;=\;\tau_{\mathcal C}(\Sigma_g\!\times\! S^1)
 \;=\;\dim Z^{\mathrm{DGGPR}}(\Sigma_g),
 \label{eq:thqghr-Z-SigmagS1}
\end{align}
with $\mathcal D(\mathcal C)^2 = |\mathrm{Hopf\ dim}(\mathfrak u_{\zeta_8})|$
the total quantum dimension (Vol~III
Theorem~\ref{V3-thm:modtr-rt-S3}), $|\Lambda_\infty|\leq 8^{129}$ the
PBW-finite rank count (Vol~III Theorem~\ref{V3-thm:modtr-zt2-hh0}),
and $Z^{\mathrm{DGGPR}}(\Sigma_g) = \mathrm{Hom}_{\mathcal C}
(\mathbf 1, \mathcal L^{\otimes g})$. At $g = 2$ the partition function
on $\Sigma_2\times S^1$ coincides with the dimension of the
Maass-Spezialschar weight-$10$ Siegel cusp-form space.
\end{theorem}

\begin{proof}
\eqref{eq:thqghr-Z-S3} is Vol~III Theorem~\ref{V3-thm:modtr-rt-S3}
(Reshetikhin--Turaev empty-link surgery for $S^3$) combined with
Witten 1989 Thm.~2.4 reading $\tau_{\mathcal C}(S^3) = \mathcal D^{-1}$
as the Chern--Simons-like path-integral normalisation. The
Hennings 1996 Thm.~3 Hopf-integral formulation gives
$\tau_{\mathcal C}(S^3) = \mu_r(\Lambda)^{-1}$, and
$\mu_r(\Lambda)^2 = |\mathrm{Hopf\ dim}(\mathfrak u_{\zeta_8})|$ up
to the ribbon-twist anomaly of Kerler--Lyubashenko 2001 \S 5.4.
\eqref{eq:thqghr-Z-S2S1} is Atiyah 1988 \S 1 Axiom 4
($\tau_{\mathcal C}(\Sigma\times S^1) = \dim Z^{\mathrm{DGGPR}}(\Sigma)$)
applied at $\Sigma = S^2$ to give
$\dim Z^{\mathrm{DGGPR}}(S^2) = \dim\mathrm{HH}_0^{\mathrm{cat}}
(\mathcal C) = |\Lambda_\infty|$ (Vol~III
Theorem~\ref{V3-thm:modtr-rt-S2timesS1}). \eqref{eq:thqghr-Z-SigmagS1}
is Atiyah Axiom 4 at genus $g$, with the $g = 2$ specialisation
matching the Fourier-coefficient decomposition $c(0, 0, 0) = 1$ of
$(\Phi_{10}^{\mathrm{un}}(Z))^{-1}$ (Vol~III Theorem~\ref{V3-thm:modtr-rt-SigmagS1}
and Eichler--Zagier 1985 \S 6).

\emph{Conditional HT interpretation.} The boundary CFT lives on
$\partial M^3$: for $M^3 = S^3$ there is no boundary, and
$\mathcal D^{-1}$ is the Euclidean partition-function
normalisation. For $M^3 = S^2\times S^1$, the boundary is
$S^2$ and the $S^1$ time-cycle computes the thermal partition
function at zero temperature, counting the closed-string
ground states via $\mathrm{HH}_0^{\mathrm{cat}}$. For $M^3 =
\Sigma_g\times S^1$, the same reading with $\Sigma_g$-sliced
boundary gives the trace of the identity on
$Z^{\mathrm{DGGPR}}(\Sigma_g)$, at $g = 2$ reproducing the
Dijkgraaf--Moore--Verlinde--Verlinde 1997 \emph{CMP}~185 \S 4
second-quantised elliptic-genus partition function
$(\Phi_{10}^{\mathrm{un}}(Z))^{-1}$.
\end{proof}

\begin{theorem}[Lens-space DGGPR invariant and the $\mu_8$ gerbe]
\label{thm:thqghr-rt-lens}
For the lens space $L(p, q)$,
\[
 Z_{\mathrm{grav}}(L(p, q); \mathbf H_{\Delta_5})
 \;=\;
 \tau_{\mathcal C}(L(p, q))
 \;=\;
 \mathcal D^{-1}\sum_{\lambda\in\Lambda_\infty}
 S_{0\lambda}\,\theta_\lambda^{q^\star}\,S_{\lambda 0},
\]
where $S$ is the DGGPR $S$-matrix, $\theta_\lambda$ the ribbon twist,
and $q^\star\equiv q^{-1}\pmod p$. At $p = 8$, the $S$-matrix coincides
with the Fricke involution $w_8\in\mathrm{Sp}_4(\mathbb Z)$ (Vol~III
Theorem~\ref{V3-thm:modtr-s-matrix-fricke}), so
\[
 Z_{\mathrm{grav}}(L(8, 1); \mathbf H_{\Delta_5})
 \;=\;
 \mathcal D^{-1}\cdot\mathrm{tr}\bigl(w_8\theta\bigm\vert
 \mathrm{HH}_0^{\mathrm{cat}}(\mathcal C)\bigr),
\]
the Fricke-twisted trace on the $\mathrm{HH}_0^{\mathrm{cat}}$ state
space. The $\mu_8$-gerbe obstruction of Vol~III
Theorem~\ref{V3-thm:modtr-banding-cocycle} is an invertible bulk
framing/gerbe phase of the closed 3d gauge theory on $L(8, 1)$. Since
the lens space is closed, there is no boundary CFT on
$\partial L(8,1)$; boundary data enter only through a cobordism with
boundary or through the state space assigned to a cutting surface. The
order-$8$ phase tracks the eighth-root ambiguity in the
Kerler--Lyubashenko ribbon twist.
\end{theorem}

\begin{proof}
The surgery formula is Vol~III Theorem~\ref{V3-thm:modtr-rt-lens}:
$L(p, q) = -q/p$-surgery on the unknot in $S^3$, giving
$\tau_{\mathcal C}(L(p, q)) = \mathcal D^{-1}(ST^{q^\star}S)_{00}$.
At $p = 8$, $q = 1$, the $S$-matrix is the Fricke involution
$w_8$ on the $\mu_8$-gerbe-refined state space (Vol~III
Theorem~\ref{V3-thm:modtr-s-matrix-fricke}), so the $(0, 0)$-entry
of $STS$ computes the Fricke-twisted trace
$\mathrm{tr}(w_8\theta|\mathrm{HH}_0^{\mathrm{cat}}(\mathcal C))$.

\emph{Conditional anomaly reading.} The lens space $L(8, 1)$ is the
$\mathbb Z/8$-quotient $S^3/\mathbb Z_8$ by the diagonal action
$(z_1, z_2)\mapsto(e^{2\pi i/8}z_1, e^{-2\pi i/8}z_2)$
(Milnor 1975 \emph{Topology}~14 \S 2). Since $L(8,1)$ is closed, no
boundary partition function lives on $\partial L(8,1)$. The
$\mu_8$-gerbe of Vol~III Theorem~\ref{V3-thm:modtr-banding-cocycle}
instead enters through the closed-state trace after cutting the
3-manifold along a surface. Under the HT interpretation, this is the
't~Hooft anomaly matching the modular anomaly of the
Kerler--Lyubashenko MTC under
$\mathrm{Sp}_4(\mathbb Z)$-action.

\emph{Multi-path verification:}
(i) Reshetikhin--Turaev 1991 \S 3 Thm.~3.3;
(ii) Vol~III Theorem~\ref{V3-thm:modtr-rt-lens};
(iii) Vol~III Theorem~\ref{V3-thm:modtr-s-matrix-fricke} (Fricke identification);
(iv) Kirby 1997 \S 4.1 calculus on $L(p, q)$.
\end{proof}

\begin{theorem}[Mapping-torus DGGPR trace under Fricke]
\label{thm:thqghr-rt-mapping-torus}
For a mapping class $\phi\in\mathrm{MCG}(\Sigma_g)$ and mapping torus
$M_\phi^3 = \Sigma_g\times_\phi S^1$,
\[
 Z_{\mathrm{grav}}(M_\phi^3; \mathbf H_{\Delta_5})
 \;=\;
 \tau_{\mathcal C}(M_\phi^3)
 \;=\;
 \mathrm{Tr}_{Z^{\mathrm{DGGPR}}(\Sigma_g)}\!\bigl(\rho_{\Sigma_g}(\phi)\bigr),
\]
with $\rho_{\Sigma_g}$ the Lyubashenko--Virelizier MCG representation.
At $\Sigma_g = \Sigma_2$, $\phi = w_8$ (Fricke involution, Vol~III
Theorem~\ref{V3-thm:modtr-s-matrix-fricke}),
\[
 Z_{\mathrm{grav}}(\Sigma_2\times_{w_8} S^1; \mathbf H_{\Delta_5})
 \;=\;
 \mathrm{tr}\bigl(w_8\bigm\vert
 \mathrm{HH}_0^{\mathrm{cat}}(\mathcal C\boxtimes\mathcal C^{\mathrm{op}})\bigr),
\]
which restricts to the Fricke trace $4$ on the $M_{24}$-invariant
Humbert block of Vol~III Theorem~\ref{V3-thm:modtr-s-matrix-fricke}
and to $-4 + (1+i)\sqrt 2$ on the $\mu_{16}$-banded block of Vol~III
Remark~\ref{V3-rem:modtr-banding-cocycle-fricke-trace}.
\end{theorem}

\begin{proof}
The mapping-torus trace formula is Atiyah 1988 \S 2 Thm.~2.4 applied
functorially: $Z(M_\phi^3) = \mathrm{tr}(Z(\phi))$ since $M_\phi^3$
is the self-composition of the cobordism $\Sigma_g\times [0, 1]$
along the gluing map $\phi$. The non-semisimple extension is
Lyubashenko 1995 \S 5.5 Thm.~5.5.1 and De~Renzi et~al.\ 2021 Thm.~5.11.
Vol~III Theorem~\ref{V3-thm:modtr-rt-mapping-torus} identifies the
values on the $M_{24}$- and $\mu_{16}$-invariant blocks.

\emph{Conditional HT reading.} $M_\phi^3 = \Sigma_g\times_\phi S^1$ is a
3d spacetime where the spatial $\Sigma_g$ evolves along a closed
time-loop under the diffeomorphism $\phi$; the DGGPR invariant
is the thermal trace $\mathrm{tr}(\phi_*)$ of the action of $\phi$ on
the Hilbert space of states $Z^{\mathrm{DGGPR}}(\Sigma_g)$. The
specific case $\phi = w_8$ on $\Sigma_2$ gives the conditional HT
partition function on the Fricke mapping torus, which by Vol~III
Theorem~\ref{V3-thm:modtr-s-matrix-fricke} computes the trace of the
Siegel Fricke involution $w_8\colon Z\mapsto-(8Z)^{-1}$ on
$\mathfrak H_2$. This trace is the $M_{24}$-Mathieu character
contribution $\chi_{1A}(\tau, z) = 24$ reduced modulo the Fricke
fixed-locus constraint $\mathrm{Fix}(w_8) = H_1\cup H_4$.
\end{proof}

\begin{theorem}[Poincar\'e sphere DGGPR invariant: non-semisimple test]
\label{thm:thqghr-rt-poincare}
For the Poincar\'e homology sphere $\Sigma(2, 3, 5)$, presented as
$(-1)$-framed $E_8$-plumbing (Kirby 1997 \S 4 Example 4.2),
\[
 Z_{\mathrm{grav}}(\Sigma(2, 3, 5); \mathbf H_{\Delta_5})
 \;=\;
 \tau_{\mathcal C}(\Sigma(2, 3, 5))
 \;=\;
 \mathcal D^{-9}\,\Delta^{8}\,
 \mathrm{CL}^{\mathrm{mod}}_{\mathcal C}(E_8; -1),
\]
with $\Delta = \sum_\lambda\theta_\lambda^{-1} d_\lambda^2$ the RT
anomaly and $\mathrm{CL}^{\mathrm{mod}}$ the Geer--Kujawa--Patureau-Mirand
2011 \emph{Compos.~Math.}~147 Thm.~1 modified-trace coloured-link
invariant. On the non-semisimple $\mathbf H_{\Delta_5}$-MTC the DGGPR
state sum is non-zero, in contrast to the semisimple
$\mathfrak{sl}_2$-level-$k$ case where
$\tau(\Sigma(2,3,5)) = \sum_{n=1}^{k+1}\sin^2(\pi n/(k+2))
e^{2\pi i(n^2-1)/8(k+2)}$ (Lawrence--Zagier 1999 \emph{Asian~J.~Math.}~3
Thm.~1): the Poincar\'e sphere is a non-trivial falsification target
for the non-semisimple renormalisation.
\end{theorem}

\begin{proof}
The surgery formula is Vol~III Theorem~\ref{V3-thm:modtr-rt-poincare}:
$b_0(E_8) = 8$, $\sigma(E_8) = -8$, so
$\tau = \mathcal D^{-9}\Delta^{8}\mathrm{CL}^{\mathrm{mod}}$. The
non-vanishing of $\mathrm{CL}^{\mathrm{mod}}$ on the $E_8$-plumbing
follows from the De~Renzi et~al.\ 2021 Thm.~6.3 non-degeneracy of the
modified trace on projective covers of
$\mathfrak u_{\zeta_8}(\widehat{\mathfrak m}_{\Delta_5})$.

\emph{Conditional HT reading.} $\Sigma(2, 3, 5) = \mathrm{SU}(2)/I^*$ is the
spherical space form with binary icosahedral symmetry $I^*$ (Milnor
1975 \S 3). Under the holographic-reconstruction hypothesis it is read
as a 3-dimensional saddle with $I^*$-isometry group, and
$\tau_{\mathcal C}(\Sigma(2, 3, 5))$ is the closed topological sector
assigned to that saddle. The $E_8$-plumbing presentation identifies the
surgery framing with the $E_8$-Cartan bilinear form; the signature
$\sigma(E_8) = -8$ contributes a phase $\Delta^{8}$ in the state sum,
matching the Chern--Simons framing anomaly (Witten 1989 \S 2.4). The
non-vanishing of
$\tau_{\mathcal C}(\Sigma(2, 3, 5))$ on the non-semisimple
$\mathbf H_{\Delta_5}$-MTC, contrasted with the Lawrence--Zagier 1999
semisimple formula, is the falsification target for the
non-semisimple renormalisation on the $E_8$-plumbing geometry.

\emph{Multi-path verification:}
(i) Vol~III Theorem~\ref{V3-thm:modtr-rt-poincare} (explicit surgery);
(ii) De~Renzi et~al.\ 2021 Thm.~6.1 (non-semisimple formula);
(iii) Lawrence--Zagier 1999 Thm.~1 (semisimple comparator);
(iv) Witten 1989 \S 2.4 (framing anomaly).
\end{proof}

\begin{remark}[Conditional dictionary $\tau_{\mathcal C}\leftrightarrow$ HT gravity]
\label{rem:thqghr-holographic-dictionary}
The table below records the conditional dictionary between the
closed 3-manifold $M^3$, its DGGPR invariant $\tau_{\mathcal C}(M^3)$,
and the bulk-HT partition-function reading under the
holographic reconstruction of Chapter~\ref{ch:programme-climax-platonic}:
\begin{center}
\begin{tabular}{l|l|l}
$M^3$ & $\tau_{\mathcal C}(M^3)$ & conditional HT reading \\ \hline
$S^3$ & $\mathcal D^{-1}$ & Euclidean normalisation \\
$S^2\!\times\! S^1$ & $|\Lambda_\infty|\leq 8^{129}$ &
 thermal ground-state count \\
$T^2\!\times\! S^1$ & $|\Lambda_\infty|$ &
 closed-string partition function \\
$\Sigma_2\!\times\! S^1$ & $\dim\mathrm{HH}_0^{\mathrm{cat}}
(\mathcal C\boxtimes\mathcal C^{\mathrm{op}})$ &
 second-quantised elliptic genus $(c(0) = 1)$ \\
$L(8, 1)$ & $\mathcal D^{-1}\mathrm{tr}(w_8\theta)$ &
 $\mu_8$-gerbe 't Hooft anomaly \\
$\Sigma_2\times_{w_8}\! S^1$ & $\mathrm{tr}(w_8|\mathrm{HH}_0^{\mathrm{cat}})$ &
 Fricke mapping-torus partition function \\
$\Sigma(2, 3, 5)$ & $\mathcal D^{-9}\Delta^{8}\mathrm{CL}^{\mathrm{mod}}(E_8)$ &
 $E_8$-plumbing saddle with binary-icosahedral $I^*$-isometry \\
\end{tabular}
\end{center}
In each row the left column is the topological data, the centre column
is the DGGPR invariant computed in Vol~III \S\ref{sec:modtr-rt-closed-3manifolds},
and the right column is the conditional HT path-integral reading with
boundary chiral-CFT fixed to $\mathbf H_{\Delta_5}$. The dictionary
extends Witten 1989 Thm.~2.4 (semisimple Chern--Simons $\leftrightarrow$
RT invariant) to the non-semisimple regime, with the
Kerler--Lyubashenko Hopf-integral serving as the path-integral
normalisation. It does not derive a Cardy, BTZ, or Page statement;
those require separate modular invariance, vacuum dominance, saddle
control, and transseries hypotheses.
\end{remark}

\begin{theorem}[DGGPR genus-$g$ state-space dimension tower]
\label{thm:thqghr-rt-genus-g-tower}
For $g\geq 0$,
\[
 Z_{\mathrm{grav}}(\Sigma_g\!\times\! S^1;\mathbf H_{\Delta_5})
 \;=\;
 \tau_{\mathcal C}(\Sigma_g\!\times\! S^1)
 \;=\;
 \dim\mathrm{Hom}_{\mathcal C}(\mathbf 1, \mathcal L^{\otimes g}),
\]
with $\mathcal L = \coend_{\mathcal C}$ the Lyubashenko coend. At low
genus,
\begin{align*}
 \tau_{\mathcal C}(S^2\!\times\! S^1) &\;=\; \dim\mathrm{HH}_0^{\mathrm{cat}}(\mathcal C)
 \;=\; |\Lambda_\infty|\leq 8^{129},\\
 \tau_{\mathcal C}(T^2\!\times\! S^1) &\;=\; \dim\mathrm{HH}_0^{\mathrm{cat}}(\mathcal C)
 \;=\; |\Lambda_\infty|,\\
 \tau_{\mathcal C}(\Sigma_2\!\times\! S^1)
 &\;=\; \dim\mathrm{HH}_0^{\mathrm{cat}}
 (\mathcal C\boxtimes\mathcal C^{\mathrm{op}})
 \;=\; c(0, 0, 0)\cdot|\Lambda_\infty|^{\!{}^{\!2}}
 \;\equiv\; 1\cdot|\Lambda_\infty|^{2}\;\mathrm{mod}\;\mathcal O(q, p, y),
\end{align*}
with $c(0, 0, 0) = 1$ the vacuum Fourier coefficient of
$(\Phi_{10}^{\mathrm{un}})^{-1}$ (Eichler--Zagier 1985 \S 6
Thm.~5.3 normalisation). At
$g\geq 3$ the Birman 1974 Ch.~4 surjection
$\mathrm{MCG}(\Sigma_g)\twoheadrightarrow\mathrm{Sp}_{2g}(\mathbb Z)$
yields
\[
 \tau_{\mathcal C}(\Sigma_g\!\times\! S^1)
 \;\leq\;
 \dim\bigl(\mathrm{HH}_0^{\mathrm{cat}}(\mathcal C)^{\otimes g}\bigr)^{\!{}^{\!\mathrm{Sp}_{2g}(\mathbb Z)}}
 \;=\;
 \dim M^{\mathrm{Sieg}}_{10}\!\bigl(\mathrm{Sp}_{2g}(\mathbb Z)\bigr),
\]
the corresponding weight-$10$ Siegel cusp-form space for genus $g$.
The numerical dimensions are those computed by Igusa 1962
\emph{Invent.}~3 Table II and Tsuyumine 1986 \emph{Mem.~AMS}~63.
\end{theorem}

\begin{proof}
The identity $\tau_{\mathcal C}(\Sigma\times S^1) = \dim Z(\Sigma)$
is Atiyah 1988 \S 1 Axiom 4. At $g = 0$ the coend pair
$\mathrm{Hom}(\mathbf 1, \mathcal L^{\otimes 0})
= \mathrm{End}(\mathbf 1)\cong\mathbb C$, collapsing to the $S^2$
evaluation; the $|\Lambda_\infty|$ count at $g = 1$ is Vol~III
Theorem~\ref{V3-thm:modtr-zt2-hh0} (PBW-truncated pro-limit). At
$g = 2$ the Kerler--Lyubashenko handlebody-vacuum identification
$\mathrm{Hom}(\mathbf 1, \mathcal L^{\otimes 2})
= \iota\bigl((\Phi_{10}^{\mathrm{un}})^{-1}\bigr)\cdot\mathrm{HH}_0^{\mathrm{cat}}
(\mathcal C\boxtimes\mathcal C^{\mathrm{op}})$ is
Theorem~\ref{thm:modbootstrap-state-sum}, and the constant-term
coefficient $c(0, 0, 0) = 1$ is Eichler--Zagier 1985 \S 6 Thm.~5.3
(Fourier--Jacobi expansion
$\Phi_{10}^{\mathrm{un}} = \sum_m\phi_{10, m}e^{2\pi im\tau'}$
with $\phi_{10, 0} = 0$ and $\phi_{10, 1}$ the weight-$10$ index-$1$
Jacobi cusp form with $c(0, 0, 0) = 1$ normalisation).

At $g\geq 3$ the Lyubashenko--Virelizier MCG-representation factors
through $\mathrm{Sp}_{2g}(\mathbb Z)$ (Birman 1974 Ch.~4 surjection
+ Lyubashenko 1995 \S 5.4 Thm.~5.4.1 MCG-action on the coend), so
the state space lies in the $\mathrm{Sp}_{2g}(\mathbb Z)$-invariant
subspace of $\mathrm{HH}_0^{\mathrm{cat}}(\mathcal C)^{\otimes g}$,
equivalently in the corresponding weight-$10$ Siegel-modular space on
$\mathrm{Sp}_{2g}(\mathbb Z)$. The numerical dimensions are those of
Igusa 1962 \emph{Invent.}~3 Table II (genus $2$) and Tsuyumine
1986 \S 4 Thm.~5 (extension to genus $3$ and beyond; cf.\
Freitag 1983 \emph{Siegelsche Modulfunktionen} \S VI.3).

\emph{Multi-path verification:}
(i) Atiyah 1988 Axiom 4;
(ii) Kerler--Lyubashenko 2001 \S 5.4 handlebody-vacuum;
(iii) Birman 1974 Ch.~4 + Lyubashenko 1995 Thm.~5.4.1
(MCG-to-Sp factor);
(iv) Igusa 1962 Table II and Tsuyumine 1986 \S 4 dimensions.
\end{proof}

\begin{theorem}[Conditional Fricke mapping-torus trace]
\label{thm:thqghr-fricke-explicit-trace}
Assume the Vol~III Fricke block decomposition of
$\mathrm{HH}_0^{\mathrm{cat}}
(\mathcal C\boxtimes\mathcal C^{\mathrm{op}})$ under the
$M_{24}$-invariant Humbert block and the $\mu_{16}$-banded block
(Vol~III Remark~\ref{V3-rem:modtr-banding-cocycle-fricke-trace}), with
block traces $4$ and $-4+(1+i)\sqrt2$. Then the conditional Fricke
mapping-torus partition function is
\[
 Z_{\mathrm{grav}}(\Sigma_2\times_{w_8} S^1; \mathbf H_{\Delta_5})
 \;=\;
 \underbrace{4}_{\text{Humbert block}}
 \;+\;
 \underbrace{-4 + (1+i)\sqrt 2}_{\text{$\mu_{16}$-banded block}}
 \;=\;
 (1+i)\sqrt 2
 \;=\;
 2\,e^{i\pi/4}.
\]
The magnitude is $|Z_{\mathrm{grav}}| = 2$, and the phase is
$\arg(Z_{\mathrm{grav}}) = \pi/4$, the $\mu_8$-gerbe-anomaly phase
tracking the Kerler--Lyubashenko ribbon-twist eighth root of unity
(Vol~III Theorem~\ref{V3-thm:modtr-banding-cocycle}).
\end{theorem}

\begin{proof}
The Atiyah--Lyubashenko mapping-torus formula
$\tau_{\mathcal C}(\Sigma_2\times_\phi S^1)
= \mathrm{tr}(\rho_{\Sigma_2}(\phi))$ (Atiyah 1988 \S 2 Thm.~2.4
extended by Lyubashenko 1995 \S 5.5 Thm.~5.5.1 and De~Renzi
et~al.\ 2021 Thm.~5.11) applied to $\phi = w_8$ (Vol~III
Theorem~\ref{V3-thm:modtr-s-matrix-fricke}) gives
$Z_{\mathrm{grav}} = \mathrm{tr}(w_8|\mathrm{HH}_0^{\mathrm{cat}})$.

The assumed block decomposition is Vol~III
Remark~\ref{V3-rem:modtr-banding-cocycle-fricke-trace}:
$\mathrm{HH}_0^{\mathrm{cat}}(\mathcal C\boxtimes\mathcal C^{\mathrm{op}})
= \mathrm{Humbert}\oplus\mathrm{Banded}$ with
$\mathrm{tr}(w_8|\mathrm{Humbert}) = 4$
(Vol~III Remark~\ref{V3-rem:modtr-s-eigenvalues}: Hirzebruch--Zagier
Euler-characteristic $\chi(H_1) + \chi(H_4) - \chi(H_1\cap H_4)
= 2 + 2 - 0 = 4$) and
$\mathrm{tr}(w_8|\mathrm{Banded}) = -4 + (1+i)\sqrt 2$. The latter is
used here as the Vol~III block trace; it is not derived from the
primitive-root Ramanujan sum alone.

Summing: $4 + (-4 + (1+i)\sqrt 2) = (1+i)\sqrt 2$.
Magnitude: $|(1+i)\sqrt 2| = \sqrt{1^2 + 1^2}\cdot\sqrt 2 = 2$.
Argument: $\arg(1+i) = \pi/4$, unchanged by positive-real scaling, so
$\arg((1+i)\sqrt 2) = \pi/4$.

\emph{Multi-path verification:}
(i) Atiyah 1988 Thm.~2.4 mapping-torus trace;
(ii) Hirzebruch--Zagier 1976 Humbert Euler-characteristic;
(iii) Vol~III Remark~\ref{V3-rem:modtr-banding-cocycle-fricke-trace}
(banded-block trace).
\end{proof}

\begin{proposition}[Poincar\'e sphere Seifert expansion, conditional]
\label{thm:thqghr-poincare-seifert-closed-form}
The Poincar\'e sphere $\Sigma(2, 3, 5) = \mathrm{SU}(2)/I^*$ is the
Seifert-fibered space with exceptional fibres of multiplicities
$(2, 3, 5)$ and Euler number $e(\Sigma(2, 3, 5)) = -1/30$. Its
binary-icosahedral isometry group $I^*\subset\mathrm{SU}(2)$ has
order $|I^*| = 120$ and is the cover of the icosahedral rotation
group $I\cong A_5$.

Assume the modified-trace Seifert surgery formula applies to the
non-semisimple MTC $\mathcal C =
\mathrm{Rep}^{\mathrm{fd}}(\mathfrak u_{\zeta_8}
(\widehat{\mathfrak m}_{\Delta_5}))$ with projective-cover label set
$\Lambda_\infty$ and local exceptional-fibre kernels
$\mathrm{SF}^{(2,3,5)}_\lambda$. Then the DGGPR invariant has the form
\[
 \tau_{\mathcal C}\!\bigl(\Sigma(2, 3, 5)\bigr)
 \;=\;
 \mathcal D^{-1}\,e^{i\pi c/4}
 \sum_{\lambda\in\Lambda_\infty}
 \mathrm{t}^{\mathrm{mod}}_{P_\lambda}
 \bigl(\mathrm{SF}^{(2,3,5)}_\lambda\bigr),
\]
where $\mathrm{t}^{\mathrm{mod}}_{P_\lambda}$ is the modified trace on
the projective cover $P_\lambda$ and $c$ is the framing anomaly. The
explicit identification of $\mathrm{SF}^{(2,3,5)}_\lambda$ with the
$S$-, $T$-, twist-, and quantum-dimension data of
$\mathcal C$ is the remaining non-semisimple Seifert kernel
calculation; in the semisimple specialization it reduces to the
Lawrence--Rozansky 1999 \emph{CMP}~205 Thm.~1.5 kernel. The
binary-icosahedral symmetry $I^*$ acts, when the small
quantum-$\mathfrak{sl}_2$ comparison is supplied, through the
$A_5$-orbits on $\Lambda_\infty$, partitioning the sum into orbit-sums
indexed by the five conjugacy classes
$\{1, [\mathrm{ord}2], [\mathrm{ord}3], [\mathrm{ord}5], [\mathrm{ord}5]^*\}$
of $A_5$.
\end{proposition}

\begin{proof}
The Seifert description and the binary-icosahedral group are standard:
$\Sigma(2,3,5)=\mathrm{SU}(2)/I^*$ and
$I^*\cong \mathrm{SL}_2(\mathbb F_5)$. Lawrence--Rozansky 1999
\S 4 Thm.~1.5 gives the corresponding semisimple RT Seifert kernel.
The displayed formula is the same statement after replacing the
semisimple quantum trace by the modified trace of
Geer--Kujawa--Patureau-Mirand 2011 and De~Renzi et~al.\ 2021, under the
explicit hypothesis that the non-semisimple Seifert kernel has been
constructed for the present category. The $A_5$ orbit decomposition is
the Lusztig--Saleur small-quantum-group comparison, again conditional on
transporting that action to the projective-cover label set
$\Lambda_\infty$.

Under the HT reading, the phase $e^{i\pi c/4}$ is the framing anomaly
of the closed topological sector. No metric-gravity saddle expansion is
deduced from this state-sum formula alone.
\end{proof}
